% !TeX program = xelatex
% 数学 LaTeX 排版 Skill v4.0.0；版式保持 v3 金标准. XeLaTeX 编译至少两次. 不分发字体.
% WZ-LOCKED-CLASS-BEGIN
\begin{filecontents*}[overwrite]{elegantbook-original-adapter.cls}
\NeedsTeXFormat{LaTeX2e}
\ProvidesClass{elegantbook-original-adapter}[2026/09/23 v3.0 fixed mathematical book environments]
\LoadClass[10pt,letterpaper,oneside]{book}
\RequirePackage[UTF8,scheme=plain,fontset=fandol]{ctex}
\RequirePackage{fontspec}
\setmainfont{cmunrm.otf}[BoldFont=cmunbx.otf,ItalicFont=cmunti.otf,BoldItalicFont=cmunbi.otf]
\setCJKfamilyfont{sourceheader}[ItalicFont=FandolKai-Regular.otf]{FandolKai-Regular.otf}
\RequirePackage{amsmath,amssymb,mathrsfs}
\RequirePackage{amsthm,etoolbox}
\RequirePackage{xcolor,geometry,setspace,fancyhdr,enumitem,needspace,xurl}
\definecolor{structurecolor}{HTML}{880E4F}
\definecolor{main}{HTML}{9C27B0}
\geometry{paperwidth=612bp,paperheight=792bp,left=20mm,right=20mm,top=95.04bp,bottom=20mm,headheight=12pt,headsep=25pt,footskip=30pt}
\setstretch{1.6}
\setlength{\parindent}{2em}
\setlength{\parskip}{0pt}
\setlist{nosep,leftmargin=2em}
\AtBeginDocument{\setlength{\abovedisplayskip}{4pt plus 1pt minus 1pt}\setlength{\belowdisplayskip}{4pt plus 1pt minus 1pt}\setlength{\abovedisplayshortskip}{2pt}\setlength{\belowdisplayshortskip}{3pt}}
\fancyhf{}
\fancyhead[L]{{\itshape\CJKfamily{sourceheader}\rightmark}}
\fancyhead[R]{{\itshape\CJKfamily{sourceheader}\leftmark}}
\fancyfoot[C]{\thepage}
\renewcommand{\headrulewidth}{0.4pt}
\renewcommand{\headrule}{{\color{structurecolor}\hrule height\headrulewidth width\headwidth\vskip-\headrulewidth}}
\pagestyle{fancy}
\fancypagestyle{cover}{\fancyhf{}\fancyfoot[C]{\thepage}\renewcommand{\headrulewidth}{0pt}\renewcommand{\headrule}{}}
\RequirePackage{hyperref}
\hypersetup{unicode=true,colorlinks=true,linkcolor=structurecolor,urlcolor=structurecolor,citecolor=structurecolor,linktoc=section,bookmarksnumbered=true,bookmarksopen=true,pdfborder={0 0 0},pdfstartview=FitH}
\setcounter{tocdepth}{1}
\renewcommand*\l@section{\@dottedtocline{1}{1.5em}{2.4em}}
\renewcommand{\thesection}{\thechapter.\arabic{section}}
\newcommand{\originalchapter}[1]{\par\addvspace{14pt}\Needspace{5\baselineskip}\refstepcounter{chapter}\setcounter{section}{0}\addcontentsline{toc}{chapter}{\protect\numberline{\thechapter}#1}\markboth{\thechapter\quad #1}{}\begingroup\centering\color{structurecolor}\bfseries\fontsize{14.4}{20}\selectfont\thechapter\quad #1\par\endgroup\nobreak\vspace{8pt}}
\newcommand{\originalsection}[1]{\par\addvspace{9pt}\Needspace{4\baselineskip}\refstepcounter{section}\addcontentsline{toc}{section}{\protect\hspace*{1.5em}\protect\numberline{\protect\textcolor{structurecolor}{\thesection}}#1}\markright{\thesection\quad #1}\noindent{\fontsize{12}{17}\selectfont\color{structurecolor}\thesection\quad}{\bfseries\fontsize{12}{17}\selectfont #1}\par\nobreak\vspace{4pt}}

% Semantic environments are registered once here, not re-designed in manuscripts.
\newtheoremstyle{wzstatement}{6pt}{5pt}
  {\normalfont\kaishu}{0pt}{\normalfont\bfseries\color{structurecolor}}
  {}{.7em}{\thmname{#1}\thmnumber{ #2}\thmnote{\normalfont\ (#3)}}
\newtheoremstyle{wzdefinition}{6pt}{5pt}
  {\normalfont\songti}{0pt}{\normalfont\bfseries\color{structurecolor}}
  {}{.7em}{\thmname{#1}\thmnumber{ #2}\thmnote{\normalfont\ (#3)}}
\newtheoremstyle{wzproblem}{7pt}{5pt}
  {\normalfont\songti}{2em}{\normalfont\bfseries\color{main}}
  {}{.7em}{\thmname{#1}\thmnumber{ #2}}
\newtheoremstyle{wzremark}{4pt}{4pt}
  {\normalfont\songti}{0pt}{\normalfont\bfseries}
  {}{.7em}{\thmname{#1}}
\theoremstyle{wzstatement}
\newtheorem{theorem}{定理}[chapter]
\newtheorem{lemma}[theorem]{引理}
\newtheorem{proposition}[theorem]{命题}
\newtheorem{corollary}[theorem]{推论}
\theoremstyle{wzdefinition}
\newtheorem{definition}[theorem]{定义}
\theoremstyle{wzproblem}
\newtheorem{example}{例题}[chapter]
\newtheorem{exercise}{习题}[chapter]
\theoremstyle{wzremark}
\newtheorem*{remark}{注}
\renewcommand{\proofname}{证明}
% Retain the native amsthm QED stack; only adjust the fixed typography.
\expandafter\patchcmd\csname\string\proof\endcsname{\normalfont \topsep6\p@\@plus6\p@\relax}
  {\normalfont\songti \topsep5\p@\@plus1\p@\relax}
  {}{\ClassError{elegantbook-original-adapter}{Cannot patch proof spacing}{Check the amsthm version.}}
\expandafter\patchcmd\csname\string\proof\endcsname{\itshape}{\normalfont\bfseries}
  {}{\ClassError{elegantbook-original-adapter}{Cannot patch proof heading}{Check the amsthm version.}}
\expandafter\patchcmd\csname\string\proof\endcsname{\ignorespaces}
  {\setlength{\parindent}{2em}\ignorespaces}
  {}{\ClassError{elegantbook-original-adapter}{Cannot patch proof paragraphs}{Check the amsthm version.}}
\AtBeginEnvironment{proof}{\par\Needspace{3\baselineskip}}
\renewcommand{\qedsymbol}{\ensuremath{\square}}
\newenvironment{solution}{\begin{proof}[解答]}{\end{proof}}
\newcommand{\wzcheckchapter}{\ifnum\value{chapter}<1
  \ClassError{elegantbook-original-adapter}{Numbered mathematics before chapter 1}{Move it after the first originalchapter.}\fi}

\newcommand{\wzregisterguard}[1]{\AtBeginEnvironment{#1}{\wzcheckchapter\par\Needspace{3\baselineskip}}}
\forcsvlist{\wzregisterguard}{theorem,lemma,proposition,corollary,definition,example,exercise}
% Front matter and bibliography are not numbered mathematical chapters.
\newcommand{\originalpreface}{\par\addvspace{10pt}\Needspace{5\baselineskip}
  \noindent{\bfseries\fontsize{12}{17}\selectfont 前言}\par\nobreak\vspace{4pt}}
\newcommand{\originalcontents}{\par\addvspace{14pt}
  \begin{center}{\color{structurecolor}\bfseries\fontsize{14.4}{20}\selectfont 目录}\end{center}
  \vspace{-10pt}\@starttoc{toc}}
\renewcommand{\bibname}{参考文献}
\renewenvironment{thebibliography}[1]{
  \par\addvspace{12pt}\Needspace{3\baselineskip}\phantomsection
  \addcontentsline{toc}{chapter}{\bibname}
  \markboth{\bibname}{}
  \noindent{\color{structurecolor}\bfseries\fontsize{14.4}{20}\selectfont\bibname}\par\nobreak\vspace{6pt}
  \list{\@biblabel{\@arabic\c@enumiv}}
    {\settowidth\labelwidth{\@biblabel{#1}}\leftmargin\labelwidth
     \advance\leftmargin\labelsep\usecounter{enumiv}
     \let\p@enumiv\@empty\renewcommand\theenumiv{\@arabic\c@enumiv}
     \setlength{\itemsep}{3pt}\setlength{\parsep}{0pt}}
  \sloppy\clubpenalty4000\widowpenalty4000\sfcode`\.=1000\relax
}{\def\@noitemerr{\@latex@warning{Empty bibliography}}\endlist}

\numberwithin{equation}{chapter}
\renewcommand{\theequation}{\thechapter.\arabic{equation}}
\allowdisplaybreaks[2]
\emergencystretch=2em
\clubpenalty=6000
\widowpenalty=6000
\raggedbottom
\endinput
\end{filecontents*}
% WZ-LOCKED-CLASS-END
\documentclass{elegantbook-original-adapter}
% ===== 元数据内容槽 =====
\newcommand{\BookTitle}{测度与积分}
\newcommand{\BookAuthor}{Jeff Viaclovsky 讲授; Ethan Brown 原讲义排录}
\newcommand{\BookDate}{MIT 18.125, Fall 2003}
\hypersetup{pdftitle={测度与积分: MIT 18.125 中文重编本},pdfauthor={Jeff Viaclovsky; Ethan Brown; 中文译编}}
\usepackage{booktabs,longtable}
\newcommand{\R}{\mathbb R}
\newcommand{\N}{\mathbb N}
\newcommand{\Z}{\mathbb Z}
\newcommand{\C}{\mathbb C}
\newcommand{\Q}{\mathbb Q}
\newcommand{\eps}{\varepsilon}
\newcommand{\dd}{\,\mathrm d}
\newcommand{\ind}{\mathbf 1}
\newcommand{\supp}{\operatorname{supp}}
\newcommand{\norm}[1]{\lVert #1\rVert}
\newcommand{\abs}[1]{\lvert #1\rvert}
\newcommand{\ip}[2]{\langle #1,#2\rangle}
\newcommand{\sourcelecture}[1]{\par\noindent 对应原课程第 #1 讲.\par}
\begin{document}
\frontmatter
\begin{center}
{\color{structurecolor}\bfseries\fontsize{18}{25}\selectfont\BookTitle\par}
{\bfseries MIT 18.125 中文重编本\par}
\BookAuthor\par
\BookDate\par
\end{center}
\originalpreface
集合的大小、函数的积分和极限交换是同一个理论的不同层次. 测度给可测集合指定大小, 可测函数把集合结构带到函数上, Lebesgue 积分由简单函数的逼近建立. 本册沿这一依赖顺序, 从一般测度空间进入欧氏空间, 继而讨论函数空间、乘积积分、卷积和几乎处处微分. 读者需掌握数列与级数、连续函数、紧性以及一元和多元微积分; 本册在用到时说明所需的集合与拓扑语言.

主要依据为 Jeff Viaclovsky 在 MIT 讲授的 \href{https://ocw.mit.edu/courses/18-125-measure-and-integration-fall-2003/pages/lecture-notes/}{18.125 Measure and Integration, Fall 2003} 的24份公开讲义, 原稿由 Ethan Brown 根据 Viaclovsky 的手写课堂笔记排成 LaTeX. 已实际下载的24份文件合计95个PDF物理页. 本册用简体中文重新编排, 合并复述, 补齐原稿省略的论证并指出影响数学结论的排印错误. 例题和习题的解答均在正文中展开; 编者增设的内容不冒称课程原文或官方答案.

第1至8章覆盖这24份讲义的主题与主要论证. 第9章为显式编者补充的 Radon--Nikodym 定理与 Lebesgue 分解, 并补证原稿只陈述的绝对连续逆向积分定理, 用来解释积分密度并接续泛函分析中的 $L^p$ 对偶. 一般赋范空间的算子理论、谱理论和一般对偶定理由泛函册承担. 官方教学大纲另列 Hausdorff 测度、面积公式与余面积公式, 但这24份下载讲义中没有相应正文; 本册不宣称已翻译这些不存在于源文件的部分. 卷首补入官方课程参考、课历及三组作业题号安排; 原书题面与官方解答未纳入, 正文33道配套练习仍为编者题.

原课程、作者、年份与逐文件下载链接见书末及 sources/real-analysis/source-manifest.json. 来源文件逐件核实页数与校验和; 没有发现单独第三方权利声明或复用外部图片. 本册正文采用原生 LaTeX, 不以源页截图代替文字. 中文译编、重排、补证、勘误与新增练习属于本次改动. 原课程与本改编均遵循 \href{https://creativecommons.org/licenses/by-nc-sa/4.0/}{CC BY-NC-SA 4.0 International}; 单独权利声明优先. 署名不表示 MIT 或原作者认可本项目. 许可依据见 \href{https://ocw.mit.edu/pages/privacy-and-terms-of-use/}{MIT OCW 使用条款}. 版本日期: 2026-10-08.
\clearpage
\phantomsection
\addcontentsline{toc}{chapter}{课程参考}
\markboth{课程参考}{}
\begin{center}
{\color{structurecolor}\bfseries\fontsize{14.4}{20}\selectfont 课程参考\par}
\end{center}

\noindent\textbf{课程与教师.} MIT 18.125 \textit{Measure and Integration}, 2003年秋季, 研究生课程; 教师 Jeff Viaclovsky. 公开讲义由 Ethan Brown 根据教师的课堂手写笔记排录. 以下课程信息译自官方\href{https://ocw.mit.edu/courses/18-125-measure-and-integration-fall-2003/pages/syllabus/}{教学大纲}, 核验日期为2026-10-08.

\noindent\textbf{授课时间与先修.} 每周2次课, 每次1.5小时. 官方先修为 Analysis I (18.100). 核实的页面未给出具体星期、钟点或逐次课的公历日期; 后列课历采用官方讲次.

\noindent\textbf{大纲.} 从抽象测度论与 Lebesgue 积分入门, 先建立积分及主要收敛定理, 再构造 $\R^n$ 上的 Lebesgue 测度. 大纲还列出 $L^p$ 空间、Radon--Nikodym 定理、Lebesgue 微分定理、Fubini 定理、Hausdorff 测度以及面积与余面积公式. 本册实际正文范围仍以24份讲义及书末来源对照为准: Radon--Nikodym 是编者补充, Hausdorff 测度与面积、余面积公式未纳入. 官网课程简介另提 Fourier 变换, 但这24份原件没有对应正文.

\noindent\textbf{教材.} 必读: Walter Rudin, \textit{Real and Complex Analysis}, McGraw-Hill International Editions: Mathematics Series, McGraw-Hill Education -- Europe, 1986, ISBN 9780070542341. 推荐: Frank Jones, \textit{Lebesgue Integration on Euclidean Space}, Jones \& Bartlett, 1993; Lawrence C. Evans 与 Ronald F. Gariepy, \textit{Measure Theory and Fine Properties of Function}, CRC Press, 1991. 这里只记录官方书目信息, 不归档外部教材内容.

\noindent\textbf{考核.} 官方大纲明确\textbf{没有考试}. 课程成绩以到课及提交作业为依据, 页面未给出比例. 作业页公开三组教材题号, 截止讲次分别为8、12、20; 完整索引见后页.

\clearpage
\phantomsection
\addcontentsline{toc}{chapter}{官方课历索引}
\markboth{官方课历索引}{}
\begin{center}
{\color{structurecolor}\bfseries\fontsize{14.4}{20}\selectfont 官方课历索引\par}
\end{center}
下表据\href{https://ocw.mit.edu/courses/18-125-measure-and-integration-fall-2003/pages/calendar/}{官方课历}译录讲次、主题与作业截止安排. 课历主题与实际讲义的对应并不完全一致; 本表记录课堂安排, 逐文件覆盖见书末.

\begin{longtable}{@{}p{.07\textwidth}p{.75\textwidth}p{.12\textwidth}@{}}
\toprule
讲次 & 官网主题 & 截止作业 \\
\midrule
\endfirsthead
\toprule
讲次 & 官网主题 & 截止作业 \\
\midrule
\endhead
\bottomrule
\endfoot
1 & 测度论的动机; 测度空间与 $\sigma$-代数; 可测函数的和、积、复合; Borel 集 & \\
2 & 实值可测函数及极限; 简单函数、正测度、Lebesgue 积分定义 & \\
3 & Riemann 积分; 可积与几乎处处连续; 两种积分的比较; 正测度与积分的性质 & \\
4 & 简单函数与非负函数积分可加; 单调收敛; 求和与积分交换; Fatou 引理 & \\
5 & 复函数积分、主导收敛、零测集、$\sigma$-代数的完备化 & \\
6 & $\R^n$ 上的 Lebesgue 测度; 矩形、多面集、开集、紧集; 内外测度 & \\
7 & 从有限外测度可测集到一般可测集; Lebesgue 可测域与测度空间 & \\
8 & Carath\'eodory 判据、Cantor 集、Lebesgue 可测而非 Borel 的集合 & 1 \\
9 & 平移、伸缩与旋转不变性; 不可测集 & \\
10 & 积分作为线性泛函; 正泛函 Riesz 表示; Riemann 与 Lebesgue 积分的衔接 & \\
11 & Lusin 定理; Vitali--Carath\'eodory 定理 & \\
12 & 连续函数逼近可测函数; 几乎处处收敛下的积分定理; 积分可数可加 & 2 \\
13 & Egorov 定理; 依测度与几乎处处收敛; 子列; 依测度主导收敛 & \\
14 & 凸函数、Jensen、H\"older 与 Minkowski 不等式 & \\
15 & $L^p$、赋范与 Banach 空间; Riesz--Fischer 完备性 & \\
16 & $C_c$ 在有限指数 $L^p$ 与 $C_0$ 中稠密 & \\
17 & $L^p$ 与 $\ell^p$ 包含; 局部 $L^p$; 范数凸性; 平滑函数稠密 & \\
18 & $\R^n$ 上非负函数的 Fubini 定理 & \\
19 & $\R^n$ 上 $L^1$ 函数的 Fubini 定理; 一般测度空间的乘积测度 & \\
20 & 乘积测度的 Fubini 定理、完备化与卷积 & 3 \\
21 & Young 不等式、磨光核、$C^\infty$ 在有限指数 $L^p$ 中稠密 & \\
22 & Lebesgue 积分的微积分基本定理; Vitali 覆盖; 最大函数; 弱 $L^1$ 估计 & \\
23 & Lebesgue 微分定理、Lebesgue 点与微积分基本定理第一部分 & \\
24 & 广义 Minkowski、另一 Young 证明、分布函数; 最大算子 $L^p$ 插值估计 & \\
\end{longtable}

\clearpage
\phantomsection
\addcontentsline{toc}{chapter}{官方作业安排索引}
\markboth{官方作业安排索引}{}
\begin{center}
{\color{structurecolor}\bfseries\fontsize{14.4}{20}\selectfont 官方作业安排索引\par}
\end{center}
官方\href{https://ocw.mit.edu/courses/18-125-measure-and-integration-fall-2003/pages/assignments/}{作业页}只列三组 Rudin 教材题号与截止讲次. 教材为前列1986年 \textit{Real and Complex Analysis}; 下表的章号、题号均指该教材, 不指本册章号或33道编者练习.

\begin{center}
\begin{tabular}{@{}p{.12\textwidth}p{.15\textwidth}p{.12\textwidth}p{.48\textwidth}@{}}
\toprule
作业 & 截止讲次 & 原书章号 & 官方所列题号 \\
\midrule
1 & 8 & 1 & 3--12 \\
2 & 12 & 2 & 5--9, 20 \\
3 & 20 & 3 & 4c, 4d, 10, 11, 13, 18b, 23, 25 \\
\bottomrule
\end{tabular}
\end{center}

本次核实的公开作业页没有提供完整题面、独立作业 PDF 或官方解答. 本册因此仅收入这份安排索引, \textbf{没有收入原书题面或官方解答}, 也不据题号重建题意. 正文保留的33道配套练习及解答为编者内容, 未改称官方作业. 原课程明确无考试, 因而没有需要补入的本课试卷.

这些课程信息与安排索引保留 MIT OCW、课程、教师及年份署名, 依 OCW 使用条款与 CC BY-NC-SA 4.0 使用. 教材仅作书目与题号定位, 不将 OCW 的许可扩展为外部教材题面的转载许可.

\clearpage
\markboth{}{}
\originalcontents
\clearpage
\mainmatter
\originalchapter{测度空间与可测函数}
\sourcelecture{1--3}

积分需要先知道集合有多大. 在 $\R^n$ 中, 长度、面积与体积给出了熟悉的答案; 在一般集合上, 我们仍希望把互不重叠的部分的大小相加. 但是, 为了使可数次分割和极限过程都能保持这一性质, 通常不能给每个子集同时规定合适的大小. 本章先确定可以测量的集合, 再确定可以积分的函数. 数值的具体来源暂不限定: 计数、一个点上的单位质量, 都可以作为测度.

\originalsection{\texorpdfstring{$\sigma$}{sigma}-代数与 Borel 集}
\begin{definition}\label{ra:c1:sigma}
设 $X$ 是集合. 其幂集 $\mathcal P(X)$ 是 $X$ 的所有子集组成的集合. 集合族 $\mathcal M\subset\mathcal P(X)$ 称为 $X$ 上的 $\sigma$-代数, 如果 $X\in\mathcal M$, 且满足下列封闭性: 若 $A\in\mathcal M$, 则 $X\setminus A\in\mathcal M$; 若 $A_n\in\mathcal M$ 对所有 $n\in\N$ 成立, 则 $\bigcup_{n=1}^{\infty}A_n\in\mathcal M$. 二元组 $(X,\mathcal M)$ 称为可测空间, $\mathcal M$ 中的集合称为可测集.
\end{definition}
这里的补集始终相对于 $X$. 由 $\varnothing=X^c$ 可知空集可测. De Morgan 法则给出 $\bigcap_n A_n=(\bigcup_n A_n^c)^c$, 所以可数交也可测. 有限并、有限交可以看成补上空集或 $X$ 的可数运算, 因而同样封闭; 差集 $A\setminus B=A\cap B^c$ 也可测.

拓扑与 $\sigma$-代数承担不同任务. 拓扑要求任意并和有限交仍为开集, 却不要求补集为开集; $\sigma$-代数要求补集和可数并仍可测, 却不要求不可数并仍可测. 因此, 从开集出发, 还要加入补集、可数并及它们反复产生的集合, 才能得到用于积分的集合族.

\begin{proposition}\label{ra:c1:generated}
对任意集合族 $\mathcal F\subset\mathcal P(X)$, 存在唯一最小的、包含 $\mathcal F$ 的 $\sigma$-代数, 记为 $\sigma(\mathcal F)$.
\end{proposition}
\begin{proof}
令 $\mathfrak S$ 是所有包含 $\mathcal F$ 的 $\sigma$-代数组成的族. 因为 $\mathcal P(X)$ 本身属于 $\mathfrak S$, 这个族非空. 取 $\mathcal M_* =\bigcap_{\mathcal A\in\mathfrak S}\mathcal A$. 每个 $\mathcal A$ 都含 $X$, 所以 $X\in\mathcal M_*$. 若 $A\in\mathcal M_*$, 则 $A$ 属于每个 $\mathcal A$, 其补集也属于每个 $\mathcal A$, 故 $A^c\in\mathcal M_*$. 对 $A_n\in\mathcal M_*$, 在每个 $\mathcal A$ 中作可数并, 得 $\bigcup_n A_n\in\mathcal M_*$. 因而 $\mathcal M_*$ 是 $\sigma$-代数. 它包含 $\mathcal F$, 且包含于任何包含 $\mathcal F$ 的 $\sigma$-代数中, 正是所求的最小者. 两个最小者必互相包含, 所以唯一.
\end{proof}

\begin{definition}
若 $X$ 是拓扑空间, 由其开集生成的 $\sigma$-代数记为 $\mathcal B(X)$, 称为 Borel $\sigma$-代数. 其中的集合称为 Borel 集. 可数个闭集的并称为 $F_\sigma$ 集, 可数个开集的交称为 $G_\delta$ 集.
\end{definition}
闭集是开集的补集, 因而闭集、$F_\sigma$ 集和 $G_\delta$ 集都为 Borel 集. 在 $\R$ 中, 单点是闭集, 所以每个可数集都是 $F_\sigma$ 集. 例如 $\mathbb Q$ 是 Borel 集, 而 $\R\setminus\mathbb Q$ 是 $G_\delta$ 集. Borel 集的定义还允许更复杂的交并过程, 并不只包括上述两类.

\begin{proposition}[编者补充]\label{ra:c1:countable-base}
$\R^d$ 中所有端点为有理数的开矩形构成可数拓扑基. 因此每个开集是可数个这样的矩形的并, 而 $\mathcal B(\R)$ 可以由开区间、由有理端点开区间, 或由所有射线 $(a,\infty)$ 生成.
\end{proposition}
\begin{proof}
有理数集合可数, 有限个有理数的有序组合也可数, 故这些矩形可数. 若 $x$ 属于开集 $U$, 则存在以 $x$ 为中心且包含于 $U$ 的小球. 在每个坐标上选取略小于和略大于 $x$ 坐标的有理数, 可以使所得矩形含 $x$ 且包含于该小球. 因而 $U$ 是所有包含于 $U$ 的有理矩形的并; 这些矩形是一个可数族的子族.

在一维情形, 这证明有理开区间生成全部开集. 射线本身开, 其生成的 $\sigma$-代数包含于 Borel $\sigma$-代数. 反过来, $[a,\infty)=\bigcap_{n=1}^{\infty}(a-1/n,\infty)$, 所以 $(-\infty,a)$ 是可测的补集. 再取交得到 $(a,b)=(a,\infty)\cap(-\infty,b)$, 故射线也生成全部开区间及 Borel 集.
\end{proof}

\originalsection{可测映射与连续运算}
\begin{definition}\label{ra:c1:measurable}
设 $(X,\mathcal M)$ 为可测空间, $Y$ 为拓扑空间. 映射 $f:X\to Y$ 称为可测, 如果对每个开集 $U\subset Y$ 都有 $f^{-1}(U)\in\mathcal M$. 当 $X$ 也是拓扑空间且 $\mathcal M=\mathcal B(X)$ 时, 称 $f$ 为 Borel 可测映射.
\end{definition}
原像 $f^{-1}(U)=\{x\in X:f(x)\in U\}$ 不要求 $f$ 有逆映射. 它保持并、交和补集, 这是可测性能够与集合运算相配合的原因. 连续函数的开集原像仍开, 因而连续函数必为 Borel 可测函数.

\begin{proposition}\label{ra:c1:preimage-borel}
设 $f:X\to Y$. 集合族 $\mathcal A_f=\{E\subset Y:f^{-1}(E)\in\mathcal M\}$ 是 $Y$ 上的 $\sigma$-代数. 如果 $f$ 可测, 则每个 Borel 集的原像都可测. 如果 $f:X\to Y$ 可测而 $g:Y\to Z$ 连续, 则 $g\circ f$ 可测.
\end{proposition}
\begin{proof}
因为 $f^{-1}(Y)=X$, 所以 $Y\in\mathcal A_f$. 恒等式 $f^{-1}(Y\setminus E)=X\setminus f^{-1}(E)$ 和 $f^{-1}(\bigcup_n E_n)=\bigcup_n f^{-1}(E_n)$ 给出其补集与可数并封闭性. 如果 $f$ 可测, $\mathcal A_f$ 包含全部开集, 因而包含它们生成的 $\mathcal B(Y)$. 最后, 对 $Z$ 中的开集 $U$, 连续性给出 $g^{-1}(U)$ 在 $Y$ 中开, 所以 $(g\circ f)^{-1}(U)=f^{-1}(g^{-1}(U))\in\mathcal M$.
\end{proof}
同一原像恒等式也证明两个连续映射的复合连续. 若把 $g$ 的假设改成 Borel 可测, 结论仍成立: 此时先用本命题知道 $f$ 对 Borel 集的原像可测, 再代入 $g^{-1}(U)$ 即可. 对一般可测空间间的复合, 必须明确中间空间的 $\sigma$-代数; 不能只省略这些结构说两个任意可测映射可以复合.

\begin{theorem}\label{ra:c1:continuous-operations}
设 $u_1,\ldots,u_d:X\to\R$ 可测, $\Phi:\R^d\to Y$ 连续. 则 $x\mapsto\Phi(u_1(x),\ldots,u_d(x))$ 可测.
\end{theorem}
\begin{proof}
先证 $F:X\to\R^d$, $F(x)=(u_1(x),\ldots,u_d(x))$ 可测. 开矩形 $I_1\times\cdots\times I_d$ 的原像为 $\bigcap_{j=1}^{d}u_j^{-1}(I_j)$, 是可测集. 任意开集由命题~\ref{ra:c1:countable-base} 写成可数个开矩形的并, 其原像也就是可数个可测集的并. 故 $F$ 可测. 命题~\ref{ra:c1:preimage-borel} 再给出 $\Phi\circ F$ 可测.
\end{proof}
证明中需要的是\emph{可数}拓扑基. 若用每个点周围的矩形形成不可数并, $\sigma$-代数的公理并不足以保证该原像可测.

\begin{corollary}\label{ra:c1:complex}
复函数 $f=u+iv:X\to\C$ 可测当且仅当实部 $u$ 与虚部 $v$ 都可测. 此时 $\abs f$ 也可测. 两个实值或复值可测函数的和、差、积可测, 可测函数乘常数仍可测. 在 $g$ 处处非零时, $f/g$ 也可测.
\end{corollary}
\begin{proof}
若 $u,v$ 可测, 用连续映射 $(a,b)\mapsto a+ib$ 和定理~\ref{ra:c1:continuous-operations}. 若 $f$ 可测, 分别复合连续函数 $z\mapsto\operatorname{Re}z$, $z\mapsto\operatorname{Im}z$ 和 $z\mapsto\abs z$, 就得到 $u,v,\abs f$ 可测. 对实函数的四则运算使用相应连续映射; 对复函数, 把实部与虚部运算展开, 再用实函数的结论. 除法中 $z\mapsto1/z$ 在 $\C\setminus\{0\}$ 连续; 把 $g$ 看成到此子空间的可测映射即可.
\end{proof}

\begin{example}
设 $E\subset X$. 证明示性函数 $\mathbf 1_E$, 在 $E$ 上取 $1$, 在 $E^c$ 上取 $0$, 可测当且仅当 $E\in\mathcal M$.
\end{example}
\begin{solution}
若函数可测, 则 $E=\mathbf 1_E^{-1}((1/2,3/2))$ 可测. 反之, 一个开集的原像只能是 $\varnothing,E,E^c,X$ 之一, 这四个集合都可测.
\end{solution}

\originalsection{扩展实值函数}
为了描述非负级数和函数的极限, 需要允许无穷值. 记 $\overline\R=[-\infty,\infty]$, 赋予序拓扑: 普通开区间、$[-\infty,a)$ 和 $(b,\infty]$ 构成拓扑基. 收敛到 $+\infty$ 是指最终大于任意实数, 收敛到 $-\infty$ 类似. 非负量的和与乘积允许无穷值, 并约定 $0\cdot\infty=0$; 但 $\infty-\infty$ 没有定义, 不可参与代数消去.

\begin{proposition}\label{ra:c1:ray-test}
映射 $f:X\to\overline\R$ 可测当且仅当对每个 $a\in\R$ 都有 $\{f>a\}\in\mathcal M$. 可以把这个条件等价地换成 $\{f\ge a\}$、$\{f<a\}$ 或 $\{f\le a\}$ 对每个实数 $a$ 可测. 只检验有理数阈值也足够.
\end{proposition}
\begin{proof}
必要性来自 $(a,\infty]$ 为开集. 反过来, 假定所有 $\{f>a\}$ 可测. 则
\[
 \{f\ge a\}=\bigcap_{n=1}^{\infty}\{f>a-1/n\},\qquad
 \{f<a\}=\bigcup_{n=1}^{\infty}\{f\le a-1/n\}.
\]
其中 $\{f\le t\}=X\setminus\{f>t\}$. 因而各型射线原像可测, 有界开区间的原像为 $\{f>a\}\cap\{f<b\}$, 也可测. $\overline\R$ 的每个开集是可数个上述基元素的并, 故 $f$ 可测. 这些恒等式及取补集同样证明四种判据等价. 若只知道有理阈值的严格上水平集可测, 用 $\{f>a\}=\bigcup_{q\in\mathbb Q,\,q>a}\{f>q\}$ 即可恢复全部实阈值. 其余型式可先以有理阈值的交并转换成严格上水平集.
\end{proof}
无穷值的水平集也可测, 因为 $\{f=+\infty\}=\bigcap_n\{f>n\}$, $\{f=-\infty\}=\bigcap_n\{f<-n\}$. 对有限实值函数, 本命题就是通常的实值可测性判据.

\begin{remark}
原讲义在扩展实数的若干射线判据中采用了不同的端点写法. 本书统一保留 $\pm\infty$ 端点; 对一般扩展实值函数, 若仅检验有限区间而未控制无穷值集合, 不能直接跳过这些端点.
\end{remark}

\originalsection{可测函数的极限}
对序列 $a_n\in\overline\R$, 定义 $b_k=\sup_{n\ge k}a_n$, $c_k=\inf_{n\ge k}a_n$. 尾部缩短时, $b_k$ 递减而 $c_k$ 递增, 因此在扩展实数中总有极限. 定义
\[
 \limsup_{n\to\infty}a_n=\inf_{k\ge1}\sup_{n\ge k}a_n,\qquad
 \liminf_{n\to\infty}a_n=\sup_{k\ge1}\inf_{n\ge k}a_n.
\]
总有 $\liminf a_n\le\limsup a_n$. 两者相等于 $L$ 当且仅当 $a_n\to L$. 对有限 $L$, 因为 $c_k\to L$ 与 $b_k\to L$, 给定 $\eps>0$ 可取 $k$ 使 $L-\eps<c_k\le a_n\le b_k<L+\eps$ 对所有 $n\ge k$ 成立. 若两者均为 $+\infty$, 则 $c_k\to+\infty$ 迫使尾部最终大于任何实数; $-\infty$ 的情形用 $b_k$ 处理. 反向则直接由收敛的定义与尾部上下界得到.

上下极限还分别是序列最大的和最小的子列极限, 在这里允许无穷极限. 例如当上极限 $L$ 有限时, 一方面整个尾部最终小于 $L+1/j$, 另一方面任意尾部均有一项大于 $L-1/j$; 递归选取严格增加的下标便得到趋于 $L$ 的子列. 当 $L=+\infty$ 时, 从各尾部取大于 $j$ 的项; 当 $L=-\infty$ 时, 原序列已经趋于 $-\infty$. 下极限的结论可对 $-a_n$ 使用同一论证.

\begin{theorem}\label{ra:c1:limits}
若每个 $f_n:X\to\overline\R$ 可测, 则 $\sup_n f_n$, $\inf_n f_n$, $\limsup_n f_n$ 与 $\liminf_n f_n$ 都可测. 因而处处存在的逐点极限可测. 有限个可测函数的最大值和最小值也可测.
\end{theorem}
\begin{proof}
令 $g(x)=\sup_n f_n(x)$. 对每个实数 $a$, 有 $\{g>a\}=\bigcup_n\{f_n>a\}$: 上确界严格大于 $a$ 时, 必有一个成员严格大于 $a$. 命题~\ref{ra:c1:ray-test} 给出 $g$ 可测. 取负号得到下确界可测. 对每个 $k$, 尾部上确界 $\sup_{n\ge k}f_n$ 可测, 再取这些函数的下确界, 得上极限可测; 下极限同理. 若逐点极限存在, 它等于上下极限, 因而可测. 最大值和最小值是有限族的上、下确界.
\end{proof}
对于复值函数的逐点极限, 分别对实部和虚部应用本定理, 再用推论~\ref{ra:c1:complex}. 原讲义曾把函数上极限写成通常极限; 上式给出了即使函数列不收敛也有意义的定义.

\begin{definition}
对实值或扩展实值函数 $f$, 定义正部 $f^+=\max(f,0)$ 与负部 $f^-=\max(-f,0)$. 若 $f$ 可测, 则二者非负且可测. 在 $f$ 有限时, $f=f^+-f^-$, $\abs f=f^++f^-$; 扩展实值情形也按逐点约定成立, 因为同一点的两部分不会同时为无穷.
\end{definition}

\originalsection{简单函数逼近}
\begin{definition}
只取有限多个有限实数值的函数称为简单函数. 若其不同取值为 $\alpha_1,\ldots,\alpha_r$, 取 $A_j=\{s=\alpha_j\}$, 则 $s=\sum_{j=1}^{r}\alpha_j\mathbf1_{A_j}$. 这些 $A_j$ 两两不交且并为 $X$, 称为函数的水平集表示. 简单函数可测当且仅当其每个水平集可测.
\end{definition}
最后的等价性可以直接从原像验证: 任意开集的原像是其中若干水平集的并; 反过来, 用足够小的开区间隔离每个不同取值, 就恢复对应水平集的原像. 复值简单函数可作完全相同的定义, 但非负积分首先只用非负实值简单函数.

\begin{theorem}\label{ra:c1:simple-approx}
若 $f:X\to[0,\infty]$ 可测, 则存在非负可测简单函数列 $s_n$, 满足 $0\le s_1\le s_2\le\cdots\le f$ 且 $s_n(x)\to f(x)$ 对每个 $x$ 成立.
\end{theorem}
\begin{proof}
将 $[0,n)$ 按长度 $2^{-n}$ 分成有限多个半开区间, 定义
\[
 \phi_n(t)=
 \begin{cases}
  2^{-n}\lfloor2^nt\rfloor,&0\le t<n,\\
  n,&n\le t\le\infty.
 \end{cases}
\]
$\phi_n$ 只取 $0,2^{-n},\ldots,n$ 中的有限多个值, 且各水平集为 Borel 集, 所以 Borel 可测. 令 $s_n=\phi_n\circ f$, 命题~\ref{ra:c1:preimage-borel} 保证 $s_n$ 可测.

逐点检验单调性. 当 $t<n$ 时, 二进制网格细分给出 $\lfloor2^{n+1}t\rfloor\ge2\lfloor2^nt\rfloor$, 故 $\phi_{n+1}(t)\ge\phi_n(t)$. 当 $n\le t<n+1$ 时, $n$ 本身是更细网格的端点, 所以 $\phi_{n+1}(t)\ge n=\phi_n(t)$. 当 $t\ge n+1$ 时, $\phi_{n+1}(t)=n+1\ge n=\phi_n(t)$. 每个 $\phi_n(t)\le t$. 对有限 $t$, 一旦 $n>t$, 就有 $0\le t-\phi_n(t)<2^{-n}$; 对 $t=\infty$, 有 $\phi_n(t)=n\to\infty$. 因而 $s_n\uparrow f$.
\end{proof}
这个构造同时作两件事: 截去很大的函数值, 再把剩下的值向下舍入到越来越细的网格. 有限截断保证每步是简单函数, 向下舍入保证始终不超过 $f$, 嵌套网格保证单调性. 后一性质正是下一章交换积分与极限的依据.

\originalsection{正测度及连续性}
\begin{definition}\label{ra:c1:measure}
可测空间 $(X,\mathcal M)$ 上的正测度是映射 $\mu:\mathcal M\to[0,\infty]$, 满足 $\mu(\varnothing)=0$, 以及对两两不交的可测集 $A_n$ 的可数可加性
\[
 \mu\left(\bigcup_{n=1}^{\infty}A_n\right)=\sum_{n=1}^{\infty}\mu(A_n).
\]
三元组 $(X,\mathcal M,\mu)$ 称为测度空间. 若 $\mu(X)<\infty$, 称为有限测度; 若 $\mu(X)=1$, 称为概率测度. 若 $X$ 可写成可数个有限测度可测集的并, 称 $\mu$ 为 $\sigma$-有限测度.
\end{definition}
正测度允许某些集合的测度为 $\infty$, 但不允许空集的测度为 $\infty$. 原讲义只先写可数可加性, 再从一个有限测度集合推出空集测度为零; 本书将 $\mu(\varnothing)=0$ 明确列入定义, 以排除恒取无穷的退化映射.

\begin{example}
在 $\mathcal P(X)$ 上定义计数测度 $\mu(E)=\#E$ 当 $E$ 有限, $\mu(E)=\infty$ 当 $E$ 无限. 固定 $x_0\in X$, 定义 Dirac 测度 $\delta_{x_0}(E)=\mathbf1_E(x_0)$. 验证二者为测度.
\end{example}
\begin{solution}
两者都给空集赋值零. 对互不相交的集合列, 若并集有限, 只有有限多个集合非空, 计数相加给出等式. 若并集无限而某个成员无限, 两边都为无穷; 若所有成员有限, 非空成员必有无限多个, 其正整数计数之和为无穷. 因而计数测度可数可加. 对 Dirac 测度, 不交集合中至多一个含 $x_0$, 且并集含 $x_0$ 当且仅当某个成员含它, 所以同样可数可加. Dirac 测度是概率测度; $\N$ 上的计数测度为 $\sigma$-有限, 因为 $\N=\bigcup_n\{1,\ldots,n\}$.
\end{solution}

\begin{proposition}\label{ra:c1:measure-properties}
设 $\mu$ 是正测度.
\begin{enumerate}
\item 若 $A\subset B$ 均可测, 则 $\mu(A)\le\mu(B)$. 若另有 $\mu(A)<\infty$, 则 $\mu(B\setminus A)=\mu(B)-\mu(A)$.
\item 对任意可测集列 $A_n$, 有 $\mu(\bigcup_n A_n)\le\sum_n\mu(A_n)$.
\item 若 $A_n\uparrow A$, 即 $A_1\subset A_2\subset\cdots$ 且 $A=\bigcup_n A_n$, 则 $\mu(A_n)\uparrow\mu(A)$.
\item 若 $A_n\downarrow A$, 即 $A_1\supset A_2\supset\cdots$ 且 $A=\bigcap_n A_n$, 并且 $\mu(A_1)<\infty$, 则 $\mu(A_n)\downarrow\mu(A)$.
\end{enumerate}
\end{proposition}
\begin{proof}
由 $B=A\mathbin{\dot\cup}(B\setminus A)$ 和有限可加性得 $\mu(B)=\mu(A)+\mu(B\setminus A)$. 这既给出单调性, 也在 $\mu(A)$ 有限时允许作减法.

令 $B_1=A_1$, $B_n=A_n\setminus\bigcup_{j<n}A_j$ 对 $n\ge2$. 则 $B_n$ 两两不交, 其并与 $A_n$ 的并相同, 且 $B_n\subset A_n$. 可数可加性与单调性给出第二项.

对于递增列, 改取 $B_1=A_1$, $B_n=A_n\setminus A_{n-1}$. 有 $A_n=\mathop{\dot\bigcup}_{j=1}^{n}B_j$ 与 $A=\mathop{\dot\bigcup}_{j=1}^{\infty}B_j$, 所以 $\mu(A_n)$ 是非负级数 $\sum_j\mu(B_j)$ 的第 $n$ 个部分和, 收敛到 $\mu(A)$.

对于递减列, 令 $C_n=A_1\setminus A_n$. 则 $C_n\uparrow A_1\setminus A$. 由前一项及 $\mu(A_1)<\infty$, 有 $\mu(A_1)-\mu(A_n)=\mu(C_n)\to\mu(A_1\setminus A)=\mu(A_1)-\mu(A)$. 所有被相减的量有限, 故得到所需结论.
\end{proof}

\begin{example}
在 $\N$ 上取计数测度, 令 $A_n=\{n,n+1,\ldots\}$. 这些集合递减, 交集为空, 但 $\mu(A_n)=\infty$ 对每个 $n$ 成立. 因而测度的向下连续性不能无条件删去有限测度假设.
\end{example}
\begin{solution}
每个整数只属于有限多个 $A_n$, 所以没有整数属于所有 $A_n$, 即 $\bigcap_n A_n=\varnothing$. 然而每个 $A_n$ 都无限, 故其测度恒为无穷, 与交集的测度零不同. 此例中 $\infty-\infty$ 的减法正是前一证明不能使用的步骤.
\end{solution}

\begin{definition}
可测集 $N$ 满足 $\mu(N)=0$ 时称为零测集. 一个性质若在某个可测零测集以外处处成立, 就称它几乎处处成立, 简记为 a.e. 可数个零测集的并仍为零测集, 这是可数次可加性或次可加性的直接结果.
\end{definition}
定义中特意要求例外集被一个\emph{可测}零测集覆盖. 一般测度空间的零测集子集未必可测; 下一章的完备化将专门处理这一点.

\originalsection{练习与解答}
以下题目及解答为编者补充.

\begin{exercise}
设 $\{E_j:j\in J\}$ 是 $X$ 的一个有限或可数分割, 即各块非空、两两不交且并为 $X$. 证明这些块的任意并组成一个 $\sigma$-代数, 它就是 $\sigma(\{E_j:j\in J\})$. 对此 $\sigma$-代数, 实值函数可测当且仅当在每个块上为常数.
\end{exercise}
\begin{solution}
这些并含 $X$, 补集是余下各块的并, 可数个这样的并仍是某些块的并, 故构成 $\sigma$-代数. 它含各块, 而因为块族可数, 每一个并都属于各块生成的 $\sigma$-代数, 所以两者相等. 若函数在各块上为常数, 任意开集的原像是一些完整块的并, 因而可测. 反过来, 若某块中有 $x,y$ 满足 $f(x)<f(y)$, 取两值之间的实数 $a$. 可测集 $\{f>a\}$ 含 $y$ 却不含 $x$, 不可能为完整块的并, 矛盾.
\end{solution}

\begin{exercise}
设 $f_n:X\to\R$ 可测. 证明 $f_n(x)$ 在 $\R$ 中收敛的点组成可测集. 注意这里不把趋于无穷视为有限实数收敛.
\end{exercise}
\begin{solution}
由实数的完备性, 实数序列收敛当且仅当为 Cauchy 序列. 因而所求集合为
\[
 C=\bigcap_{k=1}^{\infty}\ \bigcup_{N=1}^{\infty}\
    \bigcap_{m,n\ge N}\{x:\abs{f_m(x)-f_n(x)}<1/k\}.
\]
每个内层集合可测, 因为差及其绝对值可测. 双下标 $(m,n)$ 仍组成可数族, 所以上式只用了可数次并交, 给出 $C\in\mathcal M$. 这个表达式也明确区分了``对每个精度可选一个尾部''的量词次序.
\end{solution}

\begin{exercise}
设 $f:X\to[0,\infty]$ 可测. 验证简单逼近定理中的 $s_n$ 在每个 $\{f\le M\}$ 上, 当 $n>M$ 时满足 $\sup_{f\le M}\abs{f-s_n}\le2^{-n}$. 为什么不能据此断定在整个 $X$ 上一致收敛?
\end{exercise}
\begin{solution}
若 $f(x)\le M<n$, 就处于未截断的部分, 所以向下舍入的误差小于 $2^{-n}$. 对这一区域取上确界得所求估计. 整个 $X$ 上的函数可能无界, 而每个 $s_n$ 都不超过 $n$. 例如 $X=\N$, $f(k)=k$, 则任意固定 $n$ 都有 $\sup_k\abs{f(k)-s_n(k)}=\infty$. 对有限函数值的逐点逼近和对有界值域的一致逼近, 都不等于对任意无界函数的一致逼近.
\end{solution}

\begin{exercise}
设 $A_n\in\mathcal M$. 定义集合的上、下极限为 $\limsup_n A_n=\bigcap_N\bigcup_{n\ge N}A_n$ 与 $\liminf_n A_n=\bigcup_N\bigcap_{n\ge N}A_n$. 解释它们的逐点含义, 并证明其示性函数分别等于 $\limsup_n\mathbf1_{A_n}$ 与 $\liminf_n\mathbf1_{A_n}$.
\end{exercise}
\begin{solution}
点属于上极限, 是说无论从哪个下标以后看, 它总会再次出现, 即属于无穷多个 $A_n$. 点属于下极限, 是说从某个下标以后它始终出现, 即除有限多个外属于每个 $A_n$. 取值只有零与一的序列上极限为一, 恰好是有无穷多个一; 下极限为一, 恰好是最终恒为一. 这给出了所述两个示性函数等式, 也由可数并交直接证明了两个集合可测.
\end{solution}

\originalchapter{Lebesgue 积分与收敛定理}
\sourcelecture{2--5}

简单函数把值域分成有限多个层次, 每一层都有可测的原像. 它的积分只需把``函数值乘集合大小''相加. 一般非负函数则由所有不超过它的简单函数从下方逼近. 这个定义不预先要求函数有界, 也不要求整个空间有有限测度; 允许积分为无穷, 正是先建立非负积分的好处. 在此基础上, 再用正负部和实虚部定义可积函数的积分.

本章在一般测度空间 $(X,\mathcal M,\mu)$ 上进行. 只有最后与 Riemann 积分衔接的一节调用实直线上的 Lebesgue 测度; 它的构造在下一章给出, 本章其他结论不依赖这一构造.

\originalsection{简单函数的积分}
\begin{definition}\label{ra:c2:simple-integral}
设 $s=\sum_{j=1}^{r}\alpha_j\mathbf1_{A_j}$ 为非负可测简单函数, 其中 $\alpha_j\ge0$, $A_j$ 两两不交且并为 $X$. 对 $E\in\mathcal M$, 定义
\[
 \int_E s\,\dd\mu=\sum_{j=1}^{r}\alpha_j\mu(A_j\cap E),
 \qquad 0\cdot\infty=0.
\]
这里的积分允许为 $+\infty$.
\end{definition}
如果已经用所有不同取值的水平集表示 $s$, 这个表达式没有选择问题. 但计算时往往把水平集再分割, 因此还要检查不同表示给出同一积分.

\begin{proposition}\label{ra:c2:simple-properties}
简单函数积分与不交表示的选择无关. 对非负可测简单函数 $s,t$, 非负有限常数 $c$, 以及可测集 $E$, 有
\[
 \int_E(s+t)\,\dd\mu=\int_Es\,\dd\mu+\int_Et\,\dd\mu,
 \qquad \int_Ecs\,\dd\mu=c\int_Es\,\dd\mu.
\]
若 $s\le t$, 则 $\int_Es\,\dd\mu\le\int_Et\,\dd\mu$. 另外, 集合函数 $\nu_s(E)=\int_Es\,\dd\mu$ 是正测度.
\end{proposition}
\begin{proof}
设 $s$ 有两个不交表示 $\sum_i\alpha_i\mathbf1_{A_i}$ 和 $\sum_j\beta_j\mathbf1_{B_j}$. 取共同分割 $A_i\cap B_j$. 当这个交集非空时, $s$ 的值给出 $\alpha_i=\beta_j$; 交集为空时, 对积分没有贡献. 把每个 $A_i\cap E$ 分成有限个 $A_i\cap B_j\cap E$, 用测度的有限可加性, 可知两种积分都等于同一个双重有限和.

对 $s$ 和 $t$ 的水平集也取共同分割. 在每一块上, 和函数的值是两值之和, 因而积分的可加性归结为 $(\alpha_i+\beta_j)\mu(A_i\cap B_j\cap E)=\alpha_i\mu(A_i\cap B_j\cap E)+\beta_j\mu(A_i\cap B_j\cap E)$. 各量非负, 即使某项无穷, 这一等式仍合法. 齐次性同样逐项得到; $c=0$ 时两边均按约定为零. 若 $s\le t$, 在每个共同分割块上对应的常数满足同一不等式, 乘非负测度并求和便得到单调性.

最后, 若 $E_k$ 两两不交, 则
\begin{align*}
 \nu_s\left(\bigcup_k E_k\right)
 &=\sum_{j=1}^{r}\alpha_j\sum_{k=1}^{\infty}\mu(A_j\cap E_k)\\
 &=\sum_{k=1}^{\infty}\sum_{j=1}^{r}\alpha_j\mu(A_j\cap E_k)
 =\sum_{k=1}^{\infty}\nu_s(E_k).
\end{align*}
有限个非负和与一个非负级数交换次序, 可以先对 $k\le N$ 交换有限和, 再令 $N\to\infty$. 此外 $\nu_s(\varnothing)=0$, 故 $\nu_s$ 是正测度.
\end{proof}

\originalsection{非负积分}
\begin{definition}\label{ra:c2:positive-integral}
对可测 $f:X\to[0,\infty]$ 和 $E\in\mathcal M$, 定义其 Lebesgue 积分为
\[
 \int_E f\,\dd\mu
 =\sup\left\{\int_Es\,\dd\mu:
  s\text{ 为可测简单函数},\ 0\le s\le f\right\}.
\]
未注明积分区域时, 区域为 $X$. 符号 $\int_E f\,\dd\mu$ 也常写成 $\int_E f\,d\mu$.
\end{definition}
这个定义和上一节相容. 若 $f$ 本身简单, 单调性说明所有候选积分都不超过 $\int_Ef\,\dd\mu$, 而选取 $s=f$ 又使上确界达到它.

\begin{proposition}\label{ra:c2:positive-basic}
设 $f,g$ 非负可测, $E,A,B$ 可测, $c\in[0,\infty)$.
\begin{enumerate}
\item 若 $f\le g$ 于 $E$ 上, 则 $\int_E f\,\dd\mu\le\int_Eg\,\dd\mu$.
\item 若 $A\subset B$, 则 $\int_Af\,\dd\mu\le\int_Bf\,\dd\mu$.
\item $\int_Ecf\,\dd\mu=c\int_Ef\,\dd\mu$.
\item 若 $f$ 在 $E$ 上恒为零, 或 $\mu(E)=0$, 则 $\int_E f\,\dd\mu=0$.
\item $\int_E f\,\dd\mu=\int_X\mathbf1_Ef\,\dd\mu$.
\end{enumerate}
\end{proposition}
\begin{proof}
先证第五项. 对 $0\le s\le f$ 的简单函数, $\mathbf1_Es$ 是不超过 $\mathbf1_Ef$ 的简单函数, 且其在 $X$ 上的积分等于 $s$ 在 $E$ 上的积分. 这给出左边不超过右边. 反过来, 若简单函数 $0\le t\le\mathbf1_Ef$, 则 $t=0$ 于 $E^c$ 上且 $t\le f$, 故 $\int_Xt\,\dd\mu=\int_Et\,\dd\mu$ 不超过左边. 取上确界得到等式.

若 $f\le g$ 于 $E$, 则 $\mathbf1_Ef\le\mathbf1_Eg$ 在整个 $X$ 上成立. 整空间上的单调性直接来自定义中候选简单函数族的包含, 再用第五项得第一项. 第二项同理, 因为 $\mathbf1_Af\le\mathbf1_Bf$. 对第三项, $c=0$ 时两边零; $c>0$ 时, 简单函数 $s\le f$ 与简单函数 $cs\le cf$ 一一对应, 上确界随乘 $c$ 而变为原上确界的 $c$ 倍. 对第四项, 若 $f$ 在 $E$ 上零, 每个候选简单函数在 $E$ 上也零; 若 $\mu(E)=0$, 每个水平集与 $E$ 的交测度均零. 两种情形中每个候选积分都是零, 上确界也零.
\end{proof}

\begin{example}
在 $\N$ 上取计数测度. 若 $f(k)=a_k\ge0$, 证明 $\int_\N f\,\dd\mu=\sum_{k=1}^{\infty}a_k$. 对 Dirac 测度, 证明 $\int_Xf\,\dd\delta_{x_0}=f(x_0)$.
\end{example}
\begin{solution}
计数测度下, 先由简单函数定义可知 $\int s\,\dd\mu=\sum_k s(k)$, 其中等式允许无穷. 因为 $s(k)\le a_k$, 每个候选积分不超过 $\sum_k a_k$. 若某个 $a_k=\infty$, 在该点取任意大的有限值, 便知积分无穷. 若所有 $a_k$ 有限, 简单函数 $s_N=\sum_{k=1}^{N}a_k\mathbf1_{\{k\}}$ 的积分为前 $N$ 项之和, 令 $N\to\infty$ 得反向不等式. 对 Dirac 测度, 每个简单函数的积分为 $s(x_0)$; 这些值不超过 $f(x_0)$. 若 $f(x_0)$ 有限, 取简单函数 $f(x_0)\mathbf1_{\{x_0\}}$; 若它无穷, 取 $M\mathbf1_{\{x_0\}}$ 并令 $M\to\infty$. 两种情形都达到所需上确界.
\end{solution}

\originalsection{单调收敛与 Fatou 引理}
非负积分的定义直接给出单调性, 却还没有证明一般函数的可加性. 若尝试把每个函数用简单函数逼近, 就必须说明逼近的极限能否移出积分. 因此先证明单调收敛定理, 再从它推出可加性; 这个顺序避免把尚未建立的线性作为证明前提.

\begin{theorem}[单调收敛定理]\label{ra:c2:mct}
设 $f_n:X\to[0,\infty]$ 可测, $f_n\uparrow f$ 逐点成立. 则 $f$ 可测且
\[
 \lim_{n\to\infty}\int_Xf_n\,\dd\mu=\int_Xf\,\dd\mu.
\]
等式允许两边为无穷.
\end{theorem}
\begin{proof}
可测性由定理~\ref{ra:c1:limits} 得到. 单调性表明 $\int f_n\,\dd\mu$ 递增, 令其极限为 $L\in[0,\infty]$. 因为 $f_n\le f$, 已有 $L\le\int f\,\dd\mu$.

要证明反向不等式, 固定可测简单函数 $0\le s\le f$ 和常数 $0<c<1$, 令 $E_n=\{f_n\ge cs\}$. 这些集合可测, 因为可将 $s$ 的有限水平集分别与 $f_n$ 的相应上水平集取交. 它们随 $n$ 递增, 而其并为 $X$: 当 $s(x)=0$ 时每个 $E_n$ 都含 $x$; 当 $s(x)>0$ 时, $cs(x)<s(x)\le f(x)$, 由 $f_n(x)\to f(x)$ 可知最终 $f_n(x)>cs(x)$. 因此
\[
 \int_X f_n\,\dd\mu\ge\int_{E_n}f_n\,\dd\mu
 \ge c\int_{E_n}s\,\dd\mu.
\]
命题~\ref{ra:c2:simple-properties} 说明 $\nu_s(E)=\int_Es\,\dd\mu$ 是测度. 用测度的向上连续性得 $\int_{E_n}s\,\dd\mu\uparrow\int_Xs\,\dd\mu$, 故 $L\ge c\int_Xs\,\dd\mu$. 若 $\int s\,\dd\mu=\infty$, 这已经迫使 $L=\infty$; 否则令 $c\uparrow1$ 得 $L\ge\int s\,\dd\mu$. 最后对全部这样的 $s$ 取上确界, 得 $L\ge\int f\,\dd\mu$.
\end{proof}
这里必须取 $c<1$. 若直接取 $c=1$, 一列从下方趋于 $s$ 却始终小于 $s$ 的函数会使 $E_n$ 不覆盖 $X$. 原讲义在开头写到 $c=1$ 的端点, 本书按其后续论证所需改为严格小于一.

\begin{corollary}\label{ra:c2:nonnegative-additive}
对非负可测函数 $f,g$ 有 $\int(f+g)\,\dd\mu=\int f\,\dd\mu+\int g\,\dd\mu$. 若 $f_n\ge0$ 可测, 则
\[
 \int_X\sum_{n=1}^{\infty}f_n\,\dd\mu
 =\sum_{n=1}^{\infty}\int_Xf_n\,\dd\mu.
\]
对固定非负可测 $f$, 集合函数 $\nu_f(E)=\int_Ef\,\dd\mu$ 是正测度.
\end{corollary}
\begin{proof}
用定理~\ref{ra:c1:simple-approx} 取 $s_n\uparrow f$ 和 $t_n\uparrow g$. 则 $s_n+t_n\uparrow f+g$, 简单函数可加性和单调收敛定理给出
\[
 \int(f+g)\,\dd\mu
 =\lim_n\left(\int s_n\,\dd\mu+\int t_n\,\dd\mu\right)
 =\int f\,\dd\mu+\int g\,\dd\mu.
\]
两列数均非负递增, 所以它们的和的极限等于极限之和, 即使某一极限无穷也成立. 对有限项相加反复应用此结论, 再对部分和 $\sum_{n=1}^{N}f_n$ 应用单调收敛, 得级数等式. 若 $E_n$ 两两不交, 逐点有 $f\mathbf1_{\cup_nE_n}=\sum_nf\mathbf1_{E_n}$. 代入级数等式和命题~\ref{ra:c2:positive-basic} 得 $\nu_f$ 可数可加, 而空集积分为零, 所以它是测度.
\end{proof}

\begin{corollary}\label{ra:c2:nonnegative-series}
若 $a_{ij}\ge0$, 则 $\sum_i\sum_j a_{ij}=\sum_j\sum_i a_{ij}$, 两边允许为无穷.
\end{corollary}
\begin{proof}
在 $\N$ 的计数测度上取 $f_j(i)=a_{ij}$. 上一推论把 $\int\sum_jf_j$ 写成 $\sum_j\int f_j$, 再使用计数积分与级数的对应关系即可.
\end{proof}

\begin{lemma}[Fatou]\label{ra:c2:fatou}
若 $f_n:X\to[0,\infty]$ 可测, 则
\[
 \int_X\liminf_{n\to\infty}f_n\,\dd\mu
 \le\liminf_{n\to\infty}\int_Xf_n\,\dd\mu.
\]
\end{lemma}
\begin{proof}
令 $g_k=\inf_{n\ge k}f_n$. 这些函数可测、非负且递增到 $\liminf_n f_n$. 对每个 $n\ge k$, $g_k\le f_n$, 故 $\int g_k\,\dd\mu\le\inf_{n\ge k}\int f_n\,\dd\mu$. 令 $k\to\infty$, 左边由单调收敛趋于 $\int\liminf_n f_n\,\dd\mu$, 右边按数列下极限定义趋于所述右端.
\end{proof}
单调收敛给出等式, 而没有单调性时通常只能保留 Fatou 的不等式. 函数值的质量可能向空间远处移动, 因此逐点极限可能丢失原来的积分.

\begin{example}
在 $\N$ 的计数测度下取 $f_n(k)=\mathbf1_{\{n\}}(k)$. 验证 Fatou 不等式可以严格, 并证明不存在一个可积函数 $g$ 同时控制所有 $f_n$.
\end{example}
\begin{solution}
对固定 $k$, 当 $n>k$ 时 $f_n(k)=0$, 所以 $f_n\to0$ 逐点成立. 每个 $f_n$ 的积分却为一, 因而 Fatou 两端分别为零和一. 若 $g\ge f_n$ 对所有 $n$ 成立, 令 $n=k$ 就得 $g(k)\ge1$ 对所有 $k$ 成立, 从而 $\int g\,\dd\mu\ge\sum_k1=\infty$. 这也说明仅有各积分一致有界, 并不等于有一个共同的可积控制函数.
\end{solution}

\originalsection{实值与复值积分}
\begin{definition}\label{ra:c2:l1-definition}
对可测扩展实值函数 $f$, 若 $\int f^+\,\dd\mu$ 与 $\int f^-\,\dd\mu$ 至少一个有限, 就定义 $\int f\,\dd\mu=\int f^+\,\dd\mu-\int f^-\,\dd\mu$. 若两者都有限, 称 $f$ 可积, 记为 $f\in\mathcal L^1(\mu)$. 等价条件是 $\int\abs f\,\dd\mu<\infty$.

对可测复值函数 $f=u+iv$, 当 $\int\abs f\,\dd\mu<\infty$ 时称其可积, 并定义 $\int f\,\dd\mu=\int u\,\dd\mu+i\int v\,\dd\mu$.
\end{definition}
实值积分在仅一个正负部积分有限时可以取 $\pm\infty$, 但不把两部都无穷的 $\infty-\infty$ 当作一个积分值. 后文涉及线性、复值积分和范数时, 均以可积函数为对象. 对复函数, $\abs u,\abs v\le\abs f\le\abs u+\abs v$, 所以可积当且仅当实部与虚部均可积; 这也保证其定义有意义.

当前 $\mathcal L^1$ 先指实际的可积函数. 将几乎处处相等的函数认同以后, 得到真正具有范数结构的空间 $L^1$; 本章后面的零测集讨论将说明这种认同为何不改变积分.

\begin{theorem}\label{ra:c2:linearity}
若 $f,g$ 为可积的有限实值或复值函数, $\alpha,\beta\in\C$, 则 $\alpha f+\beta g$ 可积, 且
\[
 \int_X(\alpha f+\beta g)\,\dd\mu
 =\alpha\int_Xf\,\dd\mu+\beta\int_Xg\,\dd\mu.
\]
\end{theorem}
\begin{proof}
先看可积性. 可测性由前章的连续运算得到, 而 $\abs{\alpha f+\beta g}\le\abs\alpha\abs f+\abs\beta\abs g$. 非负积分的单调性、可加性与齐次性使右边积分有限.

对实值 $f,g$, 令 $h=f+g$. 逐点恒等式 $h^++f^-+g^-=f^++g^++h^-$ 只涉及非负函数. 用已经建立的非负可加性积分, 再移项, 得
\[
 \int h^+\,\dd\mu-\int h^-\,\dd\mu
 =\left(\int f^+\,\dd\mu-\int f^-\,\dd\mu\right)
 +\left(\int g^+\,\dd\mu-\int g^-\,\dd\mu\right).
\]
这里所有积分有限, 所以移项合法. 对实数 $a>0$, 有 $(af)^+=af^+$ 和 $(af)^-=af^-$, 非负齐次性给 $\int af=a\int f$. 对 $a<0$, 正负部交换为 $(af)^+=(-a)f^-$ 和 $(af)^-=(-a)f^+$, 同样得到等式; $a=0$ 时两边为零.

对复函数, 分实部与虚部使用实值可加性. 若 $\alpha=a+ib$ 与 $f=u+iv$, 则 $\alpha f=(au-bv)+i(av+bu)$, 用实值线性得到 $\int\alpha f=(a+ib)(\int u+i\int v)=\alpha\int f$. 和的线性再给出一般结论.
\end{proof}
可积扩展实值函数可能在一个零测集上取无穷. 如需进行相加等逐点代数运算, 将这些无穷值改为零, 即可选取处处有限的代表. 这种修改合法的理由在下一节证明; 线性定理的逐点代数运算针对有限代表进行.

\begin{theorem}[积分的三角不等式]\label{ra:c2:triangle}
对可积复值函数 $f$, 有 $\abs{\int_Xf\,\dd\mu}\le\int_X\abs f\,\dd\mu$.
\end{theorem}
\begin{proof}
若积分为零, 不等式立即成立. 否则写 $\int f\,\dd\mu=e^{i\theta}\abs{\int f\,\dd\mu}$. 乘以 $e^{-i\theta}$ 后实部等于这个绝对值. 由复积分定义、线性和实积分的单调性, 得
\[
 \abs{\int f\,\dd\mu}
 =\int\operatorname{Re}(e^{-i\theta}f)\,\dd\mu
 \le\int\abs{e^{-i\theta}f}\,\dd\mu
 =\int\abs f\,\dd\mu.
\]
其中实积分单调性来自可积差函数 $\abs f-\operatorname{Re}(e^{-i\theta}f)\ge0$ 的积分非负.
\end{proof}

\originalsection{零测集与几乎处处}
\begin{proposition}\label{ra:c2:zero-integral}
若非负可测 $f$ 满足 $\int_E f\,\dd\mu=0$, 则 $f=0$ 几乎处处于 $E$. 反过来, 若 $f=0$ 几乎处处于 $E$, 则此积分为零. 非负可测 $f$ 若积分有限, 则 $f<\infty$ 几乎处处.
\end{proposition}
\begin{proof}
令 $A_n=E\cap\{f>1/n\}$. 由单调性有 $0=\int_Ef\,\dd\mu\ge n^{-1}\mu(A_n)$, 所以每个 $A_n$ 零测. 因为 $E\cap\{f>0\}=\bigcup_n A_n$, 故其测度零. 反向, 选一个可测零测集 $N$ 使 $f=0$ 于 $E\setminus N$. 区域可加性将积分分为 $E\setminus N$ 与 $E\cap N$ 上的两个积分, 前者因函数恒零而为零, 后者因区域零测而为零.

若 $C=\{f=\infty\}$, 则对任意正整数 $n$, $n\mathbf1_C\le f$, 故 $n\mu(C)\le\int f\,\dd\mu<\infty$. 令 $n\to\infty$ 得 $\mu(C)=0$.
\end{proof}

\begin{proposition}\label{ra:c2:ae-invariance}
若可测函数 $f,g$ 几乎处处相等, 则非负情形下对每个 $E\in\mathcal M$ 有 $\int_Ef\,\dd\mu=\int_Eg\,\dd\mu$. 在实值或复值情形, $f$ 可积当且仅当 $g$ 可积, 且这时区域积分也相等.
\end{proposition}
\begin{proof}
选可测零测集 $N$ 使 $f=g$ 在 $X\setminus N$ 上成立. 对非负函数, 将 $E$ 分成 $E\setminus N$ 与 $E\cap N$. 第一部分的函数相同, 第二部分两者积分均零, 故得到等式. 对实值或复值函数, 先对 $\abs f,\abs g$ 应用非负结论, 得可积性等价; 再分别对正负部, 或实部与虚部的正负部应用同一结论, 得积分相等.
\end{proof}
因而几乎处处相等是可积函数上的等价关系: 自反与对称直接成立; 若 $f=g$ 与 $g=h$ 分别在零测集外成立, 则在两零测集的并之外 $f=h$, 给出传递性. 记 $L^1(\mu)=\mathcal L^1(\mu)/{\sim}$, 其中 $f\sim g$ 表示几乎处处相等. 此时 $\|f\|_1=\int\abs f\,\dd\mu$ 在等价类上与代表选择无关, 且等于零当且仅当该类为零类. 更一般的 $L^p$ 空间在后文建立.

\begin{proposition}\label{ra:c2:local-zero}
若可积复值函数 $f$ 满足 $\int_Ef\,\dd\mu=0$ 对每个可测 $E$ 成立, 则 $f=0$ 几乎处处.
\end{proposition}
\begin{proof}
写 $f=u+iv$. 取 $E=\{u>0\}$, 假设积分的实部为零, 即 $\int_Eu\,\dd\mu=\int_Xu^+\,\dd\mu=0$. 命题~\ref{ra:c2:zero-integral} 得 $u^+=0$ 几乎处处. 再取 $E=\{u<0\}$, 得 $\int_Xu^-\,\dd\mu=0$, 故 $u^-=0$ 几乎处处. 对 $v$ 分别取 $\{v>0\}$ 和 $\{v<0\}$, 得 $v^+=v^-=0$ 几乎处处. 四个例外零测集的并仍零测, 因而 $f=0$ 几乎处处.
\end{proof}

\begin{theorem}[第一 Borel--Cantelli 引理]\label{ra:c2:bc}
若可测集列 $E_k$ 满足 $\sum_{k=1}^{\infty}\mu(E_k)<\infty$, 则几乎每个点只属于有限多个 $E_k$, 即 $\mu(\limsup_k E_k)=0$.
\end{theorem}
\begin{proof}
令 $g=\sum_k\mathbf1_{E_k}$. 非负积分的级数交换给出 $\int g\,\dd\mu=\sum_k\mu(E_k)<\infty$, 所以 $g<\infty$ 几乎处处. 一个点属于无穷多个 $E_k$ 当且仅当 $g$ 在该点为无穷, 故这些点组成零测集.
\end{proof}
这里没有使用概率测度的归一化, 也没有使用事件的独立性. 编者补充的另一个直接证明是: 对任意 $N$, $\limsup_k E_k\subset\bigcup_{k\ge N}E_k$, 因而其测度不超过尾和 $\sum_{k\ge N}\mu(E_k)\to0$.

\originalsection{主导收敛定理}
\begin{theorem}[主导收敛定理]\label{ra:c2:dct}
设 $f_n:X\to\C$ 可测, 在某个可测零测集以外收敛到 $f$. 假定存在非负可积函数 $g$, 使 $\abs{f_n}\le g$ 对每个 $n$ 几乎处处成立. 在例外零测集上把 $f$ 定义为零, 则 $f$ 可测、可积, 且
\[
 \int_X\abs{f_n-f}\,\dd\mu\longrightarrow0,
 \qquad \int_Xf_n\,\dd\mu\longrightarrow\int_Xf\,\dd\mu.
\]
在第一式中的差取处处有限的代表计算.
\end{theorem}
\begin{proof}
把不收敛的零测集、各个控制不等式失效的零测集, 以及 $\{g=\infty\}$ 合并成一个可测零测集 $N$. 将 $f_n$ 和 $g$ 在 $N$ 上同时改成零. 这种在一个可测集上的分段定义保持可测性, 且由命题~\ref{ra:c2:ae-invariance} 不改变积分. 于是可以假定函数处处有限, $f_n\to f$ 处处成立, 且 $\abs{f_n}\le g$ 处处成立. 可测函数的逐点极限可测, 并有 $\abs f\le g$, 所以 $f$ 及各 $f_n$ 可积.

令 $h_n=\abs{f_n-f}$. 则 $0\le h_n\le2g$ 且 $h_n\to0$. 对非负函数 $2g-h_n$ 应用 Fatou 引理, 得
\begin{align*}
 \int_X2g\,\dd\mu
 &\le\liminf_n\int_X(2g-h_n)\,\dd\mu\\
 &=\int_X2g\,\dd\mu-\limsup_n\int_Xh_n\,\dd\mu.
\end{align*}
这里 $2g$ 可积, 故等式使用有限量的线性, 不涉及 $\infty-\infty$. 消去有限的 $\int2g\,\dd\mu$, 得 $\limsup_n\int h_n\,\dd\mu\le0$. 因各积分非负, 它们必趋于零. 再由积分三角不等式,
\[
 \abs{\int_Xf_n\,\dd\mu-\int_Xf\,\dd\mu}
 \le\int_X\abs{f_n-f}\,\dd\mu\longrightarrow0.
\]
恢复原来的代表不改变上述积分, 所以结论成立.
\end{proof}
原讲义先写处处版本; 上面的几乎处处版本通过一个共同零测集归约到该版本. 这个归约中用到了函数列只有可数多个成员. 控制函数必须同时控制整列, 并且自身可积; 各 $f_n$ 各自可积或其积分一致有界, 都不足以替代它.

\begin{corollary}[编者补充]\label{ra:c2:bounded-convergence}
若 $\mu(X)<\infty$, $f_n\to f$ 几乎处处, 且存在有限常数 $M$ 使 $\abs{f_n}\le M$ 对所有 $n$ 几乎处处成立, 则 $\int\abs{f_n-f}\,\dd\mu\to0$.
\end{corollary}
\begin{proof}
取控制函数 $g=M\mathbf1_X$. 其积分为 $M\mu(X)<\infty$, 直接应用主导收敛定理.
\end{proof}

\begin{example}
编者补充. 在 $\N$ 上定义有限测度 $\mu(\{k\})=2^{-k}$, 并按可数可加性令 $\mu(E)=\sum_{k\in E}2^{-k}$. 取 $f_n(k)=2^n\mathbf1_{\{n\}}(k)$. 证明 $f_n\to0$ 逐点成立而其积分不趋于零. 再说明 $h_n(k)=\mathbf1_{\{k\ge n\}}(k)$ 的积分为何趋于零.
\end{example}
\begin{solution}
固定 $k$ 后, $f_n(k)$ 只有在 $n=k$ 时非零, 所以逐点趋于零; 但 $\int f_n\,\dd\mu=2^n2^{-n}=1$. 若有共同可积控制函数 $g$, 逐点令 $n=k$ 得 $g(k)\ge2^k$, 因而 $\int g\,\dd\mu\ge\sum_k2^k2^{-k}=\infty$, 矛盾. 对 $h_n$, 则 $0\le h_n\le1$, 且 $\mu(\N)=1$, 故有界收敛推论给出积分趋于零. 直接计算也得到 $\int h_n\,\dd\mu=\sum_{k=n}^{\infty}2^{-k}=2^{1-n}\to0$.
\end{solution}

\originalsection{测度的完备化}
零测集中的任意子集在积分上应当没有贡献, 但在原来的 $\sigma$-代数中可能还没有它. 完备化把这些缺少的子集加入, 同时保持已有测度.

\begin{definition}
测度空间称为完备, 如果每个可测零测集的每个子集都可测. 给定 $(X,\mathcal M,\mu)$, 定义
\[
 \mathcal M^*=\{E\subset X:\text{ 存在 }A,B\in\mathcal M,
       \ A\subset E\subset B,\ \mu(B\setminus A)=0\}.
\]
对上述 $E$, 拟定义 $\mu^*(E)=\mu(A)$.
\end{definition}

\begin{theorem}\label{ra:c2:completion}
$\mathcal M^*$ 是 $\sigma$-代数, $\mu^*$ 与 $A,B$ 的选择无关, 且是扩张 $\mu$ 的正测度. 测度空间 $(X,\mathcal M^*,\mu^*)$ 完备, 称为原空间的完备化.
\end{theorem}
\begin{proof}[证明(编者补充)]
先证定义无歧义. 若 $A\subset E\subset B$ 与 $A'\subset E\subset B'$ 是两组包夹, 则 $A\setminus A'\subset B'\setminus A'$ 为零测, 同样 $A'\setminus A\subset B\setminus A$ 为零测. 这些差集本来就在 $\mathcal M$ 中, 所以由单调性其测度为零. 分别把 $A$ 和 $A'$ 分为公共部分及其零测差集, 得 $\mu(A)=\mu(A\cap A')=\mu(A')$. 这不需要从无穷测度中作减法.

取 $A=B=X$ 得 $X\in\mathcal M^*$. 若 $A\subset E\subset B$, 则 $B^c\subset E^c\subset A^c$, 而 $A^c\setminus B^c=B\setminus A$ 零测, 所以补集仍属 $\mathcal M^*$. 若 $A_n\subset E_n\subset B_n$ 是可数个包夹, 则 $\bigcup_nA_n\subset\bigcup_nE_n\subset\bigcup_nB_n$, 且
\[
 \left(\bigcup_n B_n\right)\setminus\left(\bigcup_n A_n\right)
 \subset\bigcup_n(B_n\setminus A_n).
\]
左边是原 $\sigma$-代数中的集合, 右边为可测零测集, 由单调性左边零测. 故 $\mathcal M^*$ 对可数并封闭.

显然 $\mu^*(\varnothing)=0$, 且对 $E\in\mathcal M$ 取 $A=B=E$ 得 $\mu^*(E)=\mu(E)$. 若 $E_n\in\mathcal M^*$ 两两不交, 其包夹下集 $A_n$ 也两两不交. 上一段给出了它们并集的有效包夹, 因而
\[
 \mu^*\left(\bigcup_nE_n\right)
 =\mu\left(\bigcup_nA_n\right)
 =\sum_n\mu(A_n)=\sum_n\mu^*(E_n).
\]
于是 $\mu^*$ 为正测度. 最后, 若 $N\in\mathcal M^*$ 且 $\mu^*(N)=0$, 选 $A\subset N\subset B$ 的包夹. 有 $\mu(A)=0$ 与 $\mu(B\setminus A)=0$, 所以 $\mu(B)=0$. 对任意 $S\subset N$, 包夹 $\varnothing\subset S\subset B$ 表明 $S\in\mathcal M^*$ 且 $\mu^*(S)=0$. 故扩张空间完备.
\end{proof}

\begin{proposition}[编者补充]\label{ra:c2:completion-function}
对原 $\sigma$-代数可测的非负或可积函数, 完备化不改变积分. 在完备空间中, 在零测集的任意子集上任意修改可测函数, 所得函数仍可测, 且只要原函数可积, 积分不变.
\end{proposition}
\begin{proof}
原可测简单函数的积分不变, 因为各水平集的测度不变. 对非负原可测 $f$, 取原可测简单函数 $s_n\uparrow f$, 分别在原空间与完备空间用单调收敛定理, 两边极限相同. 对可积函数再分正负部、实虚部, 得积分也相同.

设 $h=f$ 在一个可测零测集 $N$ 之外成立, $h$ 在 $N$ 上任意取值. 对任意开集 $U$, 两个原像 $h^{-1}(U)$ 与 $f^{-1}(U)$ 的差别包含于 $N$. 完备性使其中的任意子集可测, 故 $h^{-1}(U)$ 可测. 命题~\ref{ra:c2:ae-invariance} 给出积分不变.
\end{proof}
在未完备空间, ``在零测集上任意修改仍可测''不一定成立: 若零测集有不可测子集 $S$, 把零函数在 $S$ 上改成一, 就得到不可测的 $\mathbf1_S$. 原讲义关于只在几乎处处定义的函数可任意补值的说明, 必须配合完备性, 或要求所作补值本身可测. 统一在一个已知可测例外集上补成零, 则不需要原空间完备.

\originalsection{与 Riemann 积分的衔接}
本节取闭区间 $[a,b]$, $a<b$. 调用下一章构造的完备 Lebesgue 测度 $m$, 满足 $m((u,v))=v-u$, 单点零测, 且零测集能用总长度任意小的可数个开区间覆盖. 对 $[a,b]$ 上的有界实函数 $f$ 与分割 $P:a=x_0<\cdots<x_r=b$, 记各小区间上的下、上确界为 $m_j,M_j$, 定义 Darboux 和 $L(f,P)=\sum_jm_j(x_j-x_{j-1})$ 与 $U(f,P)=\sum_jM_j(x_j-x_{j-1})$. 下积分为所有下和的上确界, 上积分为所有上和的下确界; 二者相等时称 $f$ Riemann 可积.

这里的可积判据必须包含\emph{有界}这一假设. 通常的闭区间 Riemann 积分要求函数有界; 反常积分另作定义, 不能直接代入下面的判据. 原讲义陈述这一判据时没有重复写出有界条件, 本书明确补上.

\begin{lemma}[编者补充]\label{ra:c2:riemann-darboux}
有界函数 $f$ Riemann 可积当且仅当对每个 $\eps>0$, 存在分割 $P$ 使 $U(f,P)-L(f,P)<\eps$.
\end{lemma}
\begin{proof}
下和不超过任意上和: 对两个分割取共同加细, 下和随加细不减, 上和随加细不增, 且同一分割的下和不超过上和. 因而下积分不超过上积分. 若存在任意小的上下和差, 上下积分之差也不超过这个差, 故两积分相等. 反过来, 若它们同为 $I$, 取一个下和大于 $I-\eps/2$ 的分割及一个上和小于 $I+\eps/2$ 的分割. 它们的共同加细使上下和差小于 $\eps$.
\end{proof}

为描述不连续点, 在相对区间 $[a,b]$ 内定义点的振幅
\[
 \omega_f(x)=\inf_{r>0}\left(
   \sup_{\substack{y\in[a,b]\\\abs{y-x}<r}}f(y)
  -\inf_{\substack{y\in[a,b]\\\abs{y-x}<r}}f(y)\right).
\]
振幅零当且仅当 $f$ 在 $x$ 连续: 连续性给出小邻域中任意两函数值之差小; 反之, 小邻域的上下确界差小, 而 $f(x)$ 本身在其间, 所以附近值接近 $f(x)$. 对 $\eta>0$, 集合 $D_\eta=\{x:\omega_f(x)\ge\eta\}$ 闭. 确实, 若 $\omega_f(x)<\eta$, 可找到一个邻域使其中振幅小于 $\eta$; 对足够邻近的 $z$, 更小的 $z$ 邻域包含于该邻域, 于是 $\omega_f(z)<\eta$. 因而不连续点集为 Borel 集 $D=\bigcup_{n=1}^{\infty}D_{1/n}$.

\begin{theorem}[Lebesgue 的 Riemann 可积判据]\label{ra:c2:riemann-criterion}
有界实函数 $f:[a,b]\to\R$ Riemann 可积当且仅当它的不连续点集 $D$ 为 Lebesgue 零测集. 在这种情形, $f$ 为 Lebesgue 可测且可积, 两种积分相等.
\end{theorem}
\begin{proof}[证明(编者补充)]
先假设 $f$ Riemann 可积. 固定 $\eta>0$, 给定任意 $\delta>0$, 选分割 $P$ 使 $U(f,P)-L(f,P)<\eta\delta$. 将振幅 $M_j-m_j\ge\eta$ 的小区间称为坏区间. 若 $x\in D_\eta$ 且 $x$ 不是分割点, 则它在某个小区间的内部; 若该区间振幅小于 $\eta$, $x$ 的更小邻域也会如此, 矛盾. 所以 $D_\eta$ 包含于坏区间的并与有限个分割点的并, 从而
\[
 m(D_\eta)\le\sum_{\text{坏区间 }j}(x_j-x_{j-1})
 \le\frac{U(f,P)-L(f,P)}{\eta}<\delta.
\]
由于 $\delta$ 任意, $m(D_\eta)=0$. 可数并表示给出 $m(D)=0$.

反过来, 假设 $D$ 零测, 并取 $M\ge0$ 使 $\abs f\le M$. 固定 $\eta,\delta>0$. 闭集 $D_\eta\subset[a,b]$ 紧且零测, 所以可以用有限个开区间覆盖它, 这些区间总长度小于 $\delta$. 令其并为 $V$. 对每个 $x\notin D_\eta$, 由振幅定义可选相对开邻域 $U_x$, 使 $f$ 在整个 $U_x$ 上的振幅小于 $\eta$. $V$ 与这些 $U_x$ 组成 $[a,b]$ 的开覆盖, 紧性允许取有限子覆盖. 有限开覆盖存在 Lebesgue 数: 存在 $\rho>0$ 使直径小于 $\rho$ 的区间必含于某个覆盖成员. 这一基础紧性结论可如下证明: 否则取直径趋零且不含于任何成员的集合, 选其一点, 再取收敛子列; 极限点位于一个开成员内, 小直径使该子列对应集合最终都包含于这个成员, 矛盾.

取网格小于 $\rho$ 的分割. 每个小区间或者含于 $V$, 或者含于某个 $U_x$. 把前一种归入坏类, 其长度总和不超过 $m(V)<\delta$, 振幅不超过 $2M$; 其余小区间振幅小于 $\eta$. 因而 $U(f,P)-L(f,P)\le2M\delta+\eta(b-a)$. 先选 $\eta$ 再选 $\delta$, 可使右边任意小. 若 $M=0$ 结论直接成立. Darboux 判据给出 Riemann 可积性.

最后证明与 Lebesgue 积分相等, 并同时核对可测性. 对每个 $n$, 取分割 $P_n$ 使上下和差小于 $2^{-n}$. 在各小区间上分别取常数下、上确界得到 Borel 简单函数 $\ell_n,u_n$, 小区间用半开方式分割, 最后包含端点 $b$. 有 $\ell_n\le f\le u_n$ 逐点成立, 且 $\int(u_n-\ell_n)\,dm<2^{-n}$. 令 $\ell=\sup_n\ell_n$, $u=\inf_nu_n$, 则二者 Borel 可测、有限, $\ell\le f\le u$, 并且 $0\le u-\ell\le u_n-\ell_n$ 对每个 $n$ 成立. 故 $\int(u-\ell)\,dm=0$, 从而 $u=\ell$ 几乎处处. 完备性与命题~\ref{ra:c2:completion-function} 保证夹在二者之间的 $f$ Lebesgue 可测. 它有界且区间测度有限, 所以可积. 积分单调性给出 $L(f,P_n)\le\int f\,dm\le U(f,P_n)$, 另一方面 Riemann 积分也夹在同两和之间. 两和差趋零, 因而两积分相等.
\end{proof}

\begin{example}
编者补充. 证明 $\mathbf1_{\mathbb Q\cap[a,b]}$ Lebesgue 可积但不 Riemann 可积. 解释这与上述判据的一致性.
\end{example}
\begin{solution}
可数集合的 Lebesgue 测度为零, 因为每个单点零测并且测度可数次可加. 故此示性函数积分为零. 但每个非退化小区间都有有理数和无理数, 上确界为一、下确界为零, 所以上下 Darboux 和总是 $b-a$ 与零, 不可能相等. 该函数处处不连续, 不连续点集就是正测度区间 $[a,b]$, 正与判据一致.
\end{solution}

\originalsection{练习与解答}
以下题目及解答为编者补充.

\begin{exercise}
设 $0\le f_n\downarrow f$, 函数均可测且 $\int f_1\,\dd\mu<\infty$. 证明 $\int f_n\,\dd\mu\downarrow\int f\,\dd\mu$. 给出删去有限积分条件后结论失败的例子.
\end{exercise}
\begin{solution}
由 $f_1$ 可积可选共同零测集外有限的代表. 非负函数 $f_1-f_n$ 递增到 $f_1-f$, 用单调收敛及有限量的积分线性得到 $\int f_1-\int f_n\to\int f_1-\int f$. 故 $\int f_n\to\int f$, 而其单调递减性来自积分单调性. 在 $\N$ 的计数测度上取 $f_n=\mathbf1_{\{n,n+1,\ldots\}}$, 则逐点递减到零, 但每个积分均无穷, 结论失败.
\end{solution}

\begin{exercise}
若 $f_n$ 可测且 $\sum_n\int\abs{f_n}\,\dd\mu<\infty$, 证明 $\sum_nf_n$ 几乎处处绝对收敛, 其和可积, 且可以逐项积分.
\end{exercise}
\begin{solution}
非负级数交换给出 $\int\sum_n\abs{f_n}\,\dd\mu=\sum_n\int\abs{f_n}\,\dd\mu<\infty$. 故 $G=\sum_n\abs{f_n}$ 几乎处处有限. 在这个集合上级数绝对收敛, 在其零测补集上把和定义为零. 部分和可测, 所得和 $F$ 可测, 并满足 $\abs F\le G$ 几乎处处, 因而可积. 部分和也都由同一个 $G$ 控制, 主导收敛给 $\int F\,\dd\mu=\lim_N\int\sum_{n=1}^{N}f_n\,\dd\mu=\sum_n\int f_n\,\dd\mu$. 数值级数也绝对收敛, 因为 $\sum_n\abs{\int f_n\,\dd\mu}\le\sum_n\int\abs{f_n}\,\dd\mu$.
\end{solution}

\begin{exercise}
设 $f_n,f$ 可测, 且 $\sum_n\int\abs{f_n-f}\,\dd\mu<\infty$. 证明 $f_n\to f$ 几乎处处.
\end{exercise}
\begin{solution}
对正整数 $j$, 令 $E_{n,j}=\{\abs{f_n-f}>1/j\}$. 单调性给 $\mu(E_{n,j})\le j\int\abs{f_n-f}\,\dd\mu$, 所以 $\sum_n\mu(E_{n,j})<\infty$. 第一 Borel--Cantelli 引理说明对几乎每个 $x$, 固定 $j$ 后只有有限多个 $n$ 使 $x\in E_{n,j}$. 对所有 $j$ 的例外集取可数并, 在其补集上, 对每个 $j$ 都最终有 $\abs{f_n(x)-f(x)}\le1/j$, 这正是收敛到零的条件.
\end{solution}

\begin{exercise}
设 $T$ 为度量空间, $t_0\in T$. 对每个 $t\in T$ 有可测函数 $F(t,\cdot):X\to\C$. 假设对几乎每个 $x$, $F(t,x)$ 在 $t=t_0$ 连续, 并且在 $t_0$ 的一个邻域内有共同的非负可积控制函数 $g$, 即 $\abs{F(t,x)}\le g(x)$ 几乎处处, 其中允许每个 $t$ 的例外集不同. 在这个邻域内定义 $I(t)=\int F(t,x)\,\dd\mu(x)$, 证明 $I$ 在 $t_0$ 连续.
\end{exercise}
\begin{solution}
取任意趋于 $t_0$ 的序列 $t_n$. 删去有限个不在指定邻域中的成员不影响极限. 连续性在一个固定零测集以外给出 $F(t_n,x)\to F(t_0,x)$, 而控制条件对这列可数多个参数的例外集可以取并, 得到共同零测集. 对该序列应用主导收敛, 有 $I(t_n)\to I(t_0)$. 度量空间中函数在一点连续当且仅当保持所有趋于该点的序列极限: 若不连续, 可在半径 $1/n$ 的球内选取函数值与 $I(t_0)$ 相距至少某个固定正数的点, 与序列结论矛盾. 因而 $I$ 在 $t_0$ 连续.
\end{solution}

\begin{exercise}
在 Lebesgue 区间 $[0,1]$ 上取 $f_n(x)=x^n$. 计算其几乎处处极限, 并分别用主导收敛与 Riemann 积分计算检验积分极限.
\end{exercise}
\begin{solution}
当 $0\le x<1$ 时 $x^n\to0$, 而 $x=1$ 时恒为一. 所以逐点极限是 $\mathbf1_{\{1\}}$, 几乎处处极限为零. 共同控制函数一在区间上可积, 主导收敛给 $\int_0^1x^n\,dm\to0$. 每个 $x^n$ 连续, Lebesgue 与 Riemann 积分相同, 微积分计算为 $\int_0^1x^n\,dx=1/(n+1)\to0$. 单点处的不同极限值不改变积分, 而控制函数的可积性使积分极限交换有了依据.
\end{solution}

\originalchapter{Lebesgue 测度}
\label{ra:c3:chapter}
\sourcelecture{6--9}

\originalsection{矩形与有限并}

一般测度的积分理论已经建立. 现在要在 $\R^n$ 上构造一个测度, 使轴平行矩形的测度等于通常的体积. 原讲义依次处理矩形、矩形的有限并、开集、紧集, 最后用内外逼近辨认可测集. 这条路线把几何体积与抽象测度联系起来.

\begin{definition}\label{ra:c3:polygon}
轴平行闭矩形是 $I=\prod_{j=1}^n[a_j,b_j]$, 其中 $a_j\le b_j$; 其体积定义为 $v(I)=\prod_{j=1}^n(b_j-a_j)$. 有限个这类矩形的并称为矩形多面集. 将各坐标方向出现的全部端点作为分割点, 所得小矩形的内部互不相交; $v(P)$ 定义为被 $P$ 包含的小矩形的体积之和. 退化矩形的体积取 $0$.
\end{definition}

\begin{lemma}\label{ra:c3:grid}
上述定义与分割方式无关. 对矩形多面集, 体积单调、有限次可加; 若 $P_1,P_2$ 的交只在各自的矩形边界上, 则 $v(P_1\cup P_2)=v(P_1)+v(P_2)$.
\end{lemma}
\begin{proof}
将区间 $[a,b]$ 分为 $[a,t_1],\ldots,[t_m,b]$, 各段长度之和是 $b-a$. 在 $n$ 个方向分别分割, 将这些等式相乘, 得原矩形体积等于所有小矩形体积之和. 两个有限分割都有共同细分, 所以体积定义不变. 在共同网格上, 包含关系就是选择网格单元的包含关系, 从而得到单调性. 对内部互不相交的并, 每个非退化网格单元恰计一次, 故体积相加. 边界由有限个退化矩形组成, 不贡献体积.
\end{proof}

\begin{remark}
\textbf{编者修正.} 原第 6 讲将闭矩形的分解直接写成不交并. 相邻闭矩形会共享边界. 上面的共同网格约定先在有限体积层面消除这一歧义; 在测度建立后, 改用半开矩形也会给出同一数值.
\end{remark}

\originalsection{开集与紧集}

\begin{definition}\label{ra:c3:opencompact}
对开集 $G\subset\R^n$, 定义 $\lambda(G)=\sup\{v(P):P\subset G,\ P\text{ 为矩形多面集}\}$, 并取 $\lambda(\varnothing)=0$. 对紧集 $K$, 定义 $\lambda(K)=\inf\{\lambda(G):K\subset G,\ G\text{ 开}\}$.
\end{definition}

\begin{proposition}\label{ra:c3:openproperties}
非空开集的测度为正, $\lambda(\R^n)=\infty$, 且开集测度单调. 对任意开集列 $G_k$, 有 $\lambda(\bigcup_kG_k)\le\sum_k\lambda(G_k)$; 若这些开集两两不交, 则等号成立.
\end{proposition}
\begin{proof}
非空开集包含一个各边长度为正的小矩形; $\R^n$ 包含任意大的矩形. 单调性直接来自取上确界的集合包含关系.

设 $P\subset\bigcup_kG_k$ 为矩形多面集. $P$ 紧, 故有有限个 $G_k$ 覆盖它. 对这个有限开覆盖, 存在 $\delta>0$, 使 $P$ 中直径小于 $\delta$ 的每个子集包含在某一覆盖成员内. 证明这一点时, 对每个 $x\in P$ 选球 $B(x,2r_x)$ 包含在某一成员内, 再从 $B(x,r_x)$ 中取有限子覆盖, 令 $\delta$ 小于所取各半径的最小值即可. 将 $P$ 的共同网格进一步细分, 使各小矩形直径小于 $\delta$. 将每个小矩形分配给一个包含它的 $G_k$, 合并得到 $P_k\subset G_k$. 它们仅在边界上可能相交, 因而
\[
 v(P)=\sum_kv(P_k)\le\sum_k\lambda(G_k).
\]
对 $P$ 取上确界, 得到次可加性. 这补全了原稿中有限分割的论证.

若 $G_k$ 两两不交, 固定 $N$ 并选 $P_k\subset G_k$ ($1\le k\le N$). 其并包含在 $\bigcup_kG_k$ 中, 故 $\sum_{k=1}^Nv(P_k)\le\lambda(\bigcup_kG_k)$. 分别取上确界, 再令 $N\to\infty$, 得反向不等式.
\end{proof}

\begin{lemma}\label{ra:c3:consistency}
若 $P$ 为矩形多面集, 则 $\lambda(P^\circ)=v(P)$, 而按紧集定义得到的 $\lambda(P)$ 也等于 $v(P)$.
\end{lemma}
\begin{proof}
任何 $Q\subset P^\circ$ 都有 $v(Q)\le v(P)$. 反过来, 在共同网格上将每个非退化小矩形略向内缩, 得到 $P'\subset P^\circ$, 且 $v(P')>v(P)-\eps$. 于是 $\lambda(P^\circ)=v(P)$.

若开集 $G$ 包含 $P$, 则 $\lambda(G)\ge v(P)$, 所以紧集定义的数值不小于 $v(P)$. 将有限个小矩形分别略向外扩大, 使扩大的总体积小于 $v(P)+\eps$, 并使退化部分也包含在总体积小于 $\eps$ 的小矩形中. 扩大矩形内部的并是包含 $P$ 的开集; 由开集次可加性, 其测度至多为扩大的总体积. 令扩大造成的总误差趋于 $0$, 得反向不等式.
\end{proof}

\begin{proposition}\label{ra:c3:compactproperties}
紧集测度有限且单调, 满足 $\lambda(K_1\cup K_2)\le\lambda(K_1)+\lambda(K_2)$. 若 $K_1\cap K_2=\varnothing$, 则等号成立.
\end{proposition}
\begin{proof}
紧集有界, 所以可用一个有限体积的开矩形包住它. 单调性来自开邻域族的包含关系. 分别选开邻域 $G_i\supset K_i$, 利用 $\lambda(G_1\cup G_2)\le\lambda(G_1)+\lambda(G_2)$, 再取下确界, 得次可加性.

两个不交紧集间有正距离, 故有不交开邻域 $U_1,U_2$. 对任意 $G\supset K_1\cup K_2$, 令 $G_i=G\cap U_i$. 则 $\lambda(K_1)+\lambda(K_2)\le\lambda(G_1)+\lambda(G_2)=\lambda(G_1\cup G_2)\le\lambda(G)$. 对 $G$ 取下确界即可.
\end{proof}

\begin{lemma}\label{ra:c3:compactgap}
若 $K\subset G$, $K$ 紧且 $G$ 开, 则 $\lambda(G)=\lambda(K)+\lambda(G\setminus K)$.
\end{lemma}
\begin{proof}
对任意矩形多面集 $P\subset G\setminus K$, 紧集可加性给出 $\lambda(K)+v(P)=\lambda(K\cup P)\le\lambda(G)$. 对 $P$ 取上确界, 得一个方向.

另选开集 $U\supset K$, 使 $\lambda(U)<\lambda(K)+\eps$. 对任意 $P\subset G$, 有 $P\subset U\cup(G\setminus K)$, 故 $v(P)\le\lambda(U)+\lambda(G\setminus K)$. 对 $P$ 取上确界, 再令 $\eps\downarrow0$. 注意 $\lambda(K)<\infty$, 全程不使用 $\infty-\infty$.
\end{proof}

\originalsection{内外测度与有限可测集}

\begin{definition}\label{ra:c3:innerouter}
对任意 $A\subset\R^n$, 定义外测度 $\lambda^*(A)=\inf\{\lambda(G):A\subset G,\ G\text{ 开}\}$, 内测度 $\lambda_*(A)=\sup\{\lambda(K):K\subset A,\ K\text{ 紧}\}$. 若 $\lambda^*(A)<\infty$ 且 $\lambda_*(A)=\lambda^*(A)$, 称 $A$ 为有限 Lebesgue 可测集; 记全体此类集合为 $\mathcal L_0$, 并令 $\lambda(A)$ 等于上述共同值.
\end{definition}

\begin{proposition}\label{ra:c3:outerproperties}
内外测度均单调, 且 $\lambda_*(A)\le\lambda^*(A)$. 外测度满足可数次可加不等式 $\lambda^*(\bigcup_kA_k)\le\sum_k\lambda^*(A_k)$. 若 $A_k$ 两两不交, 则 $\lambda_*(\bigcup_kA_k)\ge\sum_k\lambda_*(A_k)$. 对开集或紧集 $A$, 两种测度都等于先前定义的 $\lambda(A)$.
\end{proposition}
\begin{proof}
若 $K\subset A\subset G$, 则 $\lambda(K)\le\lambda(G)$, 取上、下确界得到内外测度间的不等式. 单调性同样直接得到. 若 $\sum_k\lambda^*(A_k)=\infty$, 次可加不等式无需证明. 否则选 $G_k\supset A_k$, 使 $\lambda(G_k)<\lambda^*(A_k)+\eps2^{-k}$. 开集次可加性给出 $\lambda^*(\bigcup_kA_k)\le\sum_k\lambda^*(A_k)+\eps$, 再令 $\eps\downarrow0$.

在不交情形, 选紧集 $K_k\subset A_k$; 对有限并利用紧集可加性, 得 $\lambda_*(\bigcup_kA_k)\ge\sum_{k=1}^N\lambda(K_k)$. 分别取上确界, 再令 $N\to\infty$.

若 $G$ 开, 自身可用作外逼近, 所以 $\lambda^*(G)=\lambda(G)$. 其内部的矩形多面集是紧集, 因而 $\lambda_*(G)\ge\lambda(G)$, 其余方向已知. 若 $K$ 紧, 内测度取到 $K$ 自身, 外测度正是紧集测度的定义.
\end{proof}

\begin{theorem}[有限测度的逼近]\label{ra:c3:finiteapprox}
设 $\lambda^*(A)<\infty$. 则 $A\in\mathcal L_0$ 当且仅当对每个 $\eps>0$, 都有紧集 $K$ 与开集 $G$, 使 $K\subset A\subset G$ 且 $\lambda(G\setminus K)<\eps$.
\end{theorem}
\begin{proof}
若内外测度相等, 选 $G,K$ 使 $\lambda(G)<\lambda(A)+\eps/2$ 且 $\lambda(K)>\lambda(A)-\eps/2$. 引理 \ref{ra:c3:compactgap} 给出 $\lambda(G\setminus K)=\lambda(G)-\lambda(K)<\eps$.

反之, $\lambda^*(A)\le\lambda(G)=\lambda(K)+\lambda(G\setminus K)<\lambda_*(A)+\eps$. 由于 $\eps$ 任意, 内外测度相等.
\end{proof}

\begin{proposition}\label{ra:c3:finitealgebra}
若 $A,B\in\mathcal L_0$, 则 $A\cup B,A\cap B,A\setminus B\in\mathcal L_0$. 若 $A\cap B=\varnothing$, 则 $\lambda(A\cup B)=\lambda(A)+\lambda(B)$.
\end{proposition}
\begin{proof}
取 $K_1\subset A\subset G_1$, $K_2\subset B\subset G_2$, 两个开紧间隙的测度均小于 $\eps/2$. 并集用 $K_1\cup K_2$ 与 $G_1\cup G_2$ 夹住; 交集用 $K_1\cap K_2$ 与 $G_1\cap G_2$ 夹住; 差集用紧集 $K_1\setminus G_2$ 与开集 $G_1\setminus K_2$ 夹住. 三种夹层都包含在 $(G_1\setminus K_1)\cup(G_2\setminus K_2)$ 中. 各集合的外测度有限, 故逼近定理适用.

不交时, 外测度次可加性和内测度超可加性给出
\[
 \lambda^*(A\cup B)\le\lambda(A)+\lambda(B)
 \le\lambda_*(A\cup B)\le\lambda^*(A\cup B).
\]
所有不等式都必须成为等式.
\end{proof}

\begin{proposition}\label{ra:c3:finitecountable}
若 $A_k\in\mathcal L_0$ 且 $\lambda^*(\bigcup_kA_k)<\infty$, 则其并仍在 $\mathcal L_0$, 且 $\lambda(\bigcup_kA_k)\le\sum_k\lambda(A_k)$. 两两不交时等号成立.
\end{proposition}
\begin{proof}
不交时, 将上一命题的内外测度夹逼换成可数版本即可. 一般情形令 $B_1=A_1$, $B_k=A_k\setminus\bigcup_{j<k}A_j$. 有限代数性质保证 $B_k\in\mathcal L_0$, 它们两两不交且并等于原并. 再用 $B_k\subset A_k$.
\end{proof}

\Needspace{7\baselineskip}
\originalsection{全部可测集与正则性}

\begin{definition}\label{ra:c3:allmeasurable}
若对所有 $M\in\mathcal L_0$ 都有 $A\cap M\in\mathcal L_0$, 则称 $A$ 为 Lebesgue 可测集, 记为 $A\in\mathcal L$. 定义 $\lambda(A)=\sup_{M\in\mathcal L_0}\lambda(A\cap M)$.
\end{definition}

这个定义通过有限测度窗口考察可能无界或无限测度的集合. 先确认它没有改变此前的定义.

\begin{proposition}\label{ra:c3:extension}
$\mathcal L_0\subset\mathcal L$, 两个测度定义在 $\mathcal L_0$ 上一致. 若 $A\in\mathcal L$ 且新定义的 $\lambda(A)<\infty$, 则 $A\in\mathcal L_0$.
\end{proposition}
\begin{proof}
有限代数性质给出第一个包含关系. 对 $A\in\mathcal L_0$, 一方面 $\lambda(A\cap M)\le\lambda(A)$; 另一方面可在上确界中取 $M=A$, 故定义一致.

设新测度为有限数 $c$. 令 $Q_k=[-k,k]^n$, $A_k=A\cap Q_k\in\mathcal L_0$. 则 $\lambda(A_k)\le c$. 令 $D_1=A_1$, $D_k=A_k\setminus A_{k-1}$; 它们两两不交, 有 $\sum_{k=1}^N\lambda(D_k)=\lambda(A_N)\le c$. 外测度次可加性因而给出 $\lambda^*(A)\le\sum_k\lambda(D_k)\le c<\infty$. 现在才能使用命题 \ref{ra:c3:finitecountable}, 得 $A\in\mathcal L_0$. 这补上了原稿从有限新测度跳到有限外测度的步骤.
\end{proof}

\begin{theorem}\label{ra:c3:measure}
$\mathcal L$ 是包含 Borel $\sigma$ 代数的 $\sigma$ 代数, $\lambda$ 是其上的完备测度. 对两两不交的 $A_k\in\mathcal L$, 有 $\lambda(\bigcup_kA_k)=\sum_k\lambda(A_k)$.
\end{theorem}
\begin{proof}
对 $M\in\mathcal L_0$, 有 $A^c\cap M=M\setminus(A\cap M)\in\mathcal L_0$, 故补集可测. 又有 $(\bigcup_kA_k)\cap M=\bigcup_k(A_k\cap M)$, 这个并外测度不超过 $\lambda(M)<\infty$, 故由有限可数并性质得到可测性.

设 $A_k$ 不交, $A=\bigcup_kA_k$. 对任意 $M\in\mathcal L_0$, $\lambda(A\cap M)=\sum_k\lambda(A_k\cap M)\le\sum_k\lambda(A_k)$; 取上确界得一个方向. 固定 $N$, 任取 $M_1,\ldots,M_N\in\mathcal L_0$, 令 $M=\bigcup_{j=1}^NM_j$. 则
\[
 \lambda(A)\ge\lambda(A\cap M)
 \ge\sum_{k=1}^N\lambda(A_k\cap M_k).
\]
分别取上确界, 再令 $N\to\infty$, 得另一方向.

若 $G$ 开, 则 $G\cap(-k,k)^n$ 是有限测度开集, 属于 $\mathcal L_0$. 它们的并是 $G$, 所以开集可测, 从而全部 Borel 集可测. 若 $\lambda^*(N)=0$, 则内测度也为 $0$, 故 $N\in\mathcal L_0$. 每个 $A\subset N$ 的外测度也为 $0$, 因而可测且测度为 $0$. 特别地, 可测零集的任何子集可测, 这就是完备性.
\end{proof}

\begin{theorem}[正则逼近]\label{ra:c3:regular}
对任意 $A\subset\R^n$, 以下条件等价: $A\in\mathcal L$; 对每个 $\eps>0$, 存在闭集 $F$ 和开集 $G$, 使 $F\subset A\subset G$ 且 $\lambda(G\setminus F)<\eps$. 对可测集, 无论测度有限与否, 都有 $\lambda_*(A)=\lambda(A)=\lambda^*(A)$.
\end{theorem}
\begin{proof}
先设 $A$ 可测. 令 $E_k=\{x:k-1\le\norm{x}<k\}$ ($k\ge1$), 它们是有界 Borel 集, 两两不交且覆盖 $\R^n$. 对有限可测集 $A\cap E_k$, 选 $K_k\subset A\cap E_k\subset G_k$, 使 $\lambda(G_k\setminus K_k)<\eps2^{-k}$. 令 $F=\bigcup_kK_k$, $G=\bigcup_kG_k$. $G$ 开; 而 $K_k$ 落在越来越远的壳层内, 所以其族局部有限: 每个点的某个有界邻域只遇到有限多个 $K_k$. 局部有限闭集族的并是闭集, 因而 $F$ 闭. 又有 $G\setminus F\subset\bigcup_k(G_k\setminus K_k)$, 所以夹层测度小于 $\eps$.

反之, 对每个 $j$ 取 $F_j\subset A\subset G_j$, 夹层测度小于 $1/j$. 令 $B=\bigcup_jF_j$. $B$ 为 Borel 集, $B\subset A$, 且 $A\setminus B\subset G_j\setminus F_j$, 因而 $\lambda^*(A\setminus B)=0$. 由完备性, $A=B\cup(A\setminus B)$ 可测.

若 $\lambda(A)<\infty$, 命题 \ref{ra:c3:extension} 已给出内外测度相等. 若 $\lambda(A)=\infty$, 则由测度的从下连续性, $\lambda(A\cap Q_k)\uparrow\infty$. 各截断集均可用紧集从内逼近, 所以 $\lambda_*(A)=\infty$; 由 $\lambda_*\le\lambda^*$, 外测度也为无限大.
\end{proof}

\begin{corollary}\label{ra:c3:borelcompletion}
每个 Lebesgue 可测集 $A$ 都有 $F_\sigma$ 集 $B$ 与 $G_\delta$ 集 $C$, 使 $B\subset A\subset C$ 且 $\lambda(C\setminus B)=0$. 因而 Lebesgue $\sigma$ 代数是 Borel 测度空间的完备化.
\end{corollary}
\begin{proof}
取前一定理中的 $F_j,G_j$, 使夹层测度小于 $1/j$, 令 $B=\bigcup_jF_j$, $C=\bigcap_jG_j$. 则 $C\setminus B\subset G_j\setminus F_j$ 对每个 $j$ 都成立, 故它是 Borel 零集. 于是 $A$ 是 Borel 集 $B$ 与一个 Borel 零集的子集之并. 反向包含由完备性得到.
\end{proof}

\originalsection{Carath\'eodory 判据}

\begin{lemma}\label{ra:c3:complementidentity}
若 $B$ 可测且 $\lambda(B)<\infty$, $A\subset B$ 任意, 则 $\lambda^*(A)+\lambda_*(B\setminus A)=\lambda(B)$.
\end{lemma}
\begin{proof}
任取开集 $G\supset A$. 则 $B\setminus G\subset B\setminus A$, 所以
\[
 \lambda(G)+\lambda_*(B\setminus A)
 \ge\lambda(B\cap G)+\lambda(B\setminus G)=\lambda(B).
\]
取下确界得一个方向. 另一方面, 对任意紧集 $K\subset B\setminus A$, 有 $A\subset B\setminus K$, 从而 $\lambda^*(A)+\lambda(K)\le\lambda(B\setminus K)+\lambda(K)=\lambda(B)$. 对 $K$ 取上确界即得结论.
\end{proof}

\begin{theorem}[Carath\'eodory 判据]\label{ra:c3:caratheodory}
集合 $A\subset\R^n$ 可测, 当且仅当对每个 $E\subset\R^n$, 都有
\[
 \lambda^*(E)=\lambda^*(E\cap A)+\lambda^*(E\cap A^c).
\]
\end{theorem}
\begin{proof}
若 $A$ 可测, 对任意开集 $G\supset E$, 有 $\lambda(G)=\lambda(G\cap A)+\lambda(G\cap A^c)\ge\lambda^*(E\cap A)+\lambda^*(E\cap A^c)$. 取下确界得 $\ge$; 外测度次可加性给出 $\le$.

反之, 任取 $M\in\mathcal L_0$, 记 $D=M\cap A$. 判据给出 $\lambda(M)=\lambda^*(D)+\lambda^*(M\setminus D)$. 而引理 \ref{ra:c3:complementidentity} 应用于 $M\setminus D\subset M$, 给出 $\lambda(M)=\lambda^*(M\setminus D)+\lambda_*(D)$. 所有数值有限, 可以相减, 得 $\lambda^*(D)=\lambda_*(D)$. 因而 $D\in\mathcal L_0$, 即 $A$ 可测.
\end{proof}

判据只使用外测度, 因而可测性也能从外测度单独定义. 它所要求的不是 $A$ 自己有多大, 而是用 $A$ 切分任何测试集时, 外测度都没有损失.

\begin{example}\label{ra:c3:cantor}
\textbf{编者补充.} Cantor 集是零测集, 但存在 Lebesgue 可测而非 Borel 的集合.
\end{example}
\begin{solution}
从 $[0,1]$ 反复去掉各剩余区间的中间三分之一. 第 $m$ 步剩余 $2^m$ 个长度为 $3^{-m}$ 的闭区间, 得闭集 $C_m$, $\lambda(C_m)=(2/3)^m$. Cantor 集 $C=\bigcap_mC_m$ 为紧集, 单调性给出 $\lambda(C)=0$. 因而 $C$ 的所有子集都 Lebesgue 可测.

\textbf{编者补充.} 用只含 $0,2$ 的三进制展开将二元无限序列映到 $C$, 可知 $|C|=\mathfrak c$, 其中 $\mathfrak c=|\R|$. 不同二元序列的三进制和不同: 若首个不同位置为 $m$, 该位差为 $2\,3^{-m}$, 后续最大总差为 $3^{-m}$, 仍留下正差. $\R$ 的 Borel 集至多有 $\mathfrak c$ 个: 开集由可数有理区间基决定, 每个 Borel 集可用可数节点的并、补运算树编码, 这样的树及其基元素标记总共只有 $\mathfrak c$ 种. 但 Cantor 定理给出 $|\mathcal P(C)|=2^{\mathfrak c}>\mathfrak c$. 所以 $C$ 有非 Borel 子集; 它们仍因零测而 Lebesgue 可测.
\end{solution}

\originalsection{平移、伸缩与不可测集}

\begin{proposition}\label{ra:c3:translations}
对 $z\in\R^n$ 和 $t>0$, 内外测度满足 $\lambda^*(z+A)=\lambda^*(A)$, $\lambda_*(z+A)=\lambda_*(A)$, 以及 $\lambda^*(tA)=t^n\lambda^*(A)$, $\lambda_*(tA)=t^n\lambda_*(A)$. 可测性在这两种变换下保持, 可测集的测度也服从这些公式.
\end{proposition}
\begin{proof}
对矩形, 公式就是边长的计算. 共同网格分割使它推广到矩形多面集. 平移和正伸缩将开集、紧集分别双射到同类集合, 于是对开集的上确界、紧集的下确界、以及任意集的内外逼近依次传递公式. 变换前后有限内外测度相等的条件等价, 所以有限可测性保持. 对一般 $A$, 可用有限测度窗口的逆像检验定义 \ref{ra:c3:allmeasurable}; 逆像窗口仍在 $\mathcal L_0$, 因而一般可测性也保持.
\end{proof}

\begin{theorem}[Vitali 不可测集]\label{ra:c3:vitali}
在选择公理下, $\R^n$ 存在不可测子集. 更强地, 每个正测度可测集都含有不可测子集.
\end{theorem}
\begin{proof}
定义 $x\sim y$ 当且仅当 $x-y\in\mathbb Q^n$. 这是等价关系. 从每个等价类中选一个代表, 得集合 $V$. 枚举 $\mathbb Q^n=\{q_1,q_2,\ldots\}$, 则 $\R^n=\bigsqcup_k(q_k+V)$. 若 $\lambda^*(V)=0$, 外测度次可加性与平移不变性将给出 $\lambda^*(\R^n)=0$, 矛盾. 故 $\lambda^*(V)>0$.

任取紧集 $K\subset V$, 并取有界的无限有理向量集 $D\subset\mathbb Q^n$. 集合 $q+K$ ($q\in D$) 两两不交, 而其可数并有界且可测. 因此
\[
 \infty>\lambda\Bigl(\bigcup_{q\in D}(q+K)\Bigr)
 =\sum_{q\in D}\lambda(K).
\]
所以 $\lambda(K)=0$. 对全部 $K$ 取上确界, 得 $\lambda_*(V)=0<\lambda^*(V)$; 可测集的内外测度必相等, 故 $V$ 不可测.

若 $A$ 可测且 $\lambda(A)>0$, 则 $A=\bigcup_k[A\cap(q_k+V)]$. 次可加性说明其中至少一个集合 $B$ 外测度为正. 它的内测度不超过 $\lambda_*(q_k+V)=0$, 所以 $B$ 不可测. 这就是所求子集.
\end{proof}

\begin{remark}
原第 9 讲推论的最后一步混用了内外测度. 上面的论证使用的是所选 $B$ 的外测度为正、内测度为零; 并不声称外测度可以被内测度控制. 在 $\R^n$ 中使用 $\mathbb Q^n$ 也补全了原稿只写实数代表元的维数问题.
\end{remark}

\originalsection{线性变换与体积}

\begin{lemma}\label{ra:c3:dyadic}
每个开集 $G\subset\R^n$ 都可写成两两不交的半开立方体的可数并, 各立方体都是 $J=[0,1)^n$ 的平移和正伸缩.
\end{lemma}
\begin{proof}
先选闭包包含在 $G$ 内的全部单位网格立方体. 再在边长 $1/2$ 的网格中, 选闭包包含在 $G$ 内、且尚未被已选立方体覆盖的立方体; 对边长 $2^{-m}$ 依次继续. 同层立方体不交, 两个不同层的二进网格立方体若相交, 细者包含在粗者中, 所以已选立方体仍两两不交. 对每个 $x\in G$, 当网格足够细时, 唯一包含 $x$ 的半开立方体的闭包落在 $G$ 内. 它或者被选中, 或者已被较粗者覆盖. 故所有点都被覆盖. 每层可数, 所以总族可数.
\end{proof}

\begin{lemma}\label{ra:c3:linearconstant}
若 $T:\R^n\to\R^n$ 可逆, 令 $\rho_T=\lambda(TJ)$, 则 $0<\rho_T<\infty$, 并且对任意 $A\subset\R^n$, 有 $\lambda^*(TA)=\rho_T\lambda^*(A)$ 及 $\lambda_*(TA)=\rho_T\lambda_*(A)$. 可测性在 $T$ 下保持.
\end{lemma}
\begin{proof}
$J$ 是递增紧集 $[0,1-1/k]^n$ ($k\ge2$) 的并, 所以 $TJ$ 为 Borel 集; 它有界且包含非空开集, 故 $0<\rho_T<\infty$. 由平移伸缩公式, 若 $Q=z+tJ$, 则 $\lambda(TQ)=t^n\lambda(TJ)=\rho_T\lambda(Q)$. 对开集使用上一引理的可数不交分解, 得 $\lambda(TG)=\rho_T\lambda(G)$. $T$ 为同胚, 将开邻域与紧子集双射到相应的开邻域与紧子集, 所以内外测度公式成立. 有限可测性由内外测度相等得到; 一般可测性再用有限窗口逆像检验.
\end{proof}

\begin{theorem}\label{ra:c3:determinant}
若线性映射 $T$ 可逆, 则对所有 $A\subset\R^n$, 有 $\lambda^*(TA)=|\det T|\lambda^*(A)$, $\lambda_*(TA)=|\det T|\lambda_*(A)$. 若 $A$ 可测, 则 $TA$ 可测且 $\lambda(TA)=|\det T|\lambda(A)$. 若 $T$ 奇异, 则任意 $TA$ 都是可测零集.
\end{theorem}
\begin{proof}
只需识别 $\rho_T$. 交换两个坐标保持 $J$ 的体积. 将第 $i$ 个坐标乘以 $a\ne0$ 时, $TJ$ 为半开矩形, 体积为 $|a|$. 半开和闭矩形只差有限个超平面片, 后者可用任意薄的矩形覆盖, 因而为零测集.

剩下的初等变换是剪切 $Sx=x+a x_j e_i$, $i\ne j$. 将 $J$ 沿第 $j$ 个坐标分成 $m$ 个薄片. 在每个薄片内, 第 $i$ 个坐标的位移变化至多 $|a|/m$. 因而每个像薄片包含一个体积为 $(1-|a|/m)/m$ 的矩形, 并包含在一个体积为 $(1+|a|/m)/m$ 的矩形中, 此处取 $m>|a|$. 不同片在第 $j$ 个方向的内部互不相交. 求和得到
\[
 1-|a|/m\le\lambda(SJ)\le1+|a|/m.
\]
令 $m\to\infty$, 得 $\rho_S=1$. 这是原稿 ``检查初等矩阵'' 所省略的剪切计算.

可逆矩阵通过高斯消元分解为上述初等矩阵的乘积. 引理 \ref{ra:c3:linearconstant} 给出 $\rho_{ST}=\rho_S\rho_T$, 而行列式的绝对值具有同样的乘法性质, 并在各初等矩阵上等于 $\rho$. 故 $\rho_T=|\det T|$.

若 $T$ 奇异, 其像包含在一个真线性子空间内. 坐标超平面的每个有界部分可用厚度趋于 $0$ 的矩形覆盖, 所以坐标超平面为零集. 任意超平面是它在某个可逆线性映射下的像; 由已证公式仍是零集. 因而 $TA$ 是零集的子集, 完备性给出可测性及零测度.
\end{proof}

\begin{remark}
\textbf{编者修正.} 当 $\det T=0$ 且 $\lambda^*(A)=\infty$ 时, 乘积 $0\cdot\infty$ 在通常扩展实数运算中不定义. 将奇异情形单独陈述, 就不需要这个乘积. 若 $T$ 为正交变换, 则 $|\det T|=1$, 所以旋转和反射均保持 Lebesgue 测度; $\mathrm{SL}(n,\R)$ 中的变换也保持体积.
\end{remark}

\originalsection{练习与解答}

\begin{exercise}
证明 $\R^n$ 的每个可数子集是零集, 并给出一个稠密的零集.
\end{exercise}
\begin{solution}
单点可用各边长趋于 $0$ 的立方体覆盖, 所以外测度为 $0$. 对可数并使用外测度次可加性. $\mathbb Q^n$ 可数且在 $\R^n$ 中稠密, 是所求例子. 稠密性并不保证正测度, 而非空开集确实有正测度.
\end{solution}

\begin{exercise}
设 $A$ 可测且 $\lambda(A)<\infty$. 证明对每个 $\eps>0$, 存在矩形多面集 $P$, 使 $\lambda(A\mathbin\triangle P)<\eps$.
\end{exercise}
\begin{solution}
取 $K\subset A\subset G$, 使 $\lambda(G\setminus K)<\eps$. 对 $K$ 的每个点选一个包含它、且闭包落在 $G$ 内的小开矩形; 紧性给出有限子覆盖. 将这些开矩形的闭包取并, 得 $K\subset P\subset G$. 因为 $A$ 与 $P$ 都包含 $K$ 且包含在 $G$ 内, $A\mathbin\triangle P\subset G\setminus K$. 此处也说明矩形多面集可在对称差意义下逼近有限可测集.
\end{solution}

\begin{exercise}
设 $E$ 是 $\R^n$ 中的仿射超平面. 证明 $E$ 零测, 并说明为何这不能推出每个空内部闭集都零测.
\end{exercise}
\begin{solution}
仿射超平面是坐标超平面的可逆线性像再作平移, 因而由定理 \ref{ra:c3:determinant} 和命题 \ref{ra:c3:translations} 零测. 为说明后半句, 在 $[0,1]$ 逐步从每个剩余闭区间内部删去一个非空开区间, 使第 $k$ 轮删除的总长度不超过 $2^{-k-2}$, 并选择删去区间使剩余区间的最大长度趋于 $0$. 所得闭集不含区间, 因而内部为空; 全部删去长度不超过 $\sum_{k\ge1}2^{-k-2}=1/4$, 所以剩余测度至少 $3/4$. 例如每轮在各区间中点删去足够短的区间, 剩余单个区间长度至多上一轮最大长度的一半, 即满足要求. 这是一个正测度的胖 Cantor 集.
\end{solution}

\originalchapter{正则测度与函数逼近}
\label{ra:c4:chapter}
\sourcelecture{10--13}

\originalsection{连续函数与局部紧空间}

积分是线性的: 对可积函数 $f,h$ 与标量 $a,b$, 有 $\int(af+bh)\dd\mu=a\int f\dd\mu+b\int h\dd\mu$. 若 $f\ge0$, 则积分非负. 原第 10 讲反过来问: 一个定义在连续函数上的正线性泛函, 能否由某个测度的积分给出? 这既解释正则测度的来源, 也提供后面逼近可测函数所需的开集、紧集工具.

\begin{definition}\label{ra:c4:topology}
拓扑空间 $X$ 称为 Hausdorff 空间, 若任意两个不同点有不交的开邻域. 它称为局部紧空间, 若每个点都有闭包紧的开邻域. 函数的支集是 $\supp f=\overline{\{x:f(x)\ne0\}}$. 记 $C_c(X)$ 为复值连续且具有紧支集的函数所成的向量空间.

实值或扩展实值函数 $v$ 称为下半连续, 若每个 $\{v>a\}$ 都开; $u$ 称为上半连续, 若每个 $\{u<a\}$ 都开. 开集 $V$ 的示性函数 $\ind_V$ 下半连续, 闭集 $F$ 的示性函数 $\ind_F$ 上半连续.
\end{definition}

\begin{remark}
支集定义中的闭包不可省略. $f$ 连续且具有紧支集时, 它在支集上有界. $C_c(X)$ 对线性组合封闭, 因为新支集包含在原来两个紧支集的并内. 本节只讨论测度表示所需的正泛函; 一般 Banach 空间对偶与谱理论留给泛函分析册.
\end{remark}

\begin{definition}\label{ra:c4:cutoffnotation}
记 $K\prec f$ 表示 $K$ 紧, $f\in C_c(X)$, $0\le f\le1$, 且 $f=1$ 于 $K$. 记 $f\prec V$ 表示 $V$ 开, $f\in C_c(X)$, $0\le f\le1$, 且 $\supp f\subset V$. 正线性泛函是复线性映射 $\Lambda:C_c(X)\to\C$, 满足 $f\ge0\Rightarrow\Lambda f\ge0$.
\end{definition}

正性蕴含实函数被送到实数, 以及 $f\le h\Rightarrow\Lambda f\le\Lambda h$: 对实函数用连续的正负部分分解即可. 若 $\phi=1$ 于 $\supp f$ 且 $\phi\ge0$, 则对实函数有 $|\Lambda f|\le\norm{f}_\infty\Lambda\phi$. 因而泛函在固定紧支集上的大小可用一个截断函数控制; 不需要先假设全空间上的统一有界性.

\begin{lemma}[局部紧空间的截断函数]\label{ra:c4:urysohn}
设 $X$ 局部紧且 Hausdorff, $K\subset V$, $K$ 紧、$V$ 开. 则存在 $f\in C_c(X)$, 使 $K\prec f\prec V$. 还可以使 $f=1$ 于 $K$ 的某个开邻域.
\end{lemma}
\begin{proof}
\textbf{编者补充原稿所引用的拓扑论证.} 紧 Hausdorff 空间是正规空间. 确切地说, 对不交闭集 $A,B$, 它们也是紧集. 对每个 $a\in A$ 和 $b\in B$ 选不交邻域; 先对 $b$ 的邻域取有限子覆盖, 得到 $a$ 的一个邻域与 $B$ 的一个邻域不交, 再对 $a$ 的邻域取有限子覆盖, 得到整个 $A,B$ 的不交开邻域. 因而在紧 Hausdorff 空间内, 若闭集 $A\subset U$ 且 $U$ 开, 可选开集 $W$ 使 $A\subset W\subset\overline W\subset U$.

这个收缩性质也给出局部紧 Hausdorff 空间的相应局部性质. 对 $x\in V$, 先取闭包紧的开邻域 $H^\circ_x$, 在紧空间 $H_x=\overline{H^\circ_x}$ 中将 $\{x\}$ 的邻域收缩到 $H^\circ_x\cap V$ 内. 所得开邻域在 $X$ 中也是开集, 闭包紧且包含在 $V$ 内. 用紧性对 $K$ 取有限覆盖, 得到开集 $O$, 满足 $K\subset O\subset\overline O\subset V$, $\overline O$ 紧. 再收缩一次, 可在 $O$ 内取得 $K$ 的紧邻域 $L$.

在紧正规空间 $H=\overline O$ 上应用下面的二进有理数构造. 对不交闭集 $L$ 和 $H\setminus O$, 反复使用收缩性质, 选开集 $U_r$ ($r\in[0,1]\cap\mathbb Z[1/2]$), 使 $L\subset U_0$, $U_1=H\setminus(H\setminus O)=O$, 且 $r<s\Rightarrow\overline{U_r}\subset U_s$. 在插入每个相邻二进有理数的中点时, 只需将 $\overline{U_r}$ 的邻域收缩到 $U_s$ 内. 令 $h(x)=\inf\{r:x\in U_r\}$, 空集的下确界取 $1$. 则 $h=0$ 于 $L$, $h=1$ 于 $H\setminus O$. 对 $0<a<1$, 有 $\{h<a\}=\bigcup_{r<a}U_r$; 同时, 利用闭包的严格嵌套, $\{h>a\}=\bigcup_{r>a}(H\setminus\overline{U_r})$. 两者均开. 端点阈值 $a=0,1$ 的非平凡水平集可由这些开集取并得到; $[0,1]$ 外的阈值给出全空间或空集. 因而 $h$ 连续. 令 $f=1-h$ 于 $H$, 并在 $X\setminus H$ 上取 $0$. 它在 $H\setminus O$ 上已为 $0$, 因而延拓连续; 支集包含在紧集 $H\subset V$ 内, 且在紧邻域 $L$ 上等于 $1$.
\end{proof}

\begin{lemma}[有限单位分解]\label{ra:c4:partition}
若紧集 $K\subset V_1\cup\cdots\cup V_m$, 各 $V_i$ 开, 则存在 $h_i\prec V_i$, 使 $\sum_ih_i=1$ 于 $K$, 且 $0\le\sum_ih_i\le1$ 于整个 $X$.
\end{lemma}
\begin{proof}
对 $K$ 各点选支集落在某个 $V_i$ 内、且在该点邻域为 $1$ 的截断函数. 紧性给出有限个此类函数. 将属于同一 $V_i$ 的函数相加, 得非负连续紧支集函数 $g_i$, 使 $s=\sum_ig_i>0$ 于 $K$. 再选 $\psi$ 满足 $K\prec\psi\prec\{s>0\}$. 在 $\{s>0\}$ 上令 $h_i=\psi g_i/s$, 其他地方取 $0$. 因 $\supp\psi$ 是包含在 $\{s>0\}$ 内的紧集, 延拓连续. 各支集落在 $V_i$ 内, 且 $\sum_ih_i=\psi$, 所以满足结论.
\end{proof}

\begin{remark}
原稿的单位分解应理解为在 $K$ 上求和为 $1$. 在非紧空间上, 有限个紧支集函数不能在全空间求和为 $1$. 后面的证明同时使用 $K$ 上的等式与全空间的上界.
\end{remark}

\originalsection{Riesz 表示定理的测度构造}

\begin{theorem}[正线性泛函的 Riesz 表示]\label{ra:c4:riesz}
设 $X$ 为局部紧 Hausdorff 空间, $\Lambda:C_c(X)\to\C$ 为正线性泛函. 则存在包含全部 Borel 集的完备 $\sigma$ 代数 $\mathcal M$ 及其上的正测度 $\mu$, 使得:
\begin{enumerate}
\item 对每个 $f\in C_c(X)$, $\Lambda f=\int_X f\dd\mu$.
\item 每个紧集 $K$ 的测度有限.
\item 对 $E\in\mathcal M$, $\mu(E)=\inf\{\mu(V):E\subset V,\ V\text{ 开}\}$.
\item 若 $E$ 开, 或 $E\in\mathcal M$ 且 $\mu(E)<\infty$, 则 $\mu(E)=\sup\{\mu(K):K\subset E,\ K\text{ 紧}\}$.
\end{enumerate}
可取 $\mathcal M$ 为下述外测度的全部 Carath\'eodory 可测集. 在这个固定 $\mathcal M$ 上, 满足这些性质的测度唯一.
\end{theorem}

原第 10 讲在一般局部紧 Hausdorff 空间中陈述此定理, 并给出构造及积分表示的提纲. 以下保留这个范围, 将原稿略去的测度步骤补齐. 唯一性所比较的是正则的表示测度; 不能将 ``唯一测度'' 理解成允许任意改变可测集域后仍有字面上的同一个对象.

\begin{proof}
对开集 $V$, 定义 $m(V)=\sup\{\Lambda f:f\prec V\}$, 对任意 $E\subset X$ 定义 $\mu^*(E)=\inf\{m(V):E\subset V,\ V\text{ 开}\}$. 以下依次验证外测度、可测性、正则性和积分表示.

\textbf{开集上的运算.} $m$ 单调, $m(\varnothing)=0$. 若 $V\subset\bigcup_iV_i$, 对任意 $f\prec V$, 用 $\supp f$ 的紧性从 $V_i$ 中取有限子覆盖, 再用有限单位分解 $h_i$. 有 $f=\sum_ifh_i$, 各 $fh_i\prec V_i$, 所以 $\Lambda f\le\sum_im(V_i)$. 取上确界得 $m(V)\le\sum_im(V_i)$. 若 $V_i$ 两两不交, 固定有限个 $f_i\prec V_i$, 则 $\sum_if_i\prec\bigcup_iV_i$, 故 $m(\bigcup_iV_i)\ge\sum_im(V_i)$, 先取有限个上确界再增加项数. 因而不交开集可数可加.

$\mu^*$ 是外测度: 空集与单调性直接成立; 可数次可加性用每个 $E_i$ 的开邻域 $V_i$, 使 $m(V_i)<\mu^*(E_i)+\eps2^{-i}$, 然后使用上一段. 若右边为无限大, 不等式自动成立. 对开集, 单调性并允许取 $V$ 自身, 给出 $\mu^*(V)=m(V)$.

\textbf{紧集与开集的关系.} 对紧集 $K$, 有
\begin{equation}\label{ra:c4:compactcapacity}
 \mu^*(K)=\inf\{\Lambda f:K\prec f\}<\infty.
\end{equation}
首先选 $\phi$ 满足 $K\prec\phi$. 在 $U=\{\phi>1/2\}$ 上, 任意 $h\prec U$ 都有 $h\le2\phi$, 故 $m(U)\le2\Lambda\phi<\infty$. 这证明紧集的外测度有限. 对任意 $K\prec f$ 及 $0<\delta<1$, 开集 $U_\delta=\{f>1-\delta\}$ 包含 $K$, 且 $h\prec U_\delta$ 时 $h\le f/(1-\delta)$. 因而 $\mu^*(K)\le\Lambda f/(1-\delta)$; 令 $\delta\downarrow0$, 得 $\mu^*(K)\le\Lambda f$. 反之, 任意开集 $V\supset K$ 都允许选 $K\prec f\prec V$, 于是 $\Lambda f\le m(V)$. 取下确界即得 \eqref{ra:c4:compactcapacity}.

开集还满足
\begin{equation}\label{ra:c4:openinner}
 m(V)=\sup_{K\subset V,\ K\text{ 紧}}\mu^*(K).
\end{equation}
一个方向来自单调性. 另一个方向中, 若 $f\prec V$, 令 $K=\supp f$. 任何 $K\prec h$ 都有 $f\le h$, 所以 $\Lambda f\le\mu^*(K)$; 再对 $f$ 取上确界.

对 $K\subset V$, $K$ 紧、$V$ 开, 进一步有
\begin{equation}\label{ra:c4:opensplit}
 m(V)=\mu^*(K)+m(V\setminus K).
\end{equation}
为证 $\ge$, 任取 $h\prec V\setminus K$, 再选 $K\prec f\prec V\setminus\supp h$. $f,h$ 的支集不交, 所以 $f+h\prec V$, 从而 $m(V)\ge\mu^*(K)+\Lambda h$. 取上确界. 为证 $\le$, 给定 $\eps>0$, 选开集 $W\supset K$, 使 $m(W)<\mu^*(K)+\eps$. 在 $W\cap V$ 内取 $K$ 的紧邻域 $L$, 再取 $L\prec f\prec W\cap V$. 对任意 $h\prec V$, 有 $hf\prec W$, 而 $h(1-f)\prec V\setminus K$, 因为 $f$ 在 $K$ 的一个邻域内等于 $1$. 所以 $\Lambda h\le m(W)+m(V\setminus K)$. 取上确界再令 $\eps\downarrow0$.

\textbf{Borel 集的可测性.} 设 $U$ 开, $A\subset X$ 任意, $V\supset A$ 开. 对紧集 $K\subset V\cap U$, 式 \eqref{ra:c4:opensplit} 给出
\[
 m(V)\ge\mu^*(K)+\mu^*(A\setminus U).
\]
利用 \eqref{ra:c4:openinner} 对 $K$ 取上确界, 得 $m(V)\ge m(V\cap U)+\mu^*(A\setminus U)\ge\mu^*(A\cap U)+\mu^*(A\setminus U)$. 再对 $V$ 取下确界. 相反方向由次可加性给出, 因而每个开集满足 Carath\'eodory 等式.

这里也补全一般外测度的延拓步骤. 对任意外测度 $\nu^*$, 称 $E$ 可测若它对每个 $A$ 都满足 $\nu^*(A)=\nu^*(A\cap E)+\nu^*(A\setminus E)$. 补集封闭直接成立. 若 $E,F$ 可测, 先用 $E$ 切分 $A$, 再用 $F$ 切分剩余部分, 得
\[
 \nu^*(A)\ge\nu^*(A\cap(E\cup F))+\nu^*(A\setminus(E\cup F)),
\]
因为前两片的外测度之和不小于其并的外测度; 反向不等式总由次可加性成立. 这证明有限并封闭, 从而差集也封闭. 若 $E_i$ 两两不交且可测, 依次切分得
\[
 \nu^*(A)\ge\sum_{i=1}^N\nu^*(A\cap E_i)
       +\nu^*\Bigl(A\setminus\bigcup_{i\ge1}E_i\Bigr).
\]
令 $N\to\infty$, 再利用 $\nu^*(A\cap\bigcup_iE_i)\le\sum_i\nu^*(A\cap E_i)$, 得到可数并可测. 一般可数并先将其差集化为不交并. 对 $A=\bigcup_iE_i$ 使用同一有限切分及次可加性, 得到可数可加性. 外测度零集及其所有子集也满足切分等式, 所以所得测度完备.

将此论证应用于 $\mu^*$, 令 $\mathcal M$ 为其可测集族, $\mu=\mu^*|_{\mathcal M}$. 已证 $\mathcal M$ 包含 Borel 集, 测度完备, 紧集有限; 外正则性由定义直接成立. 开集内正则性就是 \eqref{ra:c4:openinner}.

\textbf{有限可测集的内正则性.} 设 $\mu(E)<\infty$. 给定 $\eta>0$, 取开集 $V\supset E$ 使 $\mu(V)<\mu(E)+\eta$, 再取紧集 $K\subset V$ 使 $\mu(K)>\mu(V)-\eta$. 于是 $\mu(E\setminus K)<\eta$. 取开集 $W\supset K\setminus E$, 使 $\mu(W)<\mu(K\setminus E)+\eta$. 紧集 $L=K\setminus W$ 包含在 $E$ 内, 且
\[
 \mu(E\setminus L)
 \le\mu(E\setminus K)+\mu(E\cap K\cap W)<2\eta,
\]
因为 $K\setminus E\subset W$ 且 $\mu(E\cap K\cap W)\le\mu(W)-\mu(K\setminus E)<\eta$. 任意小的损失证明有限集内正则. 特别地, 对有限 $E$ 可取 $L\subset E\subset V$, 使 $\mu(V\setminus L)$ 任意小.

\textbf{积分表示.} 保留原稿按函数值分层并用单位分解的方法. 先取实函数 $f\in C_c(X)$. 若 $f=0$, 表示等式立即成立; 以下设 $f\ne0$, 令 $K=\supp f$, $M=\norm{f}_\infty>0$. 将 $[-M,M]$ 分成有限个长度至多 $\delta$ 的半开区间, 其右端点为 $y_i$, 用适当闭合的首端点约定覆盖整个范围. 令 $E_i\subset K$ 为相应函数值层; 这些 Borel 集构成 $K$ 的不交分割, 且 $y_i-\delta\le f\le y_i$ 于 $E_i$. 由外正则性与连续性, 可选开集
\[
 E_i\subset V_i\subset\{f<y_i+\delta\},
 \qquad \mu(V_i)<\mu(E_i)+\eta/r,
\]
其中 $r$ 是层数. 对覆盖 $K$ 的 $V_i$ 取单位分解 $h_i$. 令 $a=M$. 因 $\sum_ih_i=1$ 于 $K$, 有 $f=\sum_ifh_i$; 又 $K\prec\sum_ih_i$, 故式 \eqref{ra:c4:compactcapacity} 给出 $\Lambda(\sum_ih_i)\ge\mu(K)$. 正性于是给出
\begin{align*}
 \Lambda f
 &\le\sum_i(y_i+\delta)\Lambda h_i\\
 &=\sum_i(y_i+a+\delta)\Lambda h_i-a\Lambda\Bigl(\sum_ih_i\Bigr)\\
 &\le\sum_i(y_i+a+\delta)\mu(V_i)-a\mu(K)\\
 &\le\sum_iy_i\mu(E_i)+\delta\mu(K)+(2M+\delta)\eta\\
 &\le\int_X f\dd\mu+2\delta\mu(K)+(2M+\delta)\eta.
\end{align*}
$h_i$ 的支集在 $V_i$ 内, 故 $\Lambda h_i\le\mu(V_i)$; 加上 $a$ 是为了让使用这一步的系数全部非负. 先令 $\eta\downarrow0$, 再令 $\delta\downarrow0$, 得 $\Lambda f\le\int f\dd\mu$. 对 $-f$ 应用同一不等式, 得反向不等式. 实函数表示成立, 复函数再分实部、虚部, 利用复线性得到.

\textbf{唯一性.} 若 $\mu_1,\mu_2$ 都满足定理性质, 取紧集 $K$ 及开邻域 $V$ 使 $\mu_2(V)<\mu_2(K)+\eps$, 再取 $K\prec f\prec V$. 则 $\mu_1(K)\le\int f\dd\mu_1=\Lambda f=\int f\dd\mu_2\le\mu_2(V)<\mu_2(K)+\eps$. 交换两测度并令 $\eps\downarrow0$, 得紧集上相等. 开集的内正则性使它们在开集上相等, 全部可测集的外正则性使它们在固定的 $\mathcal M$ 上相等. 若比较不同的可测集域, 此论证给出共同域上相等; 两者从开集产生的外测度及其完整 Carath\'eodory 域仍相同.
\end{proof}

\begin{definition}\label{ra:c4:regularity}
包含 Borel 集的可测空间上的测度称为 Borel 测度. 用开集外逼近称为外正则, 用紧集内逼近称为内正则; 两者都成立时称为正则. $X$ 称为 $\sigma$ 紧, 若它是可数个紧集的并; 测度称为 $\sigma$ 有限, 若 $X$ 是可数个有限测度可测集的并.
\end{definition}

一般 Riesz 定理只对开集与有限可测集保证紧集内逼近. 增加 $\sigma$ 紧条件后, 可以得到对所有可测集的内正则性以及闭集逼近.

\begin{proposition}\label{ra:c4:sigmacompact}
在定理 \ref{ra:c4:riesz} 中, 若 $X$ 还 $\sigma$ 紧, 则 $\mu$ 为 $\sigma$ 有限的正则测度. 对每个 $E\in\mathcal M$ 和 $\eps>0$, 都有闭集 $F$、开集 $V$, 使 $F\subset E\subset V$ 且 $\mu(V\setminus F)<\eps$. 此外有 $F_\sigma$ 集 $A$ 和 $G_\delta$ 集 $B$, 使 $A\subset E\subset B$ 且 $\mu(B\setminus A)=0$.
\end{proposition}
\begin{proof}
由给定的可数紧覆盖和截断引理, 可递归构造紧集 $K_j$ 使 $K_j\subset\operatorname{int}K_{j+1}$ 且 $\bigcup_jK_j=X$: 每一步用有限个相对紧开邻域覆盖前一个紧集和下一给定紧集, 取其闭包的并. 紧集测度有限, 所以测度 $\sigma$ 有限. 对任意 $E$, 有 $\mu(E\cap K_j)\uparrow\mu(E)$; 有限集内正则性使全部 $E$ 都能用紧集从内逼近, 即使其测度无限也成立.

令 $K_0=\varnothing$, $S_j=K_j\setminus\operatorname{int}K_{j-1}$. 这是局部有限的紧集族且覆盖 $X$: 若 $x\in\operatorname{int}K_m$, 它的一个邻域不会遇到 $j\ge m+2$ 的壳层. 每个 $E\cap S_j$ 有限, 可选紧集 $C_j\subset E\cap S_j$ 与开集 $V_j\supset E\cap S_j$, 使 $\mu(V_j\setminus C_j)<\eps2^{-j}$. 令 $F=\bigcup_jC_j$, $V=\bigcup_jV_j$. 局部有限性保证 $F$ 闭, 而 $V\setminus F\subset\bigcup_j(V_j\setminus C_j)$ 给出误差界.

最后对误差 $1/m$ 分别取 $F_m,V_m$, 令 $A=\bigcup_mF_m$, $B=\bigcap_mV_m$. 则 $B\setminus A$ 的测度不超过每个 $1/m$, 所以为零.
\end{proof}

\begin{example}\label{ra:c4:riemannriesz}
在 $\R^n$ 上, 令 $\Lambda f$ 为紧支集连续函数的 Riemann 积分. Riesz 构造得到的测度就是 Lebesgue 测度.
\end{example}
\begin{solution}
这个泛函正且线性. 对轴平行闭矩形 $I$, 任何 $I\prec f$ 都有 $\Lambda f\ge v(I)$, 因为非负函数 $f$ 在 $I$ 上为 $1$. 又可选截断函数支集落在任意小幅扩大后的矩形内, 因而 $\Lambda f$ 至多为扩大后的体积. 由式 \eqref{ra:c4:compactcapacity}, 得 $\mu(I)=v(I)$. 退化矩形用薄邻域同样得到零测度.

有限网格分割使 $\mu(P)=v(P)$ 对矩形多面集成立. 若 $G$ 开、$K\subset G$ 紧, 可用有限个闭包落在 $G$ 内的小开矩形覆盖 $K$, 得 $K\subset P\subset G$. 因而开集内正则性给出 $\mu(G)=\sup_{P\subset G}v(P)=\lambda(G)$. 再由外正则性, 两者产生的外测度相同, 完整可测集域与测度也相同.
\end{solution}

\originalsection{Lusin 定理}

本节取定理 \ref{ra:c4:riesz} 给出的正则性条件: 测度完备, 包含 Borel 集, 紧集有限, 所有可测集外正则, 有限可测集可由紧集从内逼近. Lebesgue 测度是主要例子. ``可测函数近乎连续'' 的准确含义是删去一个任意小测度的例外集后, 它与某个连续函数完全相同, 并非仅仅函数值接近.

\begin{theorem}[Lusin]\label{ra:c4:lusin}
设 $f:X\to\C$ 可测且处处有限, $A\in\mathcal M$, $\mu(A)<\infty$, 并且 $f=0$ 于 $X\setminus A$. 对每个 $\eps>0$, 存在 $g\in C_c(X)$, 使 $\mu(\{f\ne g\})<\eps$. 若 $f$ 有界, 还可保证 $\norm{g}_\infty\le\norm{f}_\infty$.
\end{theorem}
\begin{proof}
先设 $0\le f\le1$, 且 $f=0$ 于某个紧集 $K$ 之外. 令 $s_0=0$, 对 $n\ge1$ 定义
\[
 s_n(x)=2^{-n}\min\{\lfloor2^nf(x)\rfloor,2^n-1\}.
\]
则 $s_n\uparrow f$, 且 $t_n=s_n-s_{n-1}$ 只取 $0,2^{-n}$ 两个值. 这个上端截断使 $f=1$ 时仍有一致的二进展开. 令 $T_n=\{t_n=2^{-n}\}\subset K$, 则 $f=\sum_{n\ge1}2^{-n}\ind_{T_n}$.

取开集 $V\supset K$, 使 $\overline V$ 紧. 用有限集正则逼近选紧集 $K_n$ 和开集 $V_n$, 满足 $K_n\subset T_n\subset V_n\subset V$, $\mu(V_n\setminus K_n)<\eps2^{-n}$. 取 $K_n\prec h_n\prec V_n$, 定义 $g=\sum_{n\ge1}2^{-n}h_n$. 级数一致收敛, 所以 $g$ 连续, $0\le g\le1$, 且支集包含在 $\overline V$ 内. 在 $V_n\setminus K_n$ 之外, $h_n=\ind_{T_n}$, 故在这些夹层之并之外, $g=f$. 例外集测度小于 $\eps$.

一般有界实函数分为正负部分, 分别归一化到 $[0,1]$, 应用刚才的结论并将误差各分一半. 一般复函数再分实部、虚部. 因有限个所用紧邻域的并仍紧, 得到的近似仍在 $C_c(X)$ 中. 若 $M=\norm{f}_\infty$, 对所得函数施加连续径向截断 $P_M(z)=z$ ($|z|\le M$), $P_M(z)=Mz/|z|$ ($|z|>M$), 就可使范数不超过 $M$, 而已与 $f$ 相等的点仍保持相等. $P_M(0)=0$, 所以紧支集也保持; $M=0$ 时直接取零函数.

最后设 $A$ 只是有限测度且 $f$ 可能无界. 先取紧集 $K\subset A$, 使 $\mu(A\setminus K)<\eps/3$. 因 $f$ 处处有限, $A\cap\{|f|>M\}\downarrow\varnothing$, 从上连续性允许选 $M$ 使其测度小于 $\eps/3$. 令 $f'=f\ind_{K\cap\{|f|\le M\}}$. 它有界且在紧集 $K$ 外为 $0$, 所以存在 $g\in C_c(X)$, 使 $\mu(\{f'\ne g\})<\eps/3$. $f$ 与 $f'$ 的差异只在先前两种损失集上, 故总误差小于 $\eps$. 当原 $f$ 有界时, 最后再作范数截断即可.
\end{proof}

\begin{lemma}[可求和例外集]\label{ra:c4:bc}
在任意测度空间上, 若 $\sum_n\mu(E_n)<\infty$, 则几乎每个点只属于有限多个 $E_n$.
\end{lemma}
\begin{proof}
属于无限多个 $E_n$ 的集合是 $\limsup_nE_n=\bigcap_N\bigcup_{n\ge N}E_n$. 对每个 $N$, 它的测度不超过 $\sum_{n\ge N}\mu(E_n)$, 而后者趋于 $0$. 因而该集合零测. 也可令 $H=\sum_n\ind_{E_n}$; 单调收敛给出 $\int H=\sum_n\mu(E_n)<\infty$, 从而 $H<\infty$ 几乎处处, 这就是原稿的证明.
\end{proof}

\begin{corollary}\label{ra:c4:continuousae}
在 Lusin 定理的假设下, 若 $f$ 有界, 则存在 $g_n\in C_c(X)$, $\norm{g_n}_\infty\le\norm{f}_\infty$, 使 $g_n\to f$ 几乎处处. 更强地, 几乎每个点上最终都有 $g_n(x)=f(x)$.
\end{corollary}
\begin{proof}
取 $\mu(\{g_n\ne f\})<2^{-n}$. 应用上一引理, 几乎每个点只落在有限多个不相等的例外集中, 所以后来一直相等.
\end{proof}

\originalsection{Vitali--Carath\'eodory 定理}

Lusin 定理控制不相等的集合. 下面的定理则用半连续函数从上下夹住可积函数, 并控制夹层的积分. 示性函数的开紧逼近正好提供这种半连续性.

\begin{theorem}[Vitali--Carath\'eodory]\label{ra:c4:vitalicara}
在本章采用的 Riesz 正则测度条件下, 设 $f:X\to\R$ 可测且 $f\in L^1(\mu)$. 对每个 $\eps>0$, 存在上半连续函数 $u:X\to[-\infty,\infty)$ 与下半连续函数 $v:X\to(-\infty,\infty]$, 使 $u\le f\le v$ 处处, $u$ 有有限上界, $v$ 有有限下界, 且 $\int_X(v-u)\dd\mu<\eps$. 因而它们有限几乎处处, 并可作为 $L^1$ 函数处理.
\end{theorem}
\begin{proof}
先设 $f\ge0$. 取递增非负简单函数 $s_n\uparrow f$, 将各增量 $s_n-s_{n-1}$ 的有限个正值层枚举, 得
\[
 f=\sum_{i\ge1}c_i\ind_{E_i},\qquad
 c_i>0,\qquad\sum_{i\ge1}c_i\mu(E_i)=\int f\dd\mu<\infty.
\]
每个 $E_i$ 的测度有限. 选 $K_i\subset E_i\subset V_i$, $K_i$ 紧、$V_i$ 开, 使 $c_i\mu(V_i\setminus K_i)<\eps2^{-i-1}$. 再取 $N$ 使 $\sum_{i>N}c_i\mu(E_i)<\eps/2$. 定义
\[
 u=\sum_{i=1}^Nc_i\ind_{K_i},\qquad
 v=\sum_{i\ge1}c_i\ind_{V_i}.
\]
有限个闭集示性函数的正线性组合上半连续, 故 $u$ 上半连续且有界. 非负下半连续函数的部分和递增, 其上确界仍下半连续, 故 $v$ 下半连续. 有 $0\le u\le f\le v$. 逐点分解为非负项,
\[
 v-u=\sum_{i\ge1}c_i(\ind_{V_i}-\ind_{K_i})
             +\sum_{i>N}c_i\ind_{K_i}.
\]
用单调收敛对这些项积分, 得
\[
 \int(v-u)\dd\mu
 \le\sum_{i\ge1}c_i\mu(V_i\setminus K_i)
       +\sum_{i>N}c_i\mu(E_i)<\eps.
\]
这保留原稿的有限下逼近与无限上逼近, 并补正其夹层展开中漏掉的项.

一般实函数写成 $f=f^+-f^-$. 对两部分分别应用非负情形, 各用误差 $\eps/2$, 得 $u_\pm\le f^\pm\le v_\pm$. 令 $u=u_+-v_-$, $v=v_+-u_-$. 由于 $u_\pm$ 有限有界而 $v_\pm\ge0$, 运算不含 $\infty-\infty$. 上半连续性和下半连续性分别保持, $u\le f\le v$, 且 $v-u=(v_+-u_+)+(v_--u_-)$. 积分误差相加小于 $\eps$. 由 $f\in L^1$ 和可积夹层, $|u|,|v|\le |f|+(v-u)$ 几乎处处, 故二者可积.
\end{proof}

\begin{remark}
\textbf{编者修正.} 原稿把 $u,v$ 的值域都写为 $\R$, 但上述无限上逼近允许在零测集上取 $+\infty$, 一般实函数的下逼近相应允许 $-\infty$. 本稿采用扩展实值版本, 并由积分误差证明两者有限几乎处处; 后面的积分结论不需要它们处处有限.
\end{remark}

\originalsection{几乎处处版本的积分结论}

\begin{proposition}\label{ra:c4:aedct}
设 $f_n$ 有可测代表, $f_n\to f$ 几乎处处, 且对某个 $g\in L^1(\mu)$ 有 $|f_n|\le g$ 几乎处处. 则 $f\in L^1(\mu)$, 且 $\int f_n\dd\mu\to\int f\dd\mu$; 还可得到 $\int|f_n-f|\dd\mu\to0$.
\end{proposition}
\begin{proof}
将不收敛点、各不等式失败点、以及所取代表未定义的点合为一个可测零集 $N$. 在 $N$ 上把所有函数改为 $0$, 即可使用前章处处收敛版本的控制收敛定理; 改变零集上的值不改变积分. 不等式失败点应写成 $\{|f_n|>g\}$, 原稿的 $\ge$ 会误把等号成立的点也划入例外集. 对差值再用 $|f_n-f|\le2g$, 得 $L^1$ 收敛.
\end{proof}

\begin{theorem}[可积级数]\label{ra:c4:series}
若 $f_k\in L^1(\mu)$ 且 $\sum_k\int_X|f_k|\dd\mu<\infty$, 则 $\sum_kf_k$ 几乎处处绝对收敛, 和函数可积, 且
\[
 \int_X\Bigl(\sum_{k\ge1}f_k\Bigr)\dd\mu
   =\sum_{k\ge1}\int_Xf_k\dd\mu.
\]
\end{theorem}
\begin{proof}
令 $G=\sum_k|f_k|$. 单调收敛给出 $\int G=\sum_k\int|f_k|<\infty$, 所以 $G<\infty$ 几乎处处. 在这些点上原级数绝对收敛. 部分和 $F_N$ 满足 $|F_N|\le G$, 因而控制收敛给出积分交换. 还可用 $\int|\sum_kf_k-F_N|\le\sum_{k>N}\int|f_k|$ 直接看到 $L^1$ 收敛.
\end{proof}

\begin{proposition}[积分的可数可加性]\label{ra:c4:integraladditivity}
设 $E_k$ 为两两不交的可测集, $E=\bigcup_kE_k$. 若 $f\ge0$ 可测, 或 $\int_E|f|\dd\mu<\infty$, 则 $\int_Ef\dd\mu=\sum_k\int_{E_k}f\dd\mu$.
\end{proposition}
\begin{proof}
非负情形使用 $f\ind_E=\sum_kf\ind_{E_k}$ 和单调收敛. 可积情形先对 $|f|$ 使用非负结论, 得 $\sum_k\int|f\ind_{E_k}|=\int_E|f|<\infty$, 然后应用可积级数定理. 原第 12 讲陈述漏写两两不交; 若 $E_1=E_2$ 为正测度集且 $f=1$, 其公式便会重复计算.
\end{proof}

\originalsection{Egoroff 定理}

\begin{theorem}[Egoroff]\label{ra:c4:egoroff}
在任意测度空间上, 若 $\mu(X)<\infty$, 可测函数 $f_n:X\to\C$ 满足 $f_n\to f$ 几乎处处, 则对每个 $\eps>0$, 存在可测集 $E\subset X$, 使 $\mu(X\setminus E)<\eps$, 且 $f_n\to f$ 在 $E$ 上一致收敛.
\end{theorem}
\begin{proof}
先删去零测的不收敛集 $N$, 令 $Y=X\setminus N$. 保留原稿的一致 Cauchy 构造, 令
\[
 S(m,k)=Y\cap\bigcap_{i,j\ge m}\{|f_i-f_j|<1/k\}.
\]
对固定 $k$, $S(m,k)$ 随 $m$ 递增, 并覆盖 $Y$, 因为每个点上的收敛数列是 Cauchy 列. 从下连续性及 $\mu(Y)<\infty$ 给出 $\mu(Y\setminus S(m,k))\to0$. 选 $m_k$ 使这个测度小于 $\eps2^{-k}$, 并令 $E=\bigcap_kS(m_k,k)$. 则 $\mu(X\setminus E)<\eps$.

给定 $\delta>0$, 选 $k$ 使 $1/k<\delta$. 对所有 $x\in E$、$i,j\ge m_k$, 有 $|f_i(x)-f_j(x)|<1/k$. 固定 $i$ 并令 $j\to\infty$, 得 $|f_i(x)-f(x)|\le1/k<\delta$, 且这个上界对全部 $x\in E$ 一致. 因而确实得到向 $f$ 的一致收敛, 而不只是统一 Cauchy 条件.
\end{proof}

\begin{example}\label{ra:c4:infinitemeasure}
有限测度假设不可删去. 在 $\R$ 上令 $f_n=\ind_{[n,n+1)}$, 则 $f_n\to0$ 处处, 但不能删去测度小于 $1$ 的集合而得到一致收敛.
\end{example}
\begin{solution}
每个点至多属于一个区间, 所以点态极限为 $0$. 若在 $E$ 上一致趋于 $0$, 对误差 $1/2$, 必须有一个 $N$ 使 $f_n=0$ 于 $E$ 对所有 $n\ge N$ 成立. 因而 $\bigcup_{n\ge N}[n,n+1)\subset\R\setminus E$, 后者测度为无限大. 这直接核对了 Egoroff 结论的失败.
\end{solution}

\originalsection{依测度收敛与子列}

\begin{definition}\label{ra:c4:measureconv}
设 $f_n$ 和 $f$ 都是可测函数. 称 $f_n$ 依测度收敛到 $f$, 若对每个 $\delta>0$, 都有 $\mu(\{|f_n-f|>\delta\})\to0$. 等价地, 对任意 $\delta,\eta>0$, 存在 $N$ 使 $n\ge N$ 时该集合测度小于 $\eta$.
\end{definition}

原稿使用同一个 $\eps$ 同时控制函数值误差和集合测度. 它与上述两个参数版本等价: 取 $\eps<\min(\delta,\eta)$ 即可. 两参数写法在证明中更容易看清量词.

\begin{proposition}\label{ra:c4:l1measure}
若 $\int|f_n-f|\dd\mu\to0$, 则 $f_n\to f$ 依测度收敛. 若 $\mu(X)<\infty$ 且 $f_n\to f$ 几乎处处, 也有依测度收敛.
\end{proposition}
\begin{proof}
对 $\delta>0$, 积分单调性给出 $\delta\mu(\{|f_n-f|>\delta\})\le\int|f_n-f|\dd\mu$, 即第一结论. 第二结论中, 给定 $\delta,\eta>0$, 用 Egoroff 定理找 $E$ 使 $\mu(X\setminus E)<\eta$, 且在 $E$ 上一致收敛. 当 $n$ 足够大时, $|f_n-f|\le\delta$ 于 $E$, 所以大偏差集合包含在 $X\setminus E$ 内.
\end{proof}

\begin{example}\label{ra:c4:typewriter}
依测度收敛未必蕴含整个函数列几乎处处收敛.
\end{example}
\begin{solution}
\textbf{编者补全原稿提到而未写出的例子.} 在 $[0,1)$ 上, 按 $m=0,1,2,\ldots$ 的顺序, 将该层全部函数
\[
 f_{m,j}=\ind_{[j2^{-m},(j+1)2^{-m})},\qquad 0\le j<2^m,
\]
依次排成一个数列. 每层函数的积分是 $2^{-m}$, 所以数列依测度趋于 $0$. 每个点在每层恰使一个函数取 $1$, 因而无限次取 $1$; 对每层 $m\ge1$, 也有函数在该点取 $0$, 因而无限次取 $0$. 所以任一点都没有数列极限. 在端点 $1$ 可另取所有函数值为 $0$, 不影响测度结论.
\end{solution}

\begin{theorem}\label{ra:c4:subsequence}
在任意测度空间上, 若 $f_n\to f$ 依测度收敛, 则存在严格递增的指标列 $n_k$, 使 $f_{n_k}\to f$ 几乎处处.
\end{theorem}
\begin{proof}
递归选择 $n_k>n_{k-1}$, 使 $\mu(E_k)<2^{-k}$, 其中 $E_k=\{|f_{n_k}-f|>2^{-k}\}$. 依测度收敛保证可以在任意先前指标之后作此选择. 由可求和例外集引理, 几乎每个 $x$ 只落在有限多个 $E_k$ 内; 后来有 $|f_{n_k}(x)-f(x)|\le2^{-k}$, 因而收敛. 严格递增条件补正了原稿直接取 $N(k)$ 而未检查子列指标的问题.
\end{proof}

\begin{lemma}\label{ra:c4:subsubsequence}
复数列 $a_n$ 收敛到 $a$, 当且仅当它的每个子列都有一个进一步子列收敛到 $a$.
\end{lemma}
\begin{proof}
必要性直接成立. 若原数列不收敛到 $a$, 则有 $\delta>0$ 和无限多个指标使 $|a_n-a|\ge\delta$. 从这些指标选出的子列没有任何进一步子列能收敛到 $a$, 与假设矛盾.
\end{proof}

\begin{theorem}[依测度的控制收敛]\label{ra:c4:measuredct}
设 $f_n\to f$ 依测度收敛, 且 $|f_n|\le g$ 几乎处处, 其中 $g\in L^1(\mu)$ 非负. 则 $f\in L^1(\mu)$, $\int|f_n-f|\dd\mu\to0$, 从而 $\int f_n\dd\mu\to\int f\dd\mu$. 不要求 $\mu(X)<\infty$.
\end{theorem}
\begin{proof}
先选一个几乎处处收敛子列, 在其共同的零测例外集之外令指标趋于无限, 得 $|f|\le g$, 所以 $f$ 可积. 任取原列的一个子列, 它仍依测度收敛到 $f$, 所以有进一步子列几乎处处收敛到 $f$. 对这个进一步子列, $|f_{n_k}-f|\le2g$, 几乎处处控制收敛给出 $\int|f_{n_k}-f|\to0$. 应用引理 \ref{ra:c4:subsubsequence} 于非负实数列 $a_n=\int|f_n-f|$, 得整个数列趋于 $0$. 最后用 $|\int f_n-\int f|\le\int|f_n-f|$.
\end{proof}

\begin{remark}
原稿提到 ``收敛定理对依测度收敛仍成立''. 上一定理是控制收敛的准确版本. 对单调收敛, 若 $0\le f_n\uparrow h$ 且另已知 $f_n\to f$ 依测度, 子列定理使 $h=f$ 几乎处处, 再由前章单调收敛定理得 $\int f_n\uparrow\int f$. 单调性仍是必要的原假设, 不能只凭依测度收敛交换积分.
\end{remark}

\originalsection{练习与解答}

\begin{exercise}
设 $\mu(X)<\infty$, $|f_n|\le M$, 且 $f_n\to f$ 依测度收敛. 证明对每个 $1\le p<\infty$, $\int|f_n-f|^p\dd\mu\to0$.
\end{exercise}
\begin{solution}
子列定理给出 $|f|\le M$ 几乎处处. 对 $\delta>0$, 分大偏差集合和其补集, 得
\[
 \int|f_n-f|^p\dd\mu
 \le\delta^p\mu(X)+(2M)^p\mu(\{|f_n-f|>\delta\}).
\]
固定 $\delta$, 令 $n\to\infty$, 右边上极限不超过 $\delta^p\mu(X)$; 再令 $\delta\downarrow0$. 这也直接显示有限测度和统一有界性各自的用途.
\end{solution}

\begin{exercise}
在 $[0,1]$ 上令 $f_n=n\ind_{(0,1/n)}$. 比较它的几乎处处收敛、依测度收敛和积分收敛.
\end{exercise}
\begin{solution}
每个 $x>0$ 最终不在 $(0,1/n)$, 而 $x=0$ 的函数值始终为 $0$, 所以处处趋于 $0$. 对任意固定 $\delta>0$, 当 $n>\delta$ 时, 大偏差集合就是 $(0,1/n)$, 测度为 $1/n\to0$. 但 $\int_0^1f_n\dd x=1$, 并不趋于极限函数的积分 $0$. 若存在同一个可积控制函数 $g\ge f_n$ 几乎处处, 依测度控制收敛定理将迫使积分趋于 $0$, 因而这样的控制函数不存在.
\end{solution}

\begin{exercise}
设 $f\in L^1(\R^n)$, 证明对每个 $\eps>0$, 存在 $g\in C_c(\R^n)$, 使 $\int|f-g|\dd\lambda<\eps$.
\end{exercise}
\begin{solution}
\textbf{编者补充, 为后章 $L^p$ 逼近作衔接.} 先选有界盒子 $Q$ 及 $M<\infty$, 使 $h=f\ind_{Q\cap\{|f|\le M\}}$ 满足 $\int|f-h|<\eps/2$. 这是由 $|f|$ 可积以及控制收敛得到的空间截断与函数值截断. 在一个较大有界开盒子 $V\supset Q$ 内选共同开邻域 $W$, 使 $Q\subset W\subset\overline W\subset V$; 使用 Lusin 证明并将各层的开邻域限制在 $W$ 内, 得到 $g\in C_c(V)$, $\norm g_\infty\le M$, 且 $\lambda(\{h\ne g\})<\eps/(4M)$, 当 $M=0$ 时直接取 $g=0$. 因为不相等处 $|h-g|\le2M$, 所以 $\int|h-g|<\eps/2$. 三角不等式给出结论. 在一般局部紧空间上, 选择相对紧的共同邻域也有有限测度, 同样的论证适用.
\end{solution}

\originalchapter{\texorpdfstring{$L^p$}{Lp}空间}
\label{ra:c5:chapter}
\sourcelecture{14--17}

积分给出了函数的大小, 但不同问题需要不同的大小尺度. $\int\abs f\dd\mu$控制总量, $\int\abs f^2\dd\mu$控制平方平均, 本质上确界控制几乎处处的最大幅度. 本章把这些尺度组成函数空间, 证明其中的柯西列确实有极限, 再讨论连续函数逼近和不同指数之间的关系. 只涉及这些函数空间本身; 一般的线性算子理论留给泛函分析.

\originalsection{凸性与积分不等式}

\begin{definition}[凸函数]
设 $-\infty\le a<b\le\infty$. 函数 $\varphi:(a,b)\to\R$称为凸函数, 如果对任意 $x,y\in(a,b)$及 $0\le t\le1$, 都有 $\varphi((1-t)x+ty)\le(1-t)\varphi(x)+t\varphi(y)$.
\end{definition}

图形上的含义是: 两点之间的弦位于图形上方. 为了把这一直观用于积分, 我们需要将它改写为割线斜率的比较.

\begin{lemma}\label{ra:c5:convex-slopes}
函数 $\varphi$凸, 当且仅当对任意 $a<s<t<u<b$, 有
\[
 \frac{\varphi(t)-\varphi(s)}{t-s}
 \le \frac{\varphi(u)-\varphi(s)}{u-s}
 \le \frac{\varphi(u)-\varphi(t)}{u-t}.
\]
其中仅要求最左边的斜率不超过最右边的斜率, 也得到等价条件.
\end{lemma}
\begin{proof}
把 $t$写成 $\frac{u-t}{u-s}s+\frac{t-s}{u-s}u$, 凸性给出 $\varphi(t)\le\frac{u-t}{u-s}\varphi(s)+\frac{t-s}{u-s}\varphi(u)$. 移项并分别除以 $t-s$、$u-t$, 就得到两项比较. 反过来, 最左与最右斜率的比较乘以正数 $(t-s)(u-t)$, 正好得到上述凸性不等式. 任意两点间的严格内点都可这样表示; 端点情形是等式.
\end{proof}

\begin{proposition}\label{ra:c5:convex-continuity}
凸函数在开区间 $(a,b)$上连续. 更确切地说, 它在每个紧子区间上满足一个 Lipschitz 估计. 若 $\varphi$可微, 则凸性等价于 $\varphi'$单调不减; 若 $\varphi\in C^2((a,b))$, 则又等价于 $\varphi''\ge0$.
\end{proposition}
\begin{proof}
固定 $a<s<c<d<u<b$. 对 $c\le x<y\le d$, 割线斜率比较给出
\[
 \frac{\varphi(c)-\varphi(s)}{c-s}
 \le\frac{\varphi(y)-\varphi(x)}{y-x}
 \le\frac{\varphi(u)-\varphi(d)}{u-d}.
\]
例如左边可先与 $s,x$之间的斜率比较, 再与 $x,y$之间的斜率比较; 右边同理. 两端都是有限常数, 于是存在 $M<\infty$使 $\abs{\varphi(y)-\varphi(x)}\le M\abs{y-x}$. 这证明局部 Lipschitz 性及连续性.

若函数凸且可微, 令割线的端点分别趋近 $s<t$, 得 $\varphi'(s)\le\frac{\varphi(t)-\varphi(s)}{t-s}\le\varphi'(t)$. 反过来, 若导数单调不减, 对 $s<t<u$分别在 $[s,t]$和 $[t,u]$应用中值定理, 两条割线的斜率等于 $\varphi'(\xi)$和 $\varphi'(\eta)$, 其中 $\xi<t<\eta$, 故满足引理的条件. 最后, 对可微的 $\varphi'$, 单调不减等价于其导数 $\varphi''$非负, 同样由差商及中值定理得到.
\end{proof}

\Needspace{7\baselineskip}
\begin{theorem}[Jensen不等式]\label{ra:c5:jensen}
设 $(X,\mathcal M,\mu)$为概率测度空间, 即 $\mu(X)=1$. 设实值函数 $f\in L^1(\mu)$满足 $a<f<b$几乎处处, $\varphi:(a,b)\to\R$凸. 则 $t=\int_X f\dd\mu$属于 $(a,b)$, $\varphi\circ f$的积分有定义, 且
\[
 \varphi\left(\int_X f\dd\mu\right)
 \le \int_X\varphi(f)\dd\mu.
\]
右边允许等于 $+\infty$, 但不会出现正负部分均为无穷的情形.
\end{theorem}
\begin{proof}
若 $a$有限, 非负函数 $f-a$几乎处处严格为正, 因而其积分严格为正, 所以 $t>a$. 这里用了前章的结论: 非负函数积分为零, 当且仅当它几乎处处为零. 若 $a=-\infty$, 则 $t>a$由 $f\in L^1$直接成立. 同理 $t<b$.

现在在图形的 $(t,\varphi(t))$处寻找一条位于图形下方的直线. 设
\[
 B=\sup_{a<s<t}\frac{\varphi(t)-\varphi(s)}{t-s}.
\]
取任意 $u\in(t,b)$, 引理给出 $B\le(\varphi(u)-\varphi(t))/(u-t)$; 又任取一个 $s<t$给出有限下界, 所以 $B$是实数. 对 $s<t$由上确界定义, 对 $s>t$由割线比较, 都有 $\varphi(s)\ge\varphi(t)+B(s-t)$.

凸函数连续, 故 $\varphi(f)$可测. 由上述下界, $\varphi(f)$的负部分被可积函数 $\abs{\varphi(t)}+\abs B(\abs f+\abs t)$控制. 因而可对非负函数 $\varphi(f)-\varphi(t)-B(f-t)$积分, 得
\[
 \int_X\varphi(f)\dd\mu
 \ge\varphi(t)\mu(X)+B\left(\int_X f\dd\mu-t\mu(X)\right)
 =\varphi(t).
\]
若 $f$仅几乎处处取值于区间, 可在相应零测集上把它改为任意固定内点, 不改变各积分.
\end{proof}

\begin{remark}
以上连续性证明及 Jensen 不等式中积分存在性的核查为编者补充. 不能直接对可能不可积的 $\varphi(f)$分拆三个积分再相减; 支撑直线先保证负部分可积, 才使这个操作合法.
\end{remark}

\begin{example}\label{ra:c5:amgm}
取凸函数 $\varphi(s)=e^s$, 则 $\exp(\int f\dd\mu)\le\int e^f\dd\mu$. 特别地, 若 $y_1,\ldots,y_n>0$, $\alpha_i\ge0$且 $\sum_i\alpha_i=1$, 在有限概率空间上令 $f(i)=\log y_i$, 得 $\prod_{i=1}^n y_i^{\alpha_i}\le\sum_{i=1}^n\alpha_i y_i$.
等权情形给出 $(y_1\cdots y_n)^{1/n}\le(y_1+\cdots+y_n)/n$. 若某个 $y_i=0$, 对正权重项取 $y_i+\delta$再令 $\delta\downarrow0$, 仍得到相应不等式; 零权重项可以先删除.
\end{example}

\begin{lemma}[Young不等式]\label{ra:c5:young}
设 $1<p<\infty$, $q=p/(p-1)$, 即 $1/p+1/q=1$. 对 $A,B\ge0$, 有 $AB\le A^p/p+B^q/q$.
\end{lemma}
\begin{proof}
若 $A=0$或 $B=0$, 结论直接成立. 否则写 $A=e^{s/p}$、$B=e^{t/q}$, 在例题的加权算术几何平均不等式中取两个数 $e^s,e^t$及权重 $1/p,1/q$, 得 $e^{s/p+t/q}\le e^s/p+e^t/q$.
\end{proof}

\begin{theorem}[Hölder不等式]\label{ra:c5:holder}
设 $1<p<\infty$, $1/p+1/q=1$. 若复值可测函数 $f,g$满足 $\int_X\abs f^p\dd\mu<\infty$及 $\int_X\abs g^q\dd\mu<\infty$, 则 $fg\in L^1(\mu)$, 且
\[
 \int_X\abs{fg}\dd\mu
 \le\left(\int_X\abs f^p\dd\mu\right)^{1/p}
      \left(\int_X\abs g^q\dd\mu\right)^{1/q}.
\]
\end{theorem}
\begin{proof}
记右边两个因子为 $F,G$. 若其中一个为零, 对应的函数几乎处处为零, 所以乘积也几乎处处为零. 若 $F,G>0$, 令 $u=\abs f/F$、$v=\abs g/G$, 则 $\int u^p\dd\mu=\int v^q\dd\mu=1$. Young 不等式给出 $uv\le u^p/p+v^q/q$. 积分后右边为 $1/p+1/q=1$, 乘回 $FG$便得结论.
\end{proof}

\begin{theorem}[Minkowski不等式]\label{ra:c5:minkowski}
设 $1<p<\infty$, $f,g$为复值可测函数且 $\int\abs f^p\dd\mu$、$\int\abs g^p\dd\mu$有限. 则
\[
 \left(\int_X\abs{f+g}^p\dd\mu\right)^{1/p}
 \le\left(\int_X\abs f^p\dd\mu\right)^{1/p}
    +\left(\int_X\abs g^p\dd\mu\right)^{1/p}.
\]
\end{theorem}
\begin{proof}
先用 $s\mapsto s^p$在 $[0,\infty)$上的凸性得到 $(A+B)^p\le2^{p-1}(A^p+B^p)$, 因而 $h=\abs f+\abs g$的 $p$次方可积. 这一步保证后面除去的量是有限的. 记 $H=(\int h^p\dd\mu)^{1/p}$, 并取共轭指数 $q=p/(p-1)$. 若 $H=0$, 结论成立. 否则由 Hölder 不等式,
\begin{align*}
 H^p&=\int_X\abs f h^{p-1}\dd\mu+\int_X\abs g h^{p-1}\dd\mu\\
 &\le\left[\left(\int_X\abs f^p\dd\mu\right)^{1/p}
            +\left(\int_X\abs g^p\dd\mu\right)^{1/p}\right]
       \left(\int_X h^{(p-1)q}\dd\mu\right)^{1/q}\\
 &=\left[\left(\int_X\abs f^p\dd\mu\right)^{1/p}
            +\left(\int_X\abs g^p\dd\mu\right)^{1/p}\right]H^{p-1}.
\end{align*}
除以 $H^{p-1}$, 再用 $\abs{f+g}\le h$, 即得所求.
\end{proof}

\begin{remark}
原讲义还允许非负可测函数取值 $+\infty$, 并把上述两个不等式解释为扩展非负积分的不等式. 这可由已证有限范数情形补齐: 在 Hölder 不等式中, 若一个积分为零, 对应函数几乎处处为零, 按 $0\cdot\infty=0$的约定两边都为零; 若没有零因子而有一个因子无穷, 右边为无穷, 结论直接成立. 在 Minkowski 不等式中, 若右边有一个无穷项, 结论也直接成立. 因而真正需要证明的正是两个范数都有限的情形.
\end{remark}

\originalsection{函数空间与范数}

\begin{definition}\label{ra:c5:lp-definition}
设 $(X,\mathcal M,\mu)$为任意测度空间. 对 $1\le p<\infty$, 定义 $\norm f_p=(\int_X\abs f^p\dd\mu)^{1/p}$, 并以 $L^p(\mu)$表示所有可测复值函数中满足 $\norm f_p<\infty$的函数, 按几乎处处相等取等价类后形成的空间.

称 $f$本质有界, 如果存在 $M<\infty$使 $\abs f\le M$几乎处处. 所有本质有界可测函数的等价类组成 $L^\infty(\mu)$, 其本质上确界范数为 $\norm f_\infty=\inf\{M\ge0:\abs f\le M\text{ 几乎处处}\}$.
\end{definition}

以后仍用 $f$表示等价类或其代表元, 但凡涉及单点值, 必须说明选择了代表元. 例如某个零测集上的极大值不改变 $L^\infty$范数. 在实值函数上也可作同样的定义, 得到实向量空间.

\begin{proposition}\label{ra:c5:esssup-bound}
若 $f\in L^\infty(\mu)$, 则 $\abs{f(x)}\le\norm f_\infty$几乎处处.
\end{proposition}
\begin{proof}
由下确界定义, 对每个 $k$可取 $M_k<\norm f_\infty+1/k$使 $\abs f\le M_k$在某个零测集 $N_k$外成立. 可数并 $N=\bigcup_kN_k$仍为零测集. 对 $x\notin N$, 所有这些不等式同时成立; 令 $k\to\infty$便得结论.
\end{proof}

\begin{definition}[范数与完备性]
复向量空间 $V$上的范数是函数 $\norm{\cdot}:V\to[0,\infty)$, 满足 $\norm v=0$当且仅当 $v=0$, $\norm{\alpha v}=\abs\alpha\norm v$, 以及 $\norm{v+w}\le\norm v+\norm w$. 范数诱导距离 $d(v,w)=\norm{v-w}$.

度量空间中的序列 $(v_n)$称为柯西列, 如果对每个 $\eps>0$, 存在 $N$使 $m,n\ge N$时 $d(v_m,v_n)<\eps$. 若每个柯西列都收敛于空间中的一个点, 称该空间完备. 在范数诱导的距离下完备的赋范空间称为 Banach 空间.
\end{definition}

上述距离非负、对称, 且 $d(v,w)=0$当且仅当 $v=w$; 三角不等式直接来自范数的三角不等式. 收敛列必为柯西列, 因为 $d(v_m,v_n)\le d(v_m,v)+d(v,v_n)$. 反向推论需要完备性. 例如带通常距离的 $\R\setminus\{0\}$中, $(1/n)$是柯西列, 但它的极限 $0$不在该空间中. 这里 $\R\setminus\{0\}$是度量空间, 不是实向量空间.

\begin{proposition}\label{ra:c5:lp-norm}
对 $1\le p\le\infty$, $L^p(\mu)$是向量空间, $\norm{\cdot}_p$是其上的范数. 若 $1/p+1/q=1$, 约定 $1/\infty=0$, 则对 $f\in L^p(\mu)$及 $g\in L^q(\mu)$, 有 $fg\in L^1(\mu)$及 $\norm{fg}_1\le\norm f_p\norm g_q$.
\end{proposition}
\begin{proof}
积分或本质上确界不随零测集上的修改而改变, 所以范数在等价类上有定义. 当 $p<\infty$时, 非负积分为零的判别说明 $\norm f_p=0$等价于 $f=0$几乎处处; 当 $p=\infty$时, 命题~\ref{ra:c5:esssup-bound}给出同一结论. 齐次性由定义得到.

对 $1<p<\infty$, 向量加法的封闭性及三角不等式已经由 Minkowski 不等式证明. 对 $p=1$, 积分 $\abs{f+g}\le\abs f+\abs g$即可; 对 $p=\infty$, 在两个本质上确界估计共同成立的零测集外有 $\abs{f+g}\le\norm f_\infty+\norm g_\infty$. 这也证明了向量空间的封闭性.

Hölder 不等式已处理 $1<p<\infty$. 在端点 $p=1,q=\infty$, 有 $\abs{fg}\le\norm g_\infty\abs f$几乎处处, 积分即可. 另一个端点交换 $f,g$.
\end{proof}

\begin{remark}
对 $0<p<1$, 同一积分表达式仍可定义, 但一般不满足三角不等式, 因而不是这里所说的范数. 例如在二点计数测度空间上, $f=(1,0)$、$g=(0,1)$有 $\norm f_p=\norm g_p=1$, 而 $\norm{f+g}_p=2^{1/p}>2$. 本章的赋范空间结论始终以 $p\ge1$为条件.
\end{remark}

\begin{example}
在 $\N$上取计数测度, 非负函数的积分为 $\sum_{n\ge1}f(n)$, 所以 $L^p$就是序列空间 $\ell^p$. 对 $1\le p<\infty$, 序列 $a=(a_n)$属于 $\ell^p$当且仅当 $\sum_n\abs{a_n}^p<\infty$, 范数为该和的 $1/p$次幂; $\ell^\infty$是有界序列空间, 范数为 $\sup_n\abs{a_n}$. 计数测度没有非空零测集, 此时不需要进一步区分代表元.
\end{example}

\originalsection{完备性}

柯西性质只比较序列中已有的函数. 为了构造未知的极限, 我们选一个相邻差可以求和的子列. 点态上求和需要这些差的绝对值总和有限, 而 Fatou 引理恰能从范数估计得到这个条件.

\begin{lemma}\label{ra:c5:cauchy-subsequence}
任意度量空间中的柯西列 $(v_n)$都有子列 $(v_{n_k})$, 满足 $d(v_{n_{k+1}},v_{n_k})<2^{-k}$.
\end{lemma}
\begin{proof}
令 $n_0=0$. 对每个 $k\ge1$选 $N_k$使 $m,n\ge N_k$时 $d(v_m,v_n)<2^{-k}$. 依次选 $n_k\ge\max\{N_1,\ldots,N_k,n_{k-1}+1\}$, 则 $n_k,n_{k+1}\ge N_k$, 所求估计成立.
\end{proof}

\begin{theorem}[Riesz--Fischer定理]\label{ra:c5:riesz-fischer}
对任意测度空间 $(X,\mathcal M,\mu)$及 $1\le p\le\infty$, $L^p(\mu)$都是 Banach 空间.
\end{theorem}
\begin{proof}
先设 $1\le p<\infty$, 并设 $(f_n)$在 $L^p$中为柯西列. 选择各函数的可测代表元, 再由引理选子列使 $\norm{f_{n_{k+1}}-f_{n_k}}_p<2^{-k}$. 令 $h_k=\abs{f_{n_{k+1}}-f_{n_k}}$, $G_m=\sum_{k=1}^m h_k$, $G=\sum_{k=1}^\infty h_k$. 有限次 Minkowski 不等式给出 $\norm{G_m}_p\le\sum_{k=1}^m2^{-k}<1$. 对 $G_m^p$应用 Fatou 引理, 得 $\int G^p\dd\mu\le1$, 所以 $G<\infty$几乎处处.

于是级数 $f_{n_1}(x)+\sum_{k\ge1}(f_{n_{k+1}}(x)-f_{n_k}(x))$几乎处处绝对收敛. 在它收敛的集合上将其和定义为 $f(x)$, 在其余零测集上定义 $f(x)=0$. 收敛集合可由可测的部分和刻画, 所以 $f$可测, 且 $\abs f\le\abs{f_{n_1}}+G$几乎处处, 从而 $f\in L^p$.

级数的第 $k-1$个部分和正是 $f_{n_k}$. 对尾和, 有 $\abs{f-f_{n_k}}\le\sum_{j\ge k}h_j$几乎处处. 再对有限尾和应用 Minkowski, 然后应用 Fatou, 得
\[
 \norm{f-f_{n_k}}_p
 \le\sum_{j=k}^\infty\norm{h_j}_p
 \le\sum_{j=k}^\infty2^{-j}=2^{1-k}.
\]
因此子列在 $L^p$范数中收敛到 $f$. 为了回到原序列, 给定 $\eps>0$, 先取 $N$使 $m,n\ge N$时 $\norm{f_m-f_n}_p<\eps/2$, 再取 $k$使 $n_k\ge N$且 $\norm{f_{n_k}-f}_p<\eps/2$. 对所有 $n\ge N$, 三角不等式给出 $\norm{f_n-f}_p<\eps$.

现在设 $p=\infty$. 把以下各个估计的例外零测集取可数并: $\abs{f_n(x)}\le\norm{f_n}_\infty$, 以及对所有 $m,n$成立的 $\abs{f_m(x)-f_n(x)}\le\norm{f_m-f_n}_\infty$. 记此并集为 $N$. 在 $X\setminus N$上, 柯西条件同时给出逐点柯西性和一致的柯西估计. 由 $\C$完备, $f_n(x)$收敛到某个 $f(x)$. 将 $f$在 $N$上定义为零, 得到可测函数.

范数中的柯西列有界, 所以存在 $M$使所有 $\norm{f_n}_\infty\le M$. 在 $N$外取极限得 $\abs f\le M$, 故 $f\in L^\infty$. 又给定 $\eps>0$, 柯西条件使 $m,n\ge N_0$时 $\abs{f_m-f_n}\le\eps$在 $N$外成立. 固定 $n\ge N_0$并令 $m\to\infty$, 得 $\abs{f-f_n}\le\eps$在 $N$外成立, 即 $\norm{f-f_n}_\infty\le\eps$. 这完成端点情形的证明.
\end{proof}

\begin{remark}
引理的证明及尾和估计为编者补充. 此证明没有假设 $\mu(X)$有限或测度为 $\sigma$有限; 完备性对任意测度空间都成立.
\end{remark}

\begin{corollary}\label{ra:c5:lp-ae-subsequence}
若 $1\le p\le\infty$且 $\norm{f_n-f}_p\to0$, 则有一个子列几乎处处收敛到 $f$.
\end{corollary}
\begin{proof}
当 $p<\infty$时, 选 $n_k$使 $\norm{f_{n_k}-f}_p<2^{-k}$. 与前一定理中对 $G_m$的论证相同, $\sum_k\abs{f_{n_k}-f}$属于 $L^p$, 故几乎处处有限. 其一般项必趋于零, 即 $f_{n_k}\to f$几乎处处. 当 $p=\infty$时, 将所有 $\abs{f_n-f}\le\norm{f_n-f}_\infty$的例外零测集取并; 在其外, 原序列本身就一致收敛到 $f$.
\end{proof}

\begin{example}\label{ra:c5:spikes}
在 $\R$的 Lebesgue 测度下, 令 $f_k=k^2\ind_{(0,1/k)}$. 每个 $f_k$属于所有 $1\le p<\infty$的 $L^p$, 且对每个 $x\in\R$都有 $f_k(x)\to0$. 但是 $\norm{f_k}_p=(k^{2p}/k)^{1/p}=k^{2-1/p}\to\infty$. 因而几乎处处收敛不保证 $L^p$收敛.

另令 $g_k=k^{-1}\ind_{(k,2k)}$. 则 $g_k\to0$处处, $\norm{g_k}_p=k^{1/p-1}\to0$对每个 $p>1$成立, 而 $\norm{g_k}_1=1$. 同一列可在一种 $L^p$范数下收敛, 在另一种范数下不收敛.
\end{example}

\begin{example}\label{ra:c5:typewriter}
对 $m\ge0$及 $0\le j<2^m$, 在 $[0,1)$上令 $f_{2^m+j}=\ind_{[j2^{-m},(j+1)2^{-m})}$. 这按从左到右的顺序, 依次列出第 $m$层所有二进小区间的示性函数. 对有限的 $p\ge1$, 有 $\norm{f_{2^m+j}}_p=2^{-m/p}$, 因而整列在 $L^p$中趋于零.

但是对任意固定 $x\in[0,1)$, 每层恰有一个函数在 $x$处等于 $1$, 而从 $m\ge1$起每层也有函数在 $x$处等于 $0$. 所以整列在任何 $x\in[0,1)$处都不收敛. 这补全了原讲义留作作业的反例, 也解释了上面的推论为什么只保证子列.
\end{example}

\originalsection{连续与平滑逼近}

\begin{theorem}\label{ra:c5:simple-density}
设 $1\le p<\infty$. 所有满足 $\mu(\{s\ne0\})<\infty$的可测简单函数组成的集合 $S$, 在 $L^p(\mu)$中稠密.
\end{theorem}
\begin{proof}
若 $s\in S$, 则 $\int\abs s^p\dd\mu\le\norm s_\infty^p\mu(\{s\ne0\})<\infty$, 所以 $S\subset L^p$. 先设 $f\ge0$且 $f\in L^p$. 由前章的简单函数逼近, 可取 $0\le s_n\uparrow f$. 因为 $s_n^p\le f^p$, 每个 $s_n\in L^p$. 若 $s_n$不恒为零, 其有限个非零取值中存在最小值 $c_n>0$, 因而 $c_n^p\mu(\{s_n\ne0\})\le\int s_n^p\dd\mu<\infty$. 这就核实了 $s_n\in S$.

又有 $\abs{f-s_n}^p\le f^p$并逐点趋于零, 控制收敛定理给出 $\norm{f-s_n}_p\to0$. 对实值 $f$分别逼近 $f^+,f^-$; 对复值 $f$分别逼近实部和虚部的正负部分. 这些部分的绝对值都不超过 $\abs f$, 故都属于 $L^p$. 合并四个简单函数, 支撑仍是有限测度集合的有限并, 再用 Minkowski 不等式控制总误差.
\end{proof}

下面的连续逼近需要测度与拓扑相容. 单有局部紧 Hausdorff 拓扑, 并不足以把任意测度空间的简单函数变成连续函数.

\begin{theorem}\label{ra:c5:cc-density}
设 $X$为局部紧 Hausdorff 空间, $\mu$为正 Radon 测度, 即在紧集上有限并满足前章的内、外正则性条件; 也允许使用它的完备化. 则对 $1\le p<\infty$, $C_c(X)$在 $L^p(\mu)$中稠密.
\end{theorem}
\begin{proof}
由前一定理, 只需逼近有限测度集合的示性函数. 设 $E$可测且 $\mu(E)<\infty$, $\delta>0$. 测度正则性给出紧集 $K\subset E$和开集 $U\supset E$, 满足 $\mu(U\setminus K)<\delta$. 前章的紧集连续截断结论给出 $h\in C_c(X)$, 满足 $0\le h\le1$, 在 $K$上等于 $1$, 且 $\supp h\subset U$. 函数 $h-\ind_E$仅可能在 $U\setminus K$上非零, 幅度不超过 $1$, 因而 $\norm{h-\ind_E}_p\le\delta^{1/p}$. 在完备化情形, 先把 $E$换成与它相差零测集的 Borel 集, 再作同一构造, 其 $L^p$等价类不变.

对简单函数 $s=\sum_{i=1}^m a_i\ind_{E_i}\in S$, 每个非零系数对应的 $E_i$均有有限测度. 对各 $\ind_{E_i}$作上述逼近并取线性组合 $g\in C_c(X)$, 由 $\norm{g-s}_p\le\sum_i\abs{a_i}\norm{h_i-\ind_{E_i}}_p$, 可使误差任意小. 最后先逼近 $f$为 $s$, 再逼近 $s$为 $g$.
\end{proof}

\begin{remark}
原讲义在此调用 Lusin 定理: 对 $s\in S$可取 $g\in C_c(X)$, 满足 $\norm g_\infty\le\norm s_\infty$及 $\mu(\{g\ne s\})<\delta$, 从而 $\norm{g-s}_p\le2\norm s_\infty\delta^{1/p}$. 这条路线同样有效. 上面的正则性证明为编者补充, 并把该结论所需的测度条件写出.
\end{remark}

在 $\R^n$的 Lebesgue 测度下, 连续函数若几乎处处为零就处处为零: 否则某一点附近的绝对值有正下界, 而该邻域有正测度. 因而 $C_c(\R^n)$可直接视为 $L^p(\R^n)$的一个子空间. 由稠密性及 Riesz--Fischer 定理, $L^p(\R^n)$正是 $C_c(\R^n)$在 $L^p$范数下的完备化. 若一般 Radon 测度的支撑不是整个 $X$, 则应先按零范数取商, 才有这一表述.

\begin{definition}
连续函数 $f:X\to\C$称为在无穷远处消失, 如果对每个 $\eps>0$, 存在紧集 $K\subset X$, 使 $x\notin K$时 $\abs{f(x)}<\eps$. 这些函数组成 $C_0(X)$, 赋予通常的上确界范数 $\norm f_{\mathrm{sup}}=\sup_{x\in X}\abs{f(x)}$.
\end{definition}

这样的函数有界: 在某个紧集外幅度小于 $1$, 在紧集上又由连续性而有界. 这里的上确界看每一个点, 与依赖测度的本质上确界应当区分.

\begin{theorem}\label{ra:c5:c0-completion}
若 $X$为局部紧 Hausdorff 空间, 则 $C_c(X)$在 $C_0(X)$中按上确界范数稠密, 且 $C_0(X)$在此范数下完备. 因此 $C_c(X)$的上确界范数完备化是 $C_0(X)$.
\end{theorem}
\begin{proof}
给定 $f\in C_0(X)$及 $\eps>0$, 取紧集 $K$使 $\abs f<\eps$在 $K$外成立. 用连续截断结论取 $\chi\in C_c(X)$, 满足 $0\le\chi\le1$及 $\chi=1$在 $K$上成立. 令 $h=f\chi$. 则 $h\in C_c(X)$, 在 $K$上 $h=f$, 在 $K$外 $\abs{f-h}=\abs f\abs{1-\chi}<\eps$, 故 $\norm{f-h}_{\mathrm{sup}}\le\eps$.

若 $(f_n)$是 $C_0(X)$中的柯西列, 则对每个 $x$, $(f_n(x))$是复数柯西列, 因而有极限 $f(x)$. 在一致的柯西估计中令一个指标趋于无穷, 得 $f_n\to f$一致. 为说明 $f$连续, 固定 $x$及 $\eps>0$, 选 $n$使一致误差小于 $\eps/3$, 再在 $x$的某个邻域上用 $f_n$的连续性保证 $\abs{f_n(y)-f_n(x)}<\eps/3$; 三角不等式给出 $\abs{f(y)-f(x)}<\eps$.

再选 $n$使 $\norm{f-f_n}_{\mathrm{sup}}<\eps/2$, 并选紧集 $K$使 $\abs{f_n}<\eps/2$在 $K$外成立. 则 $\abs f<\eps$在 $K$外成立, 所以 $f\in C_0(X)$. 这证明完备性.
\end{proof}

\begin{example}\label{ra:c5:linfty-not-density}
$C_c(\R)$在 $L^\infty(\R)$中不稠密. 取 $f=\ind_{(0,\infty)}$. 对任何连续函数 $g$, 都有 $\norm{f-g}_\infty\ge1/2$. 否则可取 $c<1/2$使 $\abs{f-g}\le c$几乎处处. 连续性保证 $\abs g\le c$在整个 $(-\infty,0)$上成立, 因为一个严格违反此估计的点会带来一个正测度邻域; 同理 $\abs{1-g}\le c$在整个 $(0,\infty)$上成立. 两侧趋近零得到 $\abs{g(0)}\le c$及 $\abs{1-g(0)}\le c$, 与 $1\le2c$矛盾. 所以 $L^\infty$的完备性不能给出 $C_c$在它里面稠密.
\end{example}

\begin{proposition}[平滑逼近, 编者补充]\label{ra:c5:smooth-density}
对 $1\le p<\infty$, $C_c^\infty(\R^n)$在 $L^p(\R^n)$中稠密. 更强地, 每个 $h\in C_c(\R^n)$都可以由支撑包含在同一个紧集内的光滑函数一致逼近.
\end{proposition}
\begin{proof}
先说明所需的光滑截断. 函数 $\rho(t)=e^{-1/t}$在 $t>0$时定义, 在 $t\le0$时定义为零. 对 $t>0$, 其任意阶导数都是 $e^{-1/t}$与 $1/t$的多项式之积, 当 $t\downarrow0$时都趋于零, 所以 $\rho\in C^\infty(\R)$. 函数 $\theta(t)=\rho(t)/(\rho(t)+\rho(1-t))$光滑, 在 $t\le0$时为零, 在 $t\ge1$时为一. 由此可构造在一个闭球上等于一、在稍大球外为零且取值于 $[0,1]$的光滑函数: 将 $\theta$与适当的 $R^2-\abs{x-a}^2$的仿射函数复合即可.

因此对紧集 $K\subset U$及开集 $U\subset\R^n$, 可在每个 $x\in K$周围取两个同心球, 使大球的闭包包含在 $U$中, 并为它们取上述球形截断. 有限个小球覆盖 $K$, 对应截断的和 $H$在 $K$上至少为一, 且有紧支撑包含在 $U$中. 再取一个光滑实函数 $\beta$, 满足 $0\le\beta\le1$, 在 $(-\infty,1/2]$上为零, 在 $[1,\infty)$上为一, 则 $\beta(H)$就是在 $K$上等于一且支撑包含在 $U$中的光滑截断.

固定 $h\in C_c(\R^n)$, 令 $K=\supp h$; 若 $h=0$无需逼近. 函数 $h$在整个 $\R^n$上一致连续, 因为它在一个包含 $K$及其邻域的大闭球上连续, 并且远离该球时为零. 给定 $\eps>0$, 选 $0<r<1$使 $\abs{x-y}<2r$时 $\abs{h(x)-h(y)}<\eps$. 用有限个球 $B(x_i,r)$覆盖 $K$, 其中 $x_i\in K$. 取 $\psi_i\in C_c^\infty$使 $0\le\psi_i\le1$, $\psi_i=1$在 $B(x_i,r)$上, 且 $\supp\psi_i\subset B(x_i,2r)$, 并令 $H=\sum_i\psi_i$. 在 $K$上 $H\ge1$.

取光滑截断 $\chi$在 $K$上等于一, 且 $\supp\chi\subset\{H>0\}$. 在 $\{H>0\}$上定义
\[
 v(x)=\chi(x)\frac{\sum_i h(x_i)\psi_i(x)}{H(x)},
\]
并在其余位置定义 $v=0$. 因为 $\supp\chi$紧且包含在开集 $\{H>0\}$中, 这个延拓是光滑的, 且有紧支撑. 当 $\psi_i(x)>0$时 $\abs{x-x_i}<2r$, 所以加权平均 $A(x)=\sum_i h(x_i)\psi_i(x)/H(x)$满足 $\abs{A(x)-h(x)}<\eps$. 若 $x\in K$, 则 $v=A$; 若 $x\notin K$, 则 $h(x)=0$且 $\abs v\le\abs A<\eps$. 故 $\norm{v-h}_{\mathrm{sup}}\le\eps$.

所有这些 $v$的支撑均包含在固定紧集 $K_2=\{x:\operatorname{dist}(x,K)\le2\}$中. 于是 $\norm{v-h}_p\le\norm{v-h}_{\mathrm{sup}}\,\mu(K_2)^{1/p}\to0$. 最后先用定理~\ref{ra:c5:cc-density}把任意 $f\in L^p(\R^n)$逼近为 $h\in C_c$, 再用这里的构造逼近 $h$, 即得结论. 这一证明只用连续性、光滑截断和有限加权平均, 不需要乘积积分定理; 后面引入卷积时会得到另一种逼近构造.
\end{proof}

\originalsection{不同指数的比较}

\begin{example}\label{ra:c5:no-inclusions}
在 $(0,\infty)$上取 Lebesgue 测度. 对 $1\le p<q\le\infty$, 选 $1/q<b<1/p$, 其中 $1/\infty=0$. 函数 $f(x)=x^{-b}\ind_{(0,1)}$满足 $\int_0^1\abs f^p\dd x=1/(1-bp)<\infty$, 但当 $q<\infty$时 $bq>1$, 所以 $\int_0^1\abs f^q\dd x=\infty$; 当 $q=\infty$时, 它在零点附近本质无界. 因而 $L^p$不包含于 $L^q$.

反过来, $g(x)=x^{-b}\ind_{(1,\infty)}$在 $q<\infty$时满足 $\int_1^\infty\abs g^q\dd x=1/(bq-1)<\infty$, 在 $q=\infty$时本质有界, 但 $bp<1$使它不属于 $L^p$. 所以 $L^q$也不包含于 $L^p$. 小区间上的奇性与无穷远处的衰减, 对指数提出相反的要求.
\end{example}

\begin{theorem}\label{ra:c5:finite-inclusion}
若 $0<\mu(X)<\infty$且 $1\le p<q\le\infty$, 则 $L^q(\mu)\subset L^p(\mu)$, 并且 $\norm f_p\le\mu(X)^{1/p-1/q}\norm f_q$.
特别地, 在概率空间上 $\norm f_p\le\norm f_q$.
\end{theorem}
\begin{proof}
先设 $q<\infty$. 对 $\abs f^p$和常函数 $1$, 用共轭指数 $q/p$与 $q/(q-p)$应用 Hölder 不等式, 得 $\int\abs f^p\dd\mu\le(\int\abs f^q\dd\mu)^{p/q}\mu(X)^{1-p/q}$. 取 $1/p$次幂即得结论. 当 $q=\infty$时, 在零测集外有 $\abs f^p\le\norm f_\infty^p$, 积分后取根即可.
\end{proof}

\begin{theorem}\label{ra:c5:sequence-inclusion}
在任意集合 $X$的计数测度上, 对 $1\le p<q\le\infty$, 有 $\ell^p(X)\subset\ell^q(X)$及 $\norm f_q\le\norm f_p$.
\end{theorem}
\begin{proof}
对任意 $x\in X$, 有 $\abs{f(x)}^p\le\sum_{y\in X}\abs{f(y)}^p$, 因而 $\norm f_\infty\le\norm f_p$. 若 $q<\infty$且 $\norm f_p>0$, 将 $f$除以此范数得到 $u$; 则 $\sum_x\abs{u(x)}^p=1$且所有 $\abs{u(x)}\le1$. 因为 $q>p$, 有 $\abs{u(x)}^q\le\abs{u(x)}^p$, 求和得 $\norm u_q\le1$. 乘回 $\norm f_p$即可. 零范数情形直接成立.

若 $X$不可数, 非负和按所有有限子集上有限和的上确界定义. 当它有限时, 每个集合 $\{x:\abs{f(x)}\ge1/m\}$必有限, 所以非零取值只能出现在一个可数集合上; 上面的求和论证因此仍然成立.
\end{proof}

有限测度空间上, 大指数控制小指数; 计数测度的序列空间上, 包含方向相反. 两种结论的证明也说明了原因: 前者可以把常函数 $1$放进 Hölder 不等式, 后者则由每个点具有固定的正测度而获得逐点控制.

\begin{definition}[局部可积性]
设 $G\subset\R^n$开, $1\le p\le\infty$. 称可测函数 $f$属于 $L^p_{\mathrm{loc}}(G)$, 如果对每个紧集 $K\subset G$, 都有 $f\ind_K\in L^p(G)$. 对有限 $p$, 这就是说 $\int_K\abs f^p\dd x<\infty$; 对 $p=\infty$, 这就是说 $f$在每个这样的 $K$上本质有界.
\end{definition}

每个紧集具有有限 Lebesgue 测度, 所以定理~\ref{ra:c5:finite-inclusion}立即给出 $L^q_{\mathrm{loc}}(G)\subset L^p_{\mathrm{loc}}(G)$, 其中 $1\le p<q\le\infty$. 例如 $f(x)=x^{-1}$在 $G=(0,\infty)$上局部有界, 但在 $(0,1)$上的积分无穷; 局部条件只考察远离 $G$边界的紧集, 不控制靠近边界或无穷远处的行为.

\begin{theorem}\label{ra:c5:interpolation}
设 $1\le p<r<q\le\infty$, $f\in L^p(\mu)\cap L^q(\mu)$. 唯一选取 $0<\theta<1$使 $1/r=\theta/p+(1-\theta)/q$, 则 $f\in L^r(\mu)$且 $\norm f_r\le\norm f_p^\theta\norm f_q^{1-\theta}$.
若 $f$不是零等价类, 则 $\log\norm f_p$作为变量 $1/p$的函数, 在它有定义的指数区间上凸.
\end{theorem}
\begin{proof}
若 $q<\infty$, 则 $r\theta/p+r(1-\theta)/q=1$. 因为这两项均为正, 指数 $p/(r\theta)$及 $q/(r(1-\theta))$大于一且共轭. 对 $\abs f^{r\theta}$和 $\abs f^{r(1-\theta)}$应用 Hölder 不等式,
\[
 \int_X\abs f^r\dd\mu
 \le\left(\int_X\abs f^p\dd\mu\right)^{r\theta/p}
     \left(\int_X\abs f^q\dd\mu\right)^{r(1-\theta)/q}.
\]
取 $1/r$次幂得到范数估计. 若 $q=\infty$, 则 $\theta=p/r$. 直接估计 $\abs f^r\le\norm f_\infty^{r-p}\abs f^p$几乎处处, 即得同样结论. 这个端点为编者补充.

只要两个指数上的范数有限, 中间指数的范数也有限, 所以 $I=\{s\in[1,\infty]:f\in L^s\}$是区间, 端点可以包含或不包含. 非零等价类的这些范数严格为正, 对范数估计取对数, 得 $\log\norm f_r\le\theta\log\norm f_p+(1-\theta)\log\norm f_q$. 这正是关于 $1/s$的凸性.
\end{proof}

\begin{theorem}\label{ra:c5:limit-infty}
设 $f\in L^s(\mu)$对某个 $1\le s<\infty$成立. 对所有其他指数允许把 $\norm f_p$定义为 $+\infty$. 则在扩展实数意义下, $\lim_{p\to\infty}\norm f_p=\norm f_\infty$.
\end{theorem}
\begin{proof}
若 $f=0$几乎处处, 结论直接成立. 设 $0<t<\norm f_\infty$, 令 $A_t=\{x:\abs{f(x)}>t\}$. 本质上确界的定义保证 $\mu(A_t)>0$, 而 $f\in L^s$给出 $\mu(A_t)\le t^{-s}\norm f_s^s<\infty$. 对每个有限 $p$, 有 $\norm f_p\ge t\mu(A_t)^{1/p}$. 令 $p\to\infty$, 得 $\liminf\norm f_p\ge t$. 若 $\norm f_\infty=\infty$, 可令 $t$任意大, 便得到极限为无穷.

若 $0<\norm f_\infty<\infty$, 对 $p>s$有 $\int\abs f^p\dd\mu\le\norm f_\infty^{p-s}\int\abs f^s\dd\mu$, 故 $\norm f_p\le\norm f_s^{s/p}\norm f_\infty^{1-s/p}$. 右边趋于 $\norm f_\infty$, 与下界合并即得结论.
\end{proof}

在上面证明中使用的估计 $\mu(\{\abs f>t\})\le t^{-p}\norm f_p^p$就是 Chebyshev 不等式. 对 $p<\infty$, 它也说明 $L^p$范数收敛蕴含依测度收敛; 这与前章的几乎处处收敛子列结论相容. 假设至少有一个有限指数使 $f\in L^s$不能随意删除: 在无限测度空间上, 常函数 $1$有 $\norm1_\infty=1$, 但所有有限指数的范数都是无穷.

\originalsection{练习}

以下练习为本整理本配套练习, 不属于原课程独立作业集的逐题翻译.

\begin{exercise}\label{ra:c5:ex-holder-equality}
设 $1<p<\infty$, $1/p+1/q=1$, $f\in L^p$, $g\in L^q$, 且两个范数均非零. 证明 Hölder 不等式 $\int\abs{fg}\dd\mu\le\norm f_p\norm g_q$取等, 当且仅当 $\abs f^p/\norm f_p^p=\abs g^q/\norm g_q^q$几乎处处.
\end{exercise}
\begin{solution}
固定 $B\ge0$, 对 $A\ge0$考虑 $F(A)=A^p/p-AB+B^q/q$. 若 $B>0$, 则 $F'(A)=A^{p-1}-B$, 所以唯一最小值在 $A=B^{1/(p-1)}$处取得, 最小值为零. 若 $B=0$, 仅在 $A=0$时为零. 因而 Young 不等式取等恰当且仅当 $A^p=B^q$.

令 $u=\abs f/\norm f_p$、$v=\abs g/\norm g_q$. Hölder 证明中的差是非负函数 $u^p/p+v^q/q-uv$, 且其积分为 $1-\int uv\dd\mu$. 积分为零当且仅当此函数几乎处处为零, 再用刚得到的 Young 等号条件, 即得所需刻画.
\end{solution}

\begin{exercise}\label{ra:c5:ex-dominated-lp}
设 $1\le p<\infty$, $f_n\to f$几乎处处, 且 $\abs{f_n}\le g$几乎处处对所有 $n$成立, 其中 $g$非负可测且 $g\in L^p$. 证明 $f\in L^p$且 $\norm{f_n-f}_p\to0$. 说明 $p=\infty$时对应结论为何不成立.
\end{exercise}
\begin{solution}
把可数个例外零测集取并, 在其外令 $n\to\infty$得 $\abs f\le g$, 因而 $f\in L^p$. 又有 $\abs{f_n-f}^p\le(2g)^p$, 右边可积, 左边几乎处处趋于零. 控制收敛定理给出 $\int\abs{f_n-f}^p\dd\mu\to0$, 取根即得结论.

在 $(0,1)$上取 $f_n=\ind_{(0,1/n)}$、$g=1$. 则 $f_n\to0$处处, 且 $\abs{f_n}\le g\in L^\infty$, 但 $\norm{f_n}_\infty=1$. 所以端点不能照搬有限指数的控制收敛结论.
\end{solution}

\begin{exercise}\label{ra:c5:ex-series}
设 $1\le p\le\infty$且 $\sum_{n=1}^\infty\norm{u_n}_p<\infty$. 证明 $\sum_nu_n$在 $L^p$中收敛. 当 $p<\infty$时, 证明还可选代表元使此级数几乎处处绝对收敛.
\end{exercise}
\begin{solution}
对部分和 $s_N=\sum_{n=1}^Nu_n$, 有 $\norm{s_M-s_N}_p\le\sum_{n=N+1}^M\norm{u_n}_p$, 因而 $(s_N)$是柯西列. 由 Riesz--Fischer 定理, 它收敛到某个 $s\in L^p$.

若 $p<\infty$, 令 $G_N=\sum_{n=1}^N\abs{u_n}$, 则 $\norm{G_N}_p\le\sum_n\norm{u_n}_p$. Fatou 引理给出 $G=\sum_n\abs{u_n}\in L^p$, 故 $G<\infty$几乎处处. 点态级数因而绝对收敛到可测函数 $v$, 并且尾和的同一估计给出 $\norm{v-s_N}_p\le\sum_{n>N}\norm{u_n}_p\to0$. 极限的唯一性说明 $v=s$几乎处处. 当 $p=\infty$时, 在各个 $\abs{u_n}\le\norm{u_n}_\infty$的共同零测集外, 级数甚至一致绝对收敛.
\end{solution}

\begin{exercise}\label{ra:c5:ex-powers}
设 $\alpha,\beta>0$, 在 $(0,\infty)$上令 $f(x)=x^{-\alpha}\ind_{(0,1)}+x^{-\beta}\ind_{[1,\infty)}$. 确定所有使 $f\in L^p$的 $1\le p\le\infty$. 另求序列 $a_n=n^{-\gamma}$的 $\ell^p$归属, 其中 $\gamma>0$.
\end{exercise}
\begin{solution}
有限指数时, $\int\abs f^p\dd x=\int_0^1x^{-\alpha p}\dd x+\int_1^\infty x^{-\beta p}\dd x$. 第一个积分有限当且仅当 $\alpha p<1$, 第二个有限当且仅当 $\beta p>1$. 因而可用指数恰为 $[1,\infty)\cap(1/\beta,1/\alpha)$. 任一端点取等时都出现 $\int x^{-1}\dd x$的对数发散, 故两个端点都不能因等号而加入. 由于零点附近本质无界, $f\notin L^\infty$; 若该区间为空, 没有任何可用指数.

对序列, 积分判别法给出 $\sum_nn^{-\gamma p}<\infty$当且仅当 $\gamma p>1$, 所以有限指数为 $\{p\ge1:\gamma p>1\}$. 另一方面 $\sup_nn^{-\gamma}=1$, 故总有 $a\in\ell^\infty$. 这也具体显示了函数空间与序列空间的不同包含方向.
\end{solution}

\begin{exercise}\label{ra:c5:ex-separable}
证明对 $1\le p<\infty$, $L^p(\R^n)$存在可数稠密子集.
\end{exercise}
\begin{solution}
取所有有限和 $s=\sum_{i=1}^m a_i\ind_{Q_i}$, 其中 $a_i\in\mathbb Q+i\mathbb Q$, $Q_i$为端点坐标均有理的有界半开矩形. 因为有限个可数选择仍只有可数种, 这些函数组成可数集合 $D$.

给定 $f\in L^p$及 $\eps>0$, 先取 $h\in C_c(\R^n)$使 $\norm{f-h}_p<\eps/2$. 选有理端点的大矩形 $B$, 使 $\supp h$位于其内部. 函数 $h$在 $\overline B$上一致连续. 以足够细的有理网格把 $B$分为有限个半开小矩形, 在每个小矩形中选一点, 并用 $\mathbb Q+i\mathbb Q$中的数逼近该点的函数值. 网格与系数可选得使所得 $s\in D$满足 $\abs{s-h}<\delta$在 $B$上成立; 在 $B$外两者都为零. 于是 $\norm{s-h}_p\le\delta\mu(B)^{1/p}$. 选 $\delta<\eps/(2\mu(B)^{1/p})$, 三角不等式给出 $\norm{f-s}_p<\eps$. 矩形边界是零测集, 不影响估计.
\end{solution}

\originalchapter{乘积测度与 Fubini 定理}
\label{ra:c6}
\sourcelecture{18--19}

\originalsection{截面与欧氏空间中的积分}

把 $\R^n$ 写成 $\R^\ell\times\R^m$, 其中 $n=\ell+m$. 多重积分的一个基本问题是: 能否先固定 $y$, 对 $x$ 积分, 再对 $y$ 积分? 对连续函数和有界矩形, 这已经熟悉; 对一般可测函数, 必须先解决截面的可测性, 再解决两个积分能否交换.

\begin{definition}[截面]
若 $E\subset\R^\ell\times\R^m$, 定义 $E_y=\{x:(x,y)\in E\}$ 和 $E^x=\{y:(x,y)\in E\}$. 对函数 $f(x,y)$, 记 $f_y(x)=f(x,y)$, $f^x(y)=f(x,y)$. 这些分别是集合和函数的截面.
\end{definition}

这里的 Lebesgue 可测性包含零测集的任意子集. 正是这一点使得截面不能对每个参数都保持可测.

\begin{example}
取不可测集 $V\subset\R^\ell$ 和一点 $y_0\in\R^m$, 令 $E=V\times\{y_0\}$. 集合 $\R^\ell\times\{y_0\}$ 是零测集, 因而 $E$ 是 Lebesgue 可测的零测集. 当 $y\ne y_0$ 时 $E_y=\varnothing$, 但 $E_{y_0}=V$ 不可测.
\end{example}

零测超平面的断言可先在有界块上验证: $[-k,k]^\ell\times\{y_0\}$ 能放在体积任意小的矩形中, 再取可数并. 因此下面的结论只能要求截面在几乎处处的参数处可测.

\begin{lemma}[集合的截面积公式]\label{ra:c6:cavalieri}
若 $E\subset\R^{\ell+m}$ Lebesgue 可测, 则 $E_y$ 对几乎处处的 $y$ Lebesgue 可测. 在例外零测集上任意赋值后, $y\mapsto\lambda_\ell(E_y)$ 是 Lebesgue 可测函数, 且
\[
 \lambda_{\ell+m}(E)=\int_{\R^m}\lambda_\ell(E_y)\,\dd y.
\]
两边可以同时为 $+\infty$.
\end{lemma}
\begin{proof}
保留原讲义由规则集合逐步逼近一般集合的路线. 在下面各步中, 不仅证明积分等式, 也同时证明截面积函数可测.

若 $J=J_1\times J_2$ 是半开矩形, 则 $\lambda_\ell(J_y)=\lambda_\ell(J_1)\ind_{J_2}(y)$, 积分等于 $\lambda_{\ell+m}(J)$. 开集 $G$ 可以分解为可数个两两不交的半开矩形 $J_k$. 截面也两两不交, 所以 $\lambda_\ell(G_y)=\sum_k\lambda_\ell((J_k)_y)$; 单调收敛定理给出积分等式.

若 $K$ 紧, 取有界开集 $G\supset K$. 对每个 $y$, $K_y$ 紧, 且 $\lambda_\ell(K_y)=\lambda_\ell(G_y)-\lambda_\ell((G\setminus K)_y)$. 两个被减函数可测; 它们的积分都有限, 因而相减得到 $\int\lambda_\ell(K_y)\,\dd y=\lambda_{\ell+m}(K)$.

若 $K_j$ 递增且每个 $K_j$ 紧, 令 $B=\bigcup_jK_j$. 截面递增, 所以 $\lambda_\ell(B_y)=\lim_j\lambda_\ell((K_j)_y)$. 再用单调收敛, 公式对 $B$ 成立. 若 $G_j$ 是递减的有界开集, 令 $C=\bigcap_jG_j$, 并取紧矩形 $Q\supset G_1$. 集合 $Q\setminus C=\bigcup_j(Q\setminus G_j)$ 是递增紧集之并. 用刚才的结论, 再从 $Q$ 的截面积和体积中相减, 公式对 $C$ 成立.

现在设 $E$ 有界且可测. 测度的正则性允许选择递增紧集 $K_j$ 和递减有界开集 $G_j$, 使 $K_j\subset E\subset G_j$, 且 $\lambda(K_j)\uparrow\lambda(E)$, $\lambda(G_j)\downarrow\lambda(E)$. 记 $B=\bigcup_jK_j$, $C=\bigcap_jG_j$. 则 $B\subset E\subset C$, $\lambda(B)=\lambda(C)=\lambda(E)$. 前两步给出
\[
 \int_{\R^m}\bigl(\lambda_\ell(C_y)-\lambda_\ell(B_y)\bigr)\,\dd y=0.
\]
被积函数非负, 所以 $\lambda_\ell(C_y\setminus B_y)=0$ 对几乎处处的 $y$ 成立. 对这些 $y$, $B_y\subset E_y\subset C_y$, 由 Lebesgue 测度的完备性, $E_y$ 可测且 $\lambda_\ell(E_y)=\lambda_\ell(B_y)$. 这既给出可测性, 也给出所求等式.

最后对一般 $E$ 使用 $E\cap[-j,j]^n$. 这些集合递增到 $E$. 去掉每一步的例外零测集的可数并后, 所有截面同时可测并递增到 $E_y$. 单调收敛把公式传到 $E$.
\end{proof}

\begin{theorem}[Tonelli 定理, Lebesgue 版本]\label{ra:c6:tonelli-euclid}
设 $f:\R^\ell\times\R^m\to[0,\infty]$ Lebesgue 可测. 则 $f_y$ 对几乎处处的 $y$ 可测, $F(y)=\int_{\R^\ell}f_y(x)\,\dd x$ 在例外集上任意延拓后可测, 并且
\[
 \int_{\R^{\ell+m}}f(x,y)\,\dd(x,y)
 =\int_{\R^m}\left(\int_{\R^\ell}f(x,y)\,\dd x\right)\dd y.
\]
交换 $x,y$ 后同样成立. 各积分允许等于 $+\infty$.
\end{theorem}
\begin{proof}
对 $f=\ind_E$, 结论就是引理 \ref{ra:c6:cavalieri}. 由有限线性组合, 结论对非负简单函数成立. 选择非负简单函数 $s_j\uparrow f$. 去掉所有 $s_j$ 的截面例外集后, 对其余的 $y$, $(s_j)_y\uparrow f_y$, 因而 $f_y$ 可测, 且
$F(y)=\lim_j\int(s_j)_y\,\dd x$. 此极限递增, 所以 $F$ 可测. 对内层积分和外层积分分别用单调收敛, 得到 $\int F=\lim_j\int s_j=\int f$. 对另一种顺序重复同样论证.
\end{proof}

\begin{theorem}[Fubini 定理, Lebesgue 版本]\label{ra:c6:fubini-euclid}
若 $f\in L^1(\R^{\ell+m})$, 则 $f_y\in L^1(\R^\ell)$ 对几乎处处的 $y$ 成立. 函数 $F(y)=\int f_y\,\dd x$ 属于 $L^1(\R^m)$, 且 $\int F\,\dd y=\int f\,\dd(x,y)$. 两个积分顺序可以交换.
\end{theorem}
\begin{proof}
先对 $|f|$ 用 Tonelli 定理. 因为 $\int|f|<\infty$, 函数 $A(y)=\int|f_y|\,\dd x$ 几乎处处有限, 且 $\int A<\infty$. 所以截面几乎处处绝对可积. 若 $f$ 实值, 对 $f^+,f^-$ 用 Tonelli 后相减; 两者的积分有限, 不会出现 $\infty-\infty$. 又有 $|F(y)|\le A(y)$, 故 $F\in L^1$. 对复值函数分别处理实部和虚部.
\end{proof}

实际计算时, 非负性和绝对可积性承担不同作用. 对非负函数可先交换次序来判定积分是否有限; 对有正负变化的函数, 通常先用 $|f|$ 建立有限性, 再交换 $f$ 的积分.

\begin{example}
原讲 18 的积分计算: 在 $E=(0,\infty)\times(0,1)$ 上令 $f(x,y)=y\sin x\,e^{-xy}$. 先计算 $\int_Ef$, 再求一个一元反常积分.
\end{example}
\begin{solution}
因为 $|f(x,y)|\le ye^{-xy}$, Tonelli 定理给出
$\int_E|f|\le\int_0^1(\int_0^\infty ye^{-xy}\,\dd x)\dd y=1$.
因此可用 Fubini. 对固定的 $y>0$, 两次分部积分, 或对 $\int_0^\infty e^{(-y+i)x}\,\dd x$ 取虚部, 得到 $\int_0^\infty e^{-yx}\sin x\,\dd x=(1+y^2)^{-1}$. 故
$\int_Ef=\int_0^1y(1+y^2)^{-1}\,\dd y=\tfrac12\log2$.

另一方面, 对 $x>0$ 分部积分给出
$\int_0^1ye^{-xy}\,\dd y=(1-e^{-x})/x^2-e^{-x}/x$.
于是同一个积分等于
\[
 \int_0^\infty\frac{\sin x}{x}
 \left(\frac{1-e^{-x}}{x}-e^{-x}\right)\dd x
 =\frac12\log2.
\]
这里并非仅因两个一元积分各自存在就交换次序; 交换的依据是前面的绝对可积估计.
\end{solution}

\begin{example}\label{ra:c6:conditional}
作为编者补充, 在 $(0,1)^2$ 上令 $f(x,y)=(x^2-y^2)/(x^2+y^2)^2$. 两个迭代积分分别为 $\pi/4$ 和 $-\pi/4$, 而 $f\notin L^1((0,1)^2)$.
\end{example}
\begin{solution}
对固定 $x>0$, 有 $\partial_y[y/(x^2+y^2)]=f(x,y)$, 所以 $\int_0^1f(x,y)\,\dd y=(1+x^2)^{-1}$. 对固定 $y>0$, 有 $\partial_x[-x/(x^2+y^2)]=f(x,y)$, 所以 $\int_0^1f(x,y)\,\dd x=-(1+y^2)^{-1}$. 再积分即得所述数值.

为核实绝对不可积, 考虑 $0<x<1$, $0<y<x/2$. 此处 $x^2-y^2\ge3x^2/4$, $(x^2+y^2)^2\le25x^4/16$, 所以 $|f|\ge12/(25x^2)$. 于是
$\int_0^1\int_0^{x/2}|f|\,\dd y\,\dd x\ge(6/25)\int_0^1\dd x/x=\infty$.
\end{solution}

\originalsection{一般乘积空间}

对抽象空间, 没有开集、紧集可供逼近. 但可测矩形足以生成所需的 $\sigma$-代数. 为把矩形上的等式传到所有可测集, 先补充一个集合论工具.

\begin{lemma}[$\pi$-$\lambda$ 引理, 编者补充]\label{ra:c6:pi-lambda}
设 $\mathcal P$ 含全集并对有限交封闭. 若集合族 $\mathcal D$ 含全集, 对补集和可数不交并封闭, 且 $\mathcal P\subset\mathcal D$, 则 $\sigma(\mathcal P)\subset\mathcal D$.
\end{lemma}
\begin{proof}
把含 $\mathcal P$ 的最小此类集合族记为 $\mathcal D_0$. 这种集合族称为 Dynkin 族. 若 $A\subset B$ 且 $A,B\in\mathcal D_0$, 则 $B\setminus A=(A\cup B^c)^c\in\mathcal D_0$.

对固定的 $A\in\mathcal P$, 集合族 $\{B:A\cap B\in\mathcal D_0\}$ 是 Dynkin 族: 补集步骤使用 $A\setminus(A\cap B)$, 不交并步骤使用交对并的分配律. 它含 $\mathcal P$, 故含 $\mathcal D_0$. 因此 $A\cap B\in\mathcal D_0$ 对 $A\in\mathcal P$, $B\in\mathcal D_0$ 成立. 再固定 $B\in\mathcal D_0$, 同样考虑 $\{A:A\cap B\in\mathcal D_0\}$, 得到 $\mathcal D_0$ 对有限交封闭. 对任意可数并, 可通过依次减去前面的集合改写成不交并; 所以 $\mathcal D_0$ 是 $\sigma$-代数. 它含 $\mathcal P$, 从而含 $\sigma(\mathcal P)$.
\end{proof}

\begin{definition}[乘积 $\sigma$-代数]
给定测度空间 $(X,\mathcal M,\mu)$ 和 $(Y,\mathcal N,\nu)$, 记
$\mathcal M\otimes\mathcal N=\sigma\{A\times B:A\in\mathcal M,\ B\in\mathcal N\}$.
其生成元称为可测矩形. 测度空间称为 $\sigma$-有限, 是指全集能写成可数个有限测度的可测集之并.
\end{definition}

\begin{proposition}\label{ra:c6:sections}
若 $E\in\mathcal M\otimes\mathcal N$, 则 $E_y\in\mathcal M$ 对每个 $y$ 成立, $E^x\in\mathcal N$ 对每个 $x$ 成立. 若 $f$ 对乘积 $\sigma$-代数可测, 则每个函数截面也可测.
\end{proposition}
\begin{proof}
使所有 $E_y$ 可测的集合 $E$ 组成一个 $\sigma$-代数, 因为截面运算与补集、可数并相容; 可测矩形属于它, 所以它含 $\mathcal M\otimes\mathcal N$. 对另一种截面同理. 对函数, 对开集 $U$ 使用 $(f^{-1}(U))_y=f_y^{-1}(U)$ 即可.
\end{proof}

\begin{theorem}[乘积测度的构造]\label{ra:c6:product}
若 $\mu,\nu$ 都 $\sigma$-有限, 则对每个 $E\in\mathcal M\otimes\mathcal N$, 函数 $y\mapsto\mu(E_y)$ 和 $x\mapsto\nu(E^x)$ 都可测, 且
\[
 (\mu\times\nu)(E):=\int_Y\mu(E_y)\,\dd\nu(y)
 =\int_X\nu(E^x)\,\dd\mu(x).
\]
此式定义 $\mathcal M\otimes\mathcal N$ 上的测度, 并且它是满足
$(\mu\times\nu)(A\times B)=\mu(A)\nu(B)$ 的唯一测度. 这里约定 $0\cdot\infty=0$.
\end{theorem}
\begin{proof}
先设 $\mu(X),\nu(Y)<\infty$. 令 $\mathcal D$ 为同时满足两种截面积函数可测及两种积分相等的集合族. 可测矩形在 $\mathcal D$ 中. 全集在其中; 对补集, 使用 $\mu((E^c)_y)=\mu(X)-\mu(E_y)$ 和对应的 $\nu$ 公式, 两边均只作有限量的相减; 对可数不交并, 截面仍不交, 用可数可加性及单调收敛. 所以 $\mathcal D$ 是 Dynkin 族. 由引理 \ref{ra:c6:pi-lambda}, $\mathcal D$ 包含整个乘积 $\sigma$-代数.

一般情形下, 把 $X$ 和 $Y$ 分别分成可数个两两不交的有限测度集 $X_i,Y_j$. 对 $E\cap(X_i\times Y_j)$ 在有限空间 $X_i\times Y_j$ 上使用刚才的结论. 例如
$\mu(E_y)=\sum_i\mu((E\cap(X_i\times Y_j))_y)$ 当 $y\in Y_j$.
因此完整的截面积函数由可数和及沿 $Y_j$ 的拼接得到, 仍可测. 积分的可数可加性把每个有限块上的两种等式相加, 给出全空间上的等式.

定义的 $(\mu\times\nu)(E)$ 非负且空集的测度为零. 若 $E_k$ 两两不交, 每个截面也两两不交, 因而
$\mu((\bigcup_kE_k)_y)=\sum_k\mu((E_k)_y)$; 单调收敛证明乘积测度可数可加. 矩形公式直接从 $\mu((A\times B)_y)=\mu(A)\ind_B(y)$ 得到.

若另一测度 $\eta$ 也满足矩形公式, 在有限块 $X_i\times Y_j$ 上, $\eta$ 和 $\mu\times\nu$ 都有限. 两个有限测度相等的集合构成 Dynkin 族, 因为总测度相等并可从补集中相减; 它含该块中的可测矩形, 故由引理包含该块的全部乘积可测集. 各块再相加, 得 $\eta=\mu\times\nu$.
\end{proof}

\begin{theorem}[抽象 Tonelli 与 Fubini 定理]\label{ra:c6:abstract}
设 $\mu,\nu$ 都 $\sigma$-有限, 且 $f$ 对 $\mathcal M\otimes\mathcal N$ 可测.
若 $0\le f\le\infty$, 两个截面积分函数都可测, 并有
\[
 \int_{X\times Y}f\,\dd(\mu\times\nu)
 =\int_Y\left(\int_Xf(x,y)\,\dd\mu(x)\right)\dd\nu(y)
 =\int_X\left(\int_Yf(x,y)\,\dd\nu(y)\right)\dd\mu(x).
\]
若 $f$ 为复值且 $\int|f|\,\dd(\mu\times\nu)<\infty$, 则截面几乎处处属于相应的 $L^1$, 截面积分函数属于 $L^1$, 同样的等式成立.
\end{theorem}
\begin{proof}
指示函数的结论由定理 \ref{ra:c6:product} 得到. 对非负简单函数作有限线性组合, 再用非负简单函数递增逼近和两次单调收敛, 得到非负函数的结论. 可积情形先对 $|f|$ 用前一部分得到截面的绝对可积性, 再对实部、虚部的正负部分相减, 与定理 \ref{ra:c6:fubini-euclid} 的证明相同.
\end{proof}

\begin{example}
作为编者补充, 检查 $\sigma$-有限性的作用. 在 $X=Y=[0,1]$ 上, $X$ 取 Lebesgue 测度 $\mu$, $Y$ 取计数测度 $\nu$. 对角线 $D=\{(x,y):x=y\}$ 是乘积可测集, 但
$\int_X\nu(D^x)\,\dd\mu(x)=1$, 而 $\int_Y\mu(D_y)\,\dd\nu(y)=0$.
\end{example}
\begin{solution}
对角线是闭集, 而方形的开集是可数个开矩形之并, 所以 $D$ 乘积可测. 每个截面是单点: 其计数测度为 $1$, Lebesgue 测度为 $0$, 两个积分便得到上述值. 这里 $\nu$ 不 $\sigma$-有限, 因为有限计数测度的集合只能是有限集, 可数个有限集不能覆盖 $[0,1]$. 因此这个例子没有违反前述定理.
\end{solution}

\originalsection{完备化及 Lebesgue 乘积}

\begin{proposition}[完备乘积上的截面]\label{ra:c6:completion}
设两个因子空间完备且 $\sigma$-有限, 并把乘积测度空间完备化. 对完备化可测的非负函数或可积函数, 抽象 Tonelli 或 Fubini 仍成立, 但截面的可测性改为对几乎处处的参数成立.
\end{proposition}
\begin{proof}
完备化可测集与一个乘积可测集的对称差包含在乘积可测零测集内. 对函数的简单函数逼近逐项使用此事实, 并把例外零测集取可数并, 可得乘积可测函数 $g$ 及乘积可测零测集 $N$, 使 $f=g$ 于 $N^c$. 例如先对非负 $f$ 的递增简单逼近选代表, 再在 $N^c$ 上取极限, 在 $N$ 上置零; 实部、虚部分别处理即可.

由 Tonelli, $\int_Y\mu(N_y)\,\dd\nu(y)=0$, 所以 $\mu(N_y)=0$ 对几乎处处的 $y$ 成立. 因子 $\mu$ 完备, 保证 $f_y$ 在这些参数处的零测修改仍可测, 且积分与 $g_y$ 相等. 另一种截面同理. 全空间积分也不受零测修改影响, 因而把 $g$ 的公式传给 $f$.
\end{proof}

\begin{theorem}\label{ra:c6:lebesgue-product}
若 $\mathcal L^\ell,\mathcal L^m,\mathcal L^{\ell+m}$ 分别是对应欧氏空间的 Lebesgue $\sigma$-代数, 则 $\mathcal L^{\ell+m}$ 是 $\mathcal L^\ell\otimes\mathcal L^m$ 相对于 $\lambda_\ell\times\lambda_m$ 的完备化, 完备后的测度就是 $\lambda_{\ell+m}$.
\end{theorem}
\begin{proof}
先证 Lebesgue 可测矩形 $A\times B$ 在 $\R^{\ell+m}$ 中可测. 当 $\lambda_\ell(A),\lambda_m(B)<\infty$ 时, 对 $\delta>0$, 用正则性取 $K_1\subset A\subset G_1$, $K_2\subset B\subset G_2$, 其中 $K_i$ 紧, $G_i$ 开且有限测度, $\lambda(G_i\setminus K_i)<\delta$. 紧积和开积都 Borel 可测. 集合截面积公式给出它们的测度分别为 $\lambda(K_1)\lambda(K_2)$ 和 $\lambda(G_1)\lambda(G_2)$. 于是
\[
 \lambda(G_1\times G_2)-\lambda(K_1\times K_2)
 \le\delta\bigl(\lambda(A)+\lambda(B)+2\delta\bigr).
\]
两者夹住 $A\times B$, 差的测度能任意小, 所以正则性判别给出 $A\times B$ 的可测性. 一般情形用 $A\cap[-k,k]^\ell$ 和 $B\cap[-k,k]^m$ 的递增并. 对已证可测的矩形再用截面积公式, 得其测度为 $\lambda_\ell(A)\lambda_m(B)$. 因此 $\mathcal L^\ell\otimes\mathcal L^m\subset\mathcal L^{\ell+m}$; 乘积测度的唯一性说明 $\lambda_{\ell+m}$ 在此 $\sigma$-代数上的限制就是 $\lambda_\ell\times\lambda_m$.

另一方面, $\R^{\ell+m}$ 的开集是可数个开矩形之并, 所以全部 Borel 集都在 $\mathcal L^\ell\otimes\mathcal L^m$ 中. 每个 Lebesgue 可测集都与某个 Borel 集只差一个 Borel 零测集的子集, 因而属于该乘积空间的完备化. 反过来, Lebesgue 测度本来完备, 所以乘积可测集的零测修改仍是 Lebesgue 可测集. 两个方向合起来即得结论.
\end{proof}

\begin{remark}
原讲 19 的抽象乘积构造和完备化结论只作陈述. 上面的 $\pi$-$\lambda$ 证明、唯一性证明以及完备化论证是编者补充. 原讲的可测矩形估计中有一个把并集估计写成等式的笔误; 此处直接比较开积和紧积的体积, 避免重叠计数. “乘积可测”与“完备乘积可测”的区别解释了本章开头的异常截面.
\end{remark}

\originalsection{习题与解答}

\begin{exercise}
设 $f\in L^1(X,\mu)$, $g\in L^1(Y,\nu)$, 两个空间 $\sigma$-有限. 证明 $h(x,y)=f(x)g(y)$ 可积, 并求其积分和 $L^1$ 范数.
\end{exercise}
\begin{solution}
坐标投影可测, 所以 $h$ 乘积可测. 对 $|f(x)g(y)|$ 用 Tonelli, 得 $\int|h|=\norm{f}_1\norm{g}_1<\infty$. 再用 Fubini, 得 $\int h=(\int f\,\dd\mu)(\int g\,\dd\nu)$. 该证明对复值函数也成立.
\end{solution}

\begin{exercise}
设 $f\ge0$ 在 $\R^{\ell+m}$ 上可测. 已知 $\int_{\R^m}(\int_{\R^\ell}f\,\dd x)\dd y<\infty$. 证明 $f\in L^1$, 并解释若删去 $f\ge0$, 为什么“内层积分存在且外层积分有限”不够.
\end{exercise}
\begin{solution}
非负时, Tonelli 将所给的有限值等同于 $\int f=\int|f|$, 所以可积. 有符号时, 内层积分可能通过正负抵消变小, 而 $\int|f|$ 仍发散. 例题 \ref{ra:c6:conditional} 给出了两种迭代积分都有限但不相等的具体反例.
\end{solution}

\begin{exercise}
若 $N\subset\R^{\ell+m}$ 为 Lebesgue 零测集, 证明 $\lambda_\ell(N_y)=0$ 对几乎处处的 $y$ 成立, 并说明为何不能把“几乎处处”改成“每个”.
\end{exercise}
\begin{solution}
Tonelli 给出 $\int\lambda_\ell(N_y)\,\dd y=\lambda_{\ell+m}(N)=0$, 非负函数积分为零即几乎处处为零. 对 $N=\R^\ell\times\{y_0\}$, 在 $y_0$ 处的截面是整个 $\R^\ell$, 故不为零; 还可把此截面换成不可测集, 得到本章开头的例子.
\end{solution}

\originalchapter{卷积与光滑逼近}
\label{ra:c7}
\sourcelecture{20--21}

\originalsection{卷积的存在性}

平移不改变 Lebesgue 测度. 把函数与一族平移后的核相乘再积分, 就得到卷积. 它一方面把两个函数组合成新函数, 另一方面能用作平滑和逼近的工具. 但两个可积函数的乘积未必可积, 因而卷积定义中的积分存在性需要证明.

\begin{definition}
对 $\R^n$ 上的可测函数 $f,g$, 在积分绝对收敛的点定义
$(f*g)(x)=\int_{\R^n}f(y)g(x-y)\,\dd y$.
若 $f,g\ge0$, 允许此积分为 $+\infty$. 记 $\tau_hf(x)=f(x-h)$.
\end{definition}

作为 $(x,y)$ 的函数, $f(y)g(x-y)$ Lebesgue 可测: $f(y)$ 由坐标投影得到, 而 $g(x-y)$ 可先从乘积中的 $g(u)$ 得到, 再使用可逆线性变换 $(x,y)\mapsto(x-y,y)$. 可逆线性变换保持 Lebesgue 可测性, 所以可以使用上一章的 Tonelli 定理.

\begin{theorem}[$L^1$ 卷积]\label{ra:c7:l1}
若 $f,g\in L^1(\R^n)$, 则卷积在几乎处处的 $x$ 绝对存在, 定义一个 $L^1$ 等价类, 并有 $\norm{f*g}_1\le\norm{f}_1\norm{g}_1$. 若 $f,g\ge0$, 此不等式为等式. 此外,
$\int_{\R^n}(f*g)= (\int_{\R^n}f)(\int_{\R^n}g)$.
\end{theorem}
\begin{proof}
对非负函数 $|f(y)||g(x-y)|$ 用 Tonelli 和平移不变性, 得
\[
 \int_{\R^n}\!\int_{\R^n}|f(y)||g(x-y)|\,\dd y\,\dd x
 =\int_{\R^n}|f(y)|\left(\int_{\R^n}|g(x-y)|\,\dd x\right)\dd y
 =\norm{f}_1\norm{g}_1.
\]
右边有限, 所以内层绝对积分几乎处处有限. 又 $|f*g|\le|f|*|g|$, 所以得到范数估计. 非负时没有绝对值引起的损失, 故等号成立. 对 $f(y)g(x-y)$ 再使用 Fubini, 得积分的乘法公式.
\end{proof}

\begin{proposition}[交换律和结合律]\label{ra:c7:algebra}
对 $f,g\in L^1(\R^n)$, $f*g=g*f$ 几乎处处. 对 $f,g,h\in L^1(\R^n)$, $(f*g)*h=f*(g*h)$ 几乎处处. 卷积是双线性的, 并且不依赖代表函数的零测修改.
\end{proposition}
\begin{proof}
交换律来自代换 $u=x-y$. 要证明结合律, 先对三重绝对积分使用 Tonelli:
$\int\!\int\!\int|f(u)g(v)h(x-u-v)|\,\dd u\,\dd v\,\dd x
=\norm{f}_1\norm{g}_1\norm{h}_1<\infty$.
故对几乎处处的 $x$, 关于 $u,v$ 的积分绝对收敛. 用 Fubini 交换次序并作平移代换, 得
\[
 ((f*g)*h)(x)
 =\int\!\int f(u)g(v)h(x-u-v)\,\dd v\,\dd u
 =(f*(g*h))(x).
\]
双线性直接来自积分的线性. 若 $f$ 或 $g$ 改变在零测集上, 上述绝对积分估计应用于差函数, 表明卷积只改变在零测集上.
\end{proof}

\begin{example}
若 $f=\ind_{[0,1]}$, 则 $(f*f)(x)$ 等于区间 $[0,1]\cap[x-1,x]$ 的长度. 因而在 $0\le x\le1$ 时为 $x$, 在 $1\le x\le2$ 时为 $2-x$, 其余为零. 两个有跳跃的函数卷积后产生了连续的三角形函数, 且 $\norm{f*f}_1=1=\norm{f}_1^2$.
\end{example}

\originalsection{Young 不等式}

\begin{theorem}[Young 不等式]\label{ra:c7:young}
设 $p,q,r\in[1,\infty]$, 约定 $1/\infty=0$, 并满足
$1/r=1/p+1/q-1$. 若 $f\in L^p(\R^n)$, $g\in L^q(\R^n)$, 则卷积几乎处处绝对存在, 属于 $L^r$, 且
$\norm{f*g}_r\le\norm{f}_p\norm{g}_q$.
\end{theorem}
\begin{proof}
可先换成 $|f|,|g|$, 因为 $|f*g|\le|f|*|g|$. 范数为零的情形直接成立; 其余通过除以各自范数, 只须证明非负函数且 $\norm{f}_p=\norm{g}_q=1$ 的情形.

先设 $r<\infty$. 则 $p,q<\infty$, 且由指数关系得 $r\ge p,q$. 当 $p=q=r=1$ 时, 已由定理 \ref{ra:c7:l1} 证明. 其余情形用三个因子的 Hölder 不等式, 指数分别为 $r,q',p'$, 因为 $1/r+1/q'+1/p'=1$. 这种形式可由前章的两因子版本得到: 先以 $r,r'$ 处理第一因子和后两因子的乘积, 再对后两因子的 $r'$ 次幂使用共轭指数 $q'/r'$ 与 $p'/r'$, 得 $\norm{bc}_{r'}\le\norm{b}_{q'}\norm{c}_{p'}$; 指数为无穷时以本质上界估计相应因子. 将被积函数分解为
$[f(y)^{p/r}g(x-y)^{q/r}]\,f(y)^{1-p/r}\,g(x-y)^{1-q/r}$, 得
\[
 (f*g)(x)\le
 \left(\int f(y)^pg(x-y)^q\,\dd y\right)^{1/r}
 \left(\int f^{(1-p/r)q'}\right)^{1/q'}
 \left(\int g^{(1-q/r)p'}\right)^{1/p'}.
\]
若 $q'= \infty$ 或 $p'=\infty$, 相应因子本身是常数 $1$, 其 $L^\infty$ 范数为 $1$, 不写成积分的幂. 对有限的指数, 关系 $(1-p/r)q'=p$ 和 $(1-q/r)p'=q$ 表明后两个因子都是 $1$. 第一个积分几乎处处有限, 因为 $f^p,g^q\in L^1$, 可用定理 \ref{ra:c7:l1}. 所以
$(f*g)(x)^r\le(f^p*g^q)(x)$ 几乎处处. 再积分, 得 $\norm{f*g}_r^r\le\norm{f^p}_1\norm{g^q}_1=1$.

若 $r=\infty$, 指数关系成为 $q=p'$. 对每个 $x$ 用 Hölder 及平移不变性, 得
$\int|f(y)g(x-y)|\,\dd y\le\norm{f}_p\norm{g}_q$.
这也包含 $p=1,q=\infty$ 及 $p=\infty,q=1$ 的端点; 若其中一个指数为 $\infty$, 指数关系必然属于这些情形. 因而所有端点都已覆盖.
\end{proof}

\begin{remark}
原讲 20 的归一化一行把 $\norm{g}_q$ 误写为 $\norm{g}_p$, 并略过部分端点. 此处保留三个 Hölder 因子的证明, 补齐指数为 $1$ 或 $\infty$ 时的解释. 在第 \ref{ra:c8} 章还会从广义 Minkowski 不等式重新得到 $L^1*L^p\subset L^p$.
\end{remark}

\begin{lemma}[平移的 $L^p$ 连续性]\label{ra:c7:translation}
若 $1\le p<\infty$ 且 $f\in L^p(\R^n)$, 则 $\norm{\tau_hf-f}_p\to0$ 当 $h\to0$.
\end{lemma}
\begin{proof}
由前章连续紧支集函数的稠密性, 对 $\eps>0$ 取 $u\in C_c$ 使 $\norm{f-u}_p<\eps$. 平移是 $L^p$ 等距算子, 因而
$\norm{\tau_hf-f}_p\le2\eps+\norm{\tau_hu-u}_p$.
$u$ 一致连续, 所以 $\norm{\tau_hu-u}_\infty\to0$. 当 $|h|\le1$ 时, 差函数的支集包含在一个固定有限测度的球 $K$ 中, 因而 $\norm{\tau_hu-u}_p\le\lambda(K)^{1/p}\norm{\tau_hu-u}_\infty\to0$. 先令 $h\to0$, 再令 $\eps\to0$.
\end{proof}

\begin{theorem}[共轭指数下的连续性]\label{ra:c7:continuous}
设 $1\le p\le\infty$, $f\in L^p(\R^n)$, $g\in L^{p'}(\R^n)$. 则积分 $(f*g)(x)$ 对每个 $x$ 绝对存在, 且 $f*g$ 有界并一致连续. 若 $1<p<\infty$, 则 $f*g\in C_0(\R^n)$, 即当 $|x|\to\infty$ 时趋于零.
\end{theorem}
\begin{proof}
逐点存在和有界性已由 Hölder 给出. $p,p'$ 中至少一个有限. 若 $p'<\infty$, 则
$|(f*g)(x+h)-(f*g)(x)|\le\norm{f}_p\norm{\tau_hg-g}_{p'}$, 右边与 $x$ 无关且趋于零. 若 $p'=\infty$, 则 $p=1$; 用交换律把平移放在 $f$ 上, 得到 $\norm{g}_\infty\norm{\tau_hf-f}_1$. 因而一致连续.

若 $1<p<\infty$, 两个指数均有限, 可选 $f_k,g_k\in C_c$, 分别在 $L^p,L^{p'}$ 中逼近 $f,g$. 每个 $f_k*g_k$ 连续且紧支集. 又
$\norm{f_k*g_k-f*g}_\infty
\le\norm{f_k-f}_p\norm{g_k}_{p'}+\norm{f}_p\norm{g_k-g}_{p'}\to0$.
因为 $g_k$ 的范数有界, 右边确实趋于零. 对给定 $\eps>0$, 先取一个足够接近的 $f_k*g_k$, 再在其支集以外估计, 得 $|f*g|<\eps$. 因而 $f*g$ 在无穷远消失.
\end{proof}

\begin{remark}
原讲 20 的此定理漏写 $g\in L^{p'}$; 这是证明中 Hölder 不等式所需的条件. 两个端点不能无条件推出在无穷远消失: 取 $f\in L^1$ 且 $\int f\ne0$, $g\equiv1\in L^\infty$, 则 $f*g\equiv\int f$.
\end{remark}

\originalsection{光滑核}

要把卷积变成逼近, 应使核的积分为 $1$, 并使核的质量集中到原点附近. 若还要求核无限可微, 可以从一个在边界处所有导数都为零的函数出发.

\begin{lemma}[光滑紧支集核]\label{ra:c7:bump}
存在 $\phi\in C_c^\infty(\R^n)$, 满足 $\phi\ge0$, $\supp\phi\subset\overline{B(0,1)}$, $\int\phi=1$. 对 $a>0$, 令 $\phi_a(x)=a^{-n}\phi(x/a)$, 则 $\int\phi_a=1$, $\supp\phi_a\subset\overline{B(0,a)}$.
\end{lemma}
\begin{proof}
定义 $h(t)=e^{-1/t}$ 当 $t>0$, $h(t)=0$ 当 $t\le0$. 在 $t>0$ 上, 每个导数都是某个关于 $1/t$ 的多项式乘以 $e^{-1/t}$. 对每个整数 $k\ge0$, $s^ke^{-s}\to0$ 当 $s\to\infty$, 所以这些右侧导数在 $0$ 都趋于零. 归纳考察差商: 若前一阶导数在右边具有此形式且在左边为零, 除以 $t$ 后仍趋于零, 故其在 $0$ 的导数为零. 由此 $h\in C^\infty(\R)$.

函数 $\psi(x)=h(1-|x|^2)$ 无限可微, 在单位球外为零, 在球内为正. 它的积分有限且正. 令 $\phi=\psi/\int\psi$. 伸缩代换 $u=x/a$ 给出 $\int\phi_a=\int\phi=1$, 支集性质也直接得到.
\end{proof}

\begin{definition}[局部可积与平滑化]
若 $f$ 在每个紧集上绝对可积, 称 $f\in L^1_{\mathrm{loc}}(\R^n)$. 对这样的 $f$, 定义 $f_a=f*\phi_a$, 称为 $f$ 的平滑化.
\end{definition}

\begin{theorem}[平滑化的导数与支集]\label{ra:c7:mollify}
若 $f\in L^1_{\mathrm{loc}}(\R^n)$, 则 $f_a(x)$ 对每个 $x$ 存在, $f_a\in C^\infty(\R^n)$, 且对每个多重指标 $\alpha$,
$\partial^\alpha f_a=f*(\partial^\alpha\phi_a)$.
若 $f$ 几乎处处在紧集 $K$ 外为零, 则 $\supp f_a\subset K+\overline{B(0,a)}$, 因而 $f_a\in C_c^\infty$.
\end{theorem}
\begin{proof}
对固定 $x$, 积分只用到 $y\in\overline{B(x,a)}$, 且核有界, 所以绝对存在. 要证连续性, 将 $x$ 限制在固定的紧邻域 $V$. 所有涉及的 $y$ 都位于紧集 $V+\overline{B(0,a)}$; 被积函数被常数乘以该集上的 $|f|$ 控制. 核连续, 控制收敛给出 $f_a$ 连续.

求第 $j$ 个偏导时, 把差商写成
$\int f(y)[\phi_a(x+te_j-y)-\phi_a(x-y)]/t\,\dd y$.
当 $|t|\le1$ 时, 核的差商由 $\norm{\partial_j\phi_a}_\infty$ 控制, 所有支集仍落在一个固定紧集上. 控制收敛允许把 $t\to0$ 的极限移入积分, 得 $\partial_jf_a=f*\partial_j\phi_a$. 对右边重复连续性论证及差商论证, 归纳得到任意阶偏导, 且均连续.

若 $x\notin K+\overline{B(0,a)}$, 则对 $y\in K$ 有 $|x-y|>a$, 故核为零; 对几乎处处的 $y\notin K$ 则 $f(y)=0$. 因而 $f_a(x)=0$. 由于右侧的和集紧, 支集结论成立.
\end{proof}

\begin{remark}
原讲 21 开头把一般平滑化写成紧支集光滑函数, 并在支集证明末尾把不等号方向写反. 一般 $f$ 只能保证光滑; 例如 $f\equiv1$ 时 $f_a\equiv1$, 支集是整个空间. 紧支集结论需要 $f$ 本身具有紧支集.
\end{remark}

\originalsection{近似恒等族与稠密性}

\begin{definition}[近似恒等族]
一族 $k_a\in L^1(\R^n)$, $a>0$, 称为近似恒等族, 若 $\int k_a=1$, $\sup_a\norm{k_a}_1=A<\infty$, 且对每个 $\delta>0$,
$\int_{|y|\ge\delta}|k_a(y)|\,\dd y\to0$ 当 $a\to0$.
核可以有符号; 若均非负, 则 $A=1$. 上面构造的 $\phi_a$ 是其中一个例子.
\end{definition}

\begin{theorem}[$L^p$ 逼近]\label{ra:c7:approx}
设 $1\le p<\infty$, $f\in L^p(\R^n)$, $(k_a)$ 为近似恒等族. 则 $\norm{f*k_a-f}_p\to0$ 当 $a\to0$. 特别地, $\norm{f*\phi_a-f}_p\to0$.
\end{theorem}
\begin{proof}
Young 不等式保证 $f*k_a$ 几乎处处存在. 由 $\int k_a=1$, 差写为
$f*k_a(x)-f(x)=\int k_a(y)(f(x-y)-f(x))\,\dd y$.
记 $A_a=\norm{k_a}_1$, 有 $1\le A_a\le A$. 对概率测度 $|k_a(y)|\,\dd y/A_a$ 用 Jensen, 当 $p=1$ 时只用积分三角不等式, 得
\[
 |f*k_a(x)-f(x)|^p
 \le A_a^{p-1}\int|k_a(y)|\,|f(x-y)-f(x)|^p\,\dd y.
\]
这也可先对截断后的非负积分使用 Jensen, 再取极限. Tonelli 将它积分为
$\norm{f*k_a-f}_p^p
\le A^{p-1}\int|k_a(y)|\,\norm{\tau_yf-f}_p^p\,\dd y$.

给定 $\eps>0$, 平移连续性给出 $\delta>0$, 使 $|y|<\delta$ 时 $\norm{\tau_yf-f}_p<\eps$. 对其余 $y$, 三角不等式及平移等距性给出 $\norm{\tau_yf-f}_p\le2\norm{f}_p$. 因此
\[
 \norm{f*k_a-f}_p^p
 \le A^p\eps^p+
 A^{p-1}(2\norm{f}_p)^p\int_{|y|\ge\delta}|k_a(y)|\,\dd y.
\]
先令 $a\to0$, 再令 $\eps\to0$, 得到结论. 这里绝对值不能删掉: 有符号核或有符号 $f$ 的差不提供非负控制.
\end{proof}

\begin{proposition}[一致逼近]\label{ra:c7:uniform}
若 $u$ 有界且一致连续, 则 $\norm{u*k_a-u}_\infty\to0$. 特别地, 若 $u\in C_c$, 则 $u*\phi_a\to u$ 一致收敛, 且当 $a\le1$ 时这些函数的支集都包含在某个固定紧集内.
\end{proposition}
\begin{proof}
把 $|u(x-y)-u(x)|$ 在 $|y|<\delta$ 上用一致连续性估计为 $\eps$, 在其余部分估计为 $2\norm{u}_\infty$. 积分后得到与 $x$ 无关的上界
$A\eps+2\norm{u}_\infty\int_{|y|\ge\delta}|k_a(y)|\,\dd y$.
先令 $a\to0$, 再令 $\eps\to0$. 支集结论来自定理 \ref{ra:c7:mollify}.
\end{proof}

\begin{theorem}[光滑紧支集函数稠密]\label{ra:c7:density}
对 $1\le p<\infty$, $C_c^\infty(\R^n)$ 在 $L^p(\R^n)$ 中稠密.
\end{theorem}
\begin{proof}
给定 $f\in L^p$ 和 $\eps>0$, 先由 $C_c$ 的稠密性选 $u\in C_c$, 使 $\norm{f-u}_p<\eps/2$. 对 $a\le1$, $u*\phi_a-u$ 的支集落在固定有限测度集 $K$ 中, 因而
$\norm{u*\phi_a-u}_p\le\lambda(K)^{1/p}\norm{u*\phi_a-u}_\infty\to0$.
取足够小的 $a$ 使它小于 $\eps/2$. 定理 \ref{ra:c7:mollify} 保证 $u*\phi_a\in C_c^\infty$, 三角不等式则给出所需逼近.

也可先截断 $f$ 的支集, 再用定理 \ref{ra:c7:approx} 平滑; 这样直接得到有紧支集的平滑逼近. 以上证明保留了原讲义先逼近连续函数、再平滑连续函数的路线, 并覆盖全部有限的 $p$.
\end{proof}

\begin{example}
作为编者补充, 检查 $p=\infty$ 的障碍. 在 $\R$ 上取 $f=\ind_{(0,\infty)}$. 对每个连续函数 $u$, 都有 $\norm{u-f}_\infty\ge1/2$.
\end{example}
\begin{solution}
若本质上界小于 $1/2$, 可取 $d<1/2$ 使 $|u-f|\le d$ 几乎处处. 连续性将 $|u|\le d$ 延伸到整个 $(-\infty,0)$, 将 $|u-1|\le d$ 延伸到整个 $(0,\infty)$: 若某点不满足, 邻域内会有正测度的反例. 令两侧趋于零, 得 $|u(0)|\le d$, $|u(0)-1|\le d$, 与 $1\le2d<1$ 矛盾. 所以不能把光滑稠密性或一般的 $L^p$ 平移连续性直接延伸到 $p=\infty$.
\end{solution}

这些逼近结果是实变与泛函分析的衔接点. 在本册中, 它们用于把积分和局部平均的性质从连续函数传给可积函数. 更抽象的算子、弱拓扑和分布理论属于后续课程; 这里仅建立所需的 $L^p$ 估计与光滑逼近.

\originalsection{习题与解答}

\begin{exercise}
设 $f\in L^p(\R^n)$, $1\le p<\infty$, 且 $f=0$ 几乎处处于 $B(0,R)^c$. 证明 $f*\phi_a\in C_c^\infty$, $\supp(f*\phi_a)\subset\overline{B(0,R+a)}$, 并给出其 $L^p$ 范数和导数范数的估计.
\end{exercise}
\begin{solution}
有限测度上的 Hölder 说明 $f$ 在该球上可积, 所以满足平滑化定理的条件. 支集包含由 $|x|\le|y|+|x-y|\le R+a$ 得到. Young 不等式给出 $\norm{f*\phi_a}_p\le\norm{f}_p$, 以及
$\norm{\partial^\alpha(f*\phi_a)}_p
\le\norm{f}_p\norm{\partial^\alpha\phi_a}_1
=a^{-|\alpha|}\norm{f}_p\norm{\partial^\alpha\phi}_1$.
后一个伸缩等式来自 $\partial^\alpha\phi_a(x)=a^{-n-|\alpha|}(\partial^\alpha\phi)(x/a)$.
\end{solution}

\begin{exercise}
设 $k\in L^1(\R^n)$ 且 $\int k=1$. 证明 $k_a(x)=a^{-n}k(x/a)$ 是近似恒等族, 并说明是否必须要求 $k\ge0$ 或 $k$ 紧支集.
\end{exercise}
\begin{solution}
伸缩代换给出 $\int k_a=1$, $\norm{k_a}_1=\norm{k}_1$. 对 $\delta>0$,
$\int_{|x|\ge\delta}|k_a(x)|\,\dd x=\int_{|u|\ge\delta/a}|k(u)|\,\dd u\to0$,
因为 $k$ 绝对可积. 这些条件不要求非负或紧支集. 非负时范数恰为 $1$; 紧支集时尾部最终直接为零. 两者都只是便于估计的附加性质.
\end{solution}

\begin{exercise}
设 $f\in L^1$, $g\in L^\infty$. 证明 $f*g$ 一致连续, 并举出它不在 $C_0$ 的例子.
\end{exercise}
\begin{solution}
换元将卷积写成 $\int f(x-y)g(y)\,\dd y$, 所以其平移差的上界为 $\norm{\tau_hf-f}_1\norm{g}_\infty$, 与 $x$ 无关且随 $h\to0$ 趋于零. 取 $f=\ind_{[0,1]^n}$ 和 $g\equiv1$, 得 $f*g\equiv1$, 在无穷远不消失.
\end{solution}

\originalchapter{最大函数与 Lebesgue 微分}
\label{ra:c8}
\sourcelecture{22--24}

\originalsection{积分函数的局部变化}

若 $f$ 连续, $F(x)=\int_a^xf(t)\,\dd t$ 的导数是 $f(x)$. 对 $f\in L^1$, 差商仍是一个短区间上的平均值, 但不能再用逐点连续性. 本章要证明: 对几乎处处的 $x$, 这些平均值仍趋于 $f(x)$. 关键工具是最大函数, 它控制所有尺度上的平均值.

先说明积分函数本身具有的正则性. 对 $f\in L^1([a,b])$, 集合 $E$ 测度很小时, $\int_E|f|$ 也很小. 事实上, 取 $M$ 足够大, 使 $\int_{\{|f|>M\}}|f|<\eps/2$, 则
$\int_E|f|\le\eps/2+M\lambda(E)$; 当 $\lambda(E)<\eps/(2M)$ 时即小于 $\eps$. 这称为积分对测度的绝对连续性.

\begin{definition}
实值或复值函数 $F:[a,b]\to\R$ 或 $\C$ 称为绝对连续, 若对每个 $\eps>0$ 都存在 $\delta>0$, 使任意有限个两两不交的区间 $(\alpha_j,\beta_j)\subset[a,b]$ 满足 $\sum_j(\beta_j-\alpha_j)<\delta$ 时, 有 $\sum_j|F(\beta_j)-F(\alpha_j)|<\eps$. 其中绝对值在复值情形表示复数的模.
\end{definition}

\begin{proposition}\label{ra:c8:primitive-ac}
若 $f\in L^1([a,b])$, 则 $F(x)=\int_a^xf(t)\,\dd t$ 绝对连续, 因而连续.
\end{proposition}
\begin{proof}
对上述不交区间, $\sum_j|F(\beta_j)-F(\alpha_j)|\le\int_{\bigcup_j(\alpha_j,\beta_j)}|f|$. 右边在区间总长度趋于零时趋于零, 正是刚才证明的积分绝对连续性. 只取一个区间就得到连续性.
\end{proof}

连续性之外的可导结论将在后面证明. 反过来, “几乎处处可导且导数可积”本身还不足以用积分恢复函数.

\begin{example}\label{ra:c8:cantor}
作为编者补充, 构造 Cantor 函数 $C:[0,1]\to[0,1]$: 它连续非减, $C(0)=0$, $C(1)=1$, 但 $C'=0$ 几乎处处. 所以 $C(1)-C(0)\ne\int_0^1C'$.
\end{example}
\begin{solution}
从三分 Cantor 集的构造出发. 第 $k$ 步留下 $2^k$ 个长度为 $3^{-k}$ 的闭区间. 定义连续非减的分段线性函数 $C_k$: 在按从左至右排序的每个留下区间上增加 $2^{-k}$, 在已经删去的开区间上保持常数, 且端点值为 $0,1$. 细分到下一步只在这些留下区间内改变函数, 并保留其端点值. 因而 $\norm{C_{k+1}-C_k}_\infty\le2^{-k}$, 该估计的尾和趋于零, 所以 $C_k$ 一致收敛到连续非减函数 $C$.

对任一在有限步骤中删去的开区间, 后续所有函数都在其上等于同一个常数, 故 $C'=0$ 于该开区间. 这些区间的并是 Cantor 集的补集, 其测度为 $1$, 所以 $C'=0$ 几乎处处. 端点值则在极限下仍为 $0,1$.

此函数不绝对连续. 第 $k$ 步留下区间的总长度为 $(2/3)^k\to0$, 但 $C$ 在各区间两端的增量之和为 $2^k2^{-k}=1$. 这直接违反绝对连续的定义.
\end{solution}

\originalsection{覆盖选择与最大函数}

\begin{lemma}[Vitali 覆盖选择]\label{ra:c8:vitali}
设 $E\subset\R^n$ 有界, $\mathcal F$ 是一族开球, 每个球的中心属于 $E$, 且每个 $x\in E$ 是其中至少一个球的中心. 则可从 $\mathcal F$ 中选出有限或可数个两两不交的球 $B_j$, 使 $E\subset\bigcup_j3B_j$. 这里 $3B_j$ 与 $B_j$ 同心, 半径是其三倍.
\end{lemma}
\begin{proof}
若球的半径无上界, 由于所有中心在有界集 $E$ 中, 一个足够大的球就覆盖 $E$, 结论成立. 因而可设半径有统一上界 $R$. 已选 $B_1,\ldots,B_{k-1}$ 后, 令 $d_k$ 为所有与它们均不相交的候选球半径的上确界. 若无此候选球, 就停止; 否则选 $B_k$ 使其半径 $r_k>d_k/2$. 此选择无需假定上确界能取到.

选出的球两两不交. 若选择无限进行, 所有球都位于一个固定有界球中, 所以 $\sum_kr_k^n<\infty$, 特别地 $r_k\to0$.
固定 $x\in E$, 取以它为中心的候选球 $B$ 并记其半径为 $\rho>0$. $B$ 必与某个已选球相交: 若与所有已选球都不相交, 在有限停止时会与停止条件矛盾; 在无限选择时则每一步 $d_k\ge\rho$, 导致 $r_k>\rho/2$, 与 $r_k\to0$ 矛盾.

设 $k$ 是第一次相交的编号. 在第 $k$ 步 $B$ 仍是候选球, 所以 $\rho\le d_k<2r_k$. 若 $z_k$ 为 $B_k$ 的中心, 取交点 $w\in B\cap B_k$, 得
$|x-z_k|\le|x-w|+|w-z_k|<\rho+r_k<3r_k$.
因此 $x\in3B_k$, 覆盖结论成立.
\end{proof}

\begin{remark}
本引理覆盖的是原球的中心集合 $E$. 因而选择半径大于候选上确界一半时, 放大倍数 $3$ 足够. 不能把结论未经检查改为“全部原球都包含在这些三倍球之并中”. 本章只需要中心集合的覆盖.
\end{remark}

\begin{definition}[Hardy--Littlewood 最大函数]
对 $f\in L^1_{\mathrm{loc}}(\R^n)$, 定义
$Mf(x)=\sup_{r>0}\lambda(B(x,r))^{-1}\int_{B(x,r)}|f(y)|\,\dd y$.
它可以取值 $+\infty$.
\end{definition}

\begin{proposition}\label{ra:c8:max-meas}
$Mf$ 下半连续, 因而 Borel 可测. 此外, $M(f+g)\le Mf+Mg$, $M(cf)=|c|Mf$, 且对 $f\in L^\infty$, $\norm{Mf}_\infty\le\norm{f}_\infty$.
\end{proposition}
\begin{proof}
若 $Mf(x)>t$, 存在 $r>0$ 使 $\int_{B(x,r)}|f|>t\lambda(B(x,r))$. 当 $t\ge0$ 时, 可取稍大的 $r'>r$, 使同一个积分仍大于 $t\lambda(B(x,r'))$. 若 $|x'-x|<r'-r$, 则 $B(x,r)\subset B(x',r')$, 且两个半径为 $r'$ 的球体积相同. 故 $Mf(x')>t$. 当 $t<0$ 时超水平集是全空间. 因而 $\{Mf>t\}$ 开, 即下半连续. 其余三个性质来自每个球上的积分三角不等式、齐次性和本质上界.
\end{proof}

\begin{definition}[弱 $L^1$ 估计]
可测函数 $g$ 满足弱 $L^1$ 估计, 是指存在 $A<\infty$, 使 $\lambda\{|g|>t\}\le A/t$ 对所有 $t>0$ 成立. 若算子 $T$ 满足 $A\le C\norm{f}_1$ 对 $g=Tf$ 成立, 就说它满足弱 $(1,1)$ 估计.
\end{definition}

Chebyshev 不等式说明 $L^1$ 函数满足此估计, 反向却不成立. 例如令 $g(x)=|x|^{-n}$ 于 $x\ne0$, 并在原点任意赋值, 则 $\lambda\{g>t\}=v_n/t$, 其中 $v_n=\lambda(B(0,1))$. 它在原点附近和无穷远都不可积.

\begin{theorem}[最大函数的弱 $(1,1)$ 估计]\label{ra:c8:weak}
若 $f\in L^1(\R^n)$, 则对每个 $t>0$,
$\lambda\{Mf>t\}\le3^n\norm{f}_1/t$.
\end{theorem}
\begin{proof}
先对有界集 $E_k=\{Mf>t\}\cap B(0,k)$ 估计. 对每个 $x\in E_k$, 选择一个以 $x$ 为中心的球 $B$, 使
$t\lambda(B)<\int_B|f|$.
这些球的半径有统一上界, 因为 $t\lambda(B)<\norm{f}_1$. 用引理 \ref{ra:c8:vitali} 选出两两不交的 $B_j$, 使 $E_k\subset\bigcup_j3B_j$. 于是
\[
 \lambda(E_k)\le3^n\sum_j\lambda(B_j)
 \le\frac{3^n}{t}\sum_j\int_{B_j}|f|
 \le\frac{3^n}{t}\norm{f}_1.
\]
最后 $E_k\uparrow\{Mf>t\}$, 用测度的下连续性即可. 若 $\norm{f}_1=0$, 每个球上的积分都为零, 结论也直接成立.
\end{proof}

同一估计也对 $\{Mf\ge t\}$ 成立: 对 $s<t$ 使用 $\{Mf\ge t\}\subset\{Mf>s\}$, 再令 $s\uparrow t$. 特别地, $Mf<\infty$ 几乎处处.

\begin{proposition}
若 $f\in L^1(\R^n)$ 且 $f$ 不几乎处处为零, 则 $Mf\notin L^1(\R^n)$.
\end{proposition}
\begin{proof}
选择 $a>0$ 使 $c=\int_{B(0,a)}|f|>0$. 当 $|x|>a$ 时, $B(0,a)\subset B(x,2|x|)$, 所以
$Mf(x)\ge c/(v_n2^n|x|^n)$.
函数 $|x|^{-n}$ 在无穷远不可积. 不用极坐标也能看出: 在每个环带 $2^ka<|x|<2^{k+1}a$ 上, 其积分至少为
$(2^{k+1}a)^{-n}v_n((2^{k+1}a)^n-(2^ka)^n)=v_n(1-2^{-n})$.
将无穷多个不交环带相加即发散. 故 $Mf$ 不可积.
\end{proof}

\begin{remark}
原讲 22 的证明中把“若 $Mf$ 可积”误写成“若 $Mf$ 不可积”. 正确的推论是 $Mf\in L^1$ 迫使 $f=0$ 几乎处处. 因此弱型估计不能直接提升到 $L^1$ 范数估计.
\end{remark}

\originalsection{Lebesgue 微分定理}

连续函数在小球上的平均振荡趋于零. 可积函数可以由连续函数在 $L^1$ 中逼近, 但 $L^1$ 小误差不能逐点控制. 最大函数的弱型估计正好控制误差平均值较大的点集.

\begin{theorem}[Lebesgue 微分定理]\label{ra:c8:differentiation}
若 $f\in L^1_{\mathrm{loc}}(\R^n)$, 则对几乎处处的 $x$,
\[
 \lim_{r\downarrow0}\frac1{\lambda(B(x,r))}
 \int_{B(x,r)}|f(y)-f(x)|\,\dd y=0.
\]
因而 $\lambda(B(x,r))^{-1}\int_{B(x,r)}f(y)\,\dd y\to f(x)$ 几乎处处.
\end{theorem}
\begin{proof}
先设 $f\in L^1(\R^n)$, 并在一个零测集上修改, 使其处处有限. 定义平均振荡的上极限
\[
 f^*(x)=\limsup_{r\downarrow0}\frac1{\lambda(B(x,r))}
 \int_{B(x,r)}|f(y)-f(x)|\,\dd y.
\]
目标是证明 $f^*=0$ 几乎处处. 下面的估计用外测度, 因而无需先证明这个上极限可测.

积分三角不等式给出 $(f+g)^*\le f^*+g^*$. 若 $g$ 在 $x$ 连续, 则对每个 $\eps>0$, 所有足够小的球上 $|g(y)-g(x)|<\eps$, 因而 $g^*(x)=0$. 将三角不等式分别用于 $f=(f-g)+g$ 和 $f-g=f+(-g)$, 得当 $g$ 连续时 $(f-g)^*=f^*$.

另一方面, 对每个球有平均振荡不超过 $Mf(x)+|f(x)|$, 所以 $f^*\le Mf+|f|$. 因而对 $t>0$,
$\{f^*>t\}\subset\{Mf>t/2\}\cup\{|f|>t/2\}$.
用最大函数的弱型估计和 Chebyshev 不等式, 得
\[
 \lambda^*\{f^*>t\}\le\frac{2(3^n+1)}{t}\norm{f}_1.
\]
对任意 $\eps>0$, 由 $C_c$ 在 $L^1$ 中稠密, 选连续函数 $g$ 使 $\norm{f-g}_1<\eps$. 对 $f-g$ 使用刚才的估计并利用 $(f-g)^*=f^*$, 得
$\lambda^*\{f^*>t\}\le2(3^n+1)\eps/t$.
让 $\eps\downarrow0$, 得每个超水平集的外测度为零. 又
$\{f^*>0\}=\bigcup_{k\ge1}\{f^*>1/k\}$,
故外测度为零. 由于平均振荡非负, 上极限为零就意味着其极限为零.

一般局部可积的情形, 令 $f_j=f\ind_{B(0,j)}\in L^1$. 在 $B(0,j)$ 的每个点, 足够小的球完全位于该球中, 故 $f_j$ 与 $f$ 的平均振荡相同. 每个 $j$ 的例外集为零测集, 去掉它们的可数并即可.

最后, 平均值与 $f(x)$ 的差的绝对值不超过平均振荡, 从而得到第二个结论.
\end{proof}

\begin{definition}[Lebesgue 点与精确值]\label{ra:c8:lebpoint}
对局部可积等价类 $[f]$, 若存在数 $A$ 使
$\lambda(B(x,r))^{-1}\int_{B(x,r)}|f(y)-A|\,\dd y\to0$,
则称 $x$ 为该等价类的 Lebesgue 点, 称 $A$ 为在 $x$ 处的精确值, 记为 $\widetilde f(x)$.
\end{definition}

精确值唯一: 若 $A,B$ 都满足定义, 则对每个球有
$|A-B|\le\lambda(B(x,r))^{-1}\int(|f-A|+|f-B|)$,
令 $r\downarrow0$ 就得 $A=B$. 零测修改不改变任何这些积分, 所以 Lebesgue 点集合和精确值只取决于等价类. 微分定理说明几乎处处的点都是 Lebesgue 点, 并且对任一代表 $f$, $\widetilde f(x)=f(x)$ 几乎处处. 这不保证在每个 Lebesgue 点都等于任意选取的代表值.

\begin{example}
函数 $\ind_{\mathbb Q}$ 与零函数属于同一个局部可积等价类. 每个点都是 Lebesgue 点, 精确值都是 $0$, 但在有理点处所选代表值是 $1$.
\end{example}

\begin{example}
原讲 23 的振荡例子: 在 $\R$ 上令 $g(x)=\sin(1/x)$ 当 $x\ne0$, $g(0)=0$. 原点不是其 Lebesgue 点.
\end{example}
\begin{solution}
$g$ 是奇函数, 所以每个关于原点对称的区间上的平均值为零. 若存在精确值 $A$, 平均绝对偏差趋于零必推出平均值趋于 $A$, 因而只能有 $A=0$.

但平均绝对值不趋于零. 换元 $u=1/x$ 得
$r^{-1}\int_0^r|\sin(1/x)|\,\dd x=r^{-1}\int_{1/r}^\infty|\sin u|u^{-2}\,\dd u$.
记 $q(u)=|\sin u|-2/\pi$. 它是周期函数, 每个周期的积分为零, 所以有有界原函数 $Q$. 分部积分给出
$\int_A^\infty q(u)u^{-2}\,\dd u
=-Q(A)A^{-2}+2\int_A^\infty Q(u)u^{-3}\,\dd u=O(A^{-2})$.
于是上面的平均绝对值趋于 $2/\pi$. 在 $(-r,r)$ 上的平均绝对值相同, 故 $A=0$ 也不满足定义.
\end{solution}

\begin{definition}[正则地趋近一点]
可测集序列 $E_k$ 称为正则地趋近 $x$, 若存在 $c<\infty$ 和正数序列 $r_k\to0$, 使 $0<\lambda(E_k)$, $E_k\subset B(x,r_k)$, 且 $\lambda(B(x,r_k))\le c\lambda(E_k)$.
\end{definition}

\begin{theorem}\label{ra:c8:regular}
若 $x$ 是 $f\in L^1_{\mathrm{loc}}(\R^n)$ 的 Lebesgue 点, $E_k$ 正则地趋近 $x$, 则 $\lambda(E_k)^{-1}\int_{E_k}f\to\widetilde f(x)$.
\end{theorem}
\begin{proof}
设 $A=\widetilde f(x)$. 则
\[
 \left|\frac1{\lambda(E_k)}\int_{E_k}f-A\right|
 \le\frac1{\lambda(E_k)}\int_{E_k}|f-A|
 \le\frac c{\lambda(B(x,r_k))}\int_{B(x,r_k)}|f-A|\longrightarrow0.
\]
\end{proof}

这里的体积比条件防止 $E_k$ 在球内过于狭细, 从而只挑出平均意义下很小的一部分异常区域.

\begin{corollary}[Lebesgue 密度定理]\label{ra:c8:density}
若 $E\subset\R^n$ Lebesgue 可测, 则对几乎处处的 $x\in E$,
$\lambda(E\cap B(x,r))/\lambda(B(x,r))\to1$;
对几乎处处的 $x\notin E$, 此比值趋于 $0$.
\end{corollary}
\begin{proof}
$\ind_E$ 局部可积, 对它使用微分定理即可. 不要求 $E$ 的总测度有限.
\end{proof}

\originalsection{微积分基本定理}

\begin{theorem}[Lebesgue 积分的微积分基本定理 I]\label{ra:c8:ftc}
设 $f\in L^1([a,b])$, $F(x)=\int_a^xf(t)\,\dd t$. 则 $F$ 绝对连续, 在几乎处处的 $x\in(a,b)$ 可导, 且 $F'(x)=f(x)$ 几乎处处. 更精确地, 若 $x$ 是 $f$ 的 Lebesgue 点, 则 $F'(x)=\widetilde f(x)$.
\end{theorem}
\begin{proof}
绝对连续性已由命题 \ref{ra:c8:primitive-ac} 证明. 将 $f$ 在 $[a,b]$ 外置零. 若 $h>0$ 且 $x+h\le b$, 则
$[F(x+h)-F(x)]/h=h^{-1}\int_x^{x+h}f(t)\,\dd t$.
区间 $(x,x+h)$ 包含在 $B(x,h)$ 中, 后者长度为 $2h$, 所以它正则地趋近 $x$, 体积比常数为 $2$. 定理 \ref{ra:c8:regular} 说明, 对任意 $h_k\downarrow0$, 右侧趋于 $\widetilde f(x)$. 任意趋于零的正数序列都可同样处理, 所以右导数存在并等于该值.

当 $h<0$ 时, 差商是 $(x+h,x)$ 上的平均值, 同理得左导数为同一个精确值. 因而在每个内部 Lebesgue 点 $F$ 可导. 微分定理说明这些点几乎处处存在, 且精确值与代表 $f(x)$ 几乎处处相等.
\end{proof}

\begin{theorem}[绝对连续版微积分基本定理 II]\label{ra:c8:ac-inverse}
设 $F:[a,b]\to\R$ 绝对连续. 则 $F$ 几乎处处可导, $F'\in L^1([a,b])$, 且对每个 $x\in[a,b]$,
$F(x)-F(a)=\int_a^xF'(t)\,\dd t$.
因此绝对连续函数恰是可积函数的不定积分加一个常数.
\end{theorem}

原讲 22 陈述了这一定理, 没有给出其证明. 本册保留其完整结论, 并在定理 \ref{ra:c9:ac-ftc} 中用正测度、Radon--Nikodym 密度和本章微分定理补全证明. 最后一条等价刻画的反向, 已由命题 \ref{ra:c8:primitive-ac} 证明. 例题 \ref{ra:c8:cantor} 说明绝对连续性不能只换成连续性和几乎处处可导.

\begin{proposition}[处处可导时的充分条件]\label{ra:c8:everywhere}
若 $F$ 在 $[a,b]$ 上连续, 在每个 $x\in(a,b)$ 可导, 且 $F'\in L^1([a,b])$, 则对每个 $x\in[a,b]$, 有
$F(x)-F(a)=\int_a^xF'(t)\,\dd t$.
\end{proposition}

这是原讲 22 同时指出的另一个充分条件. 几乎处处可导不够, 但处处可导与导数可积结合则足够. 原稿没有给出这一结论的证明; 本册在定理 \ref{ra:c9:everywhere-ftc} 中用 Vitali--Carathéodory 的逐点上逼近和单侧 Dini 导数补证, 不把它当作前述微分定理的前提.

\begin{remark}
原讲 23 把积分函数几乎处处求导的结论标为 \textit{FTOC II}, 而课程目录将它列为微积分基本定理 I. 本册按通常的内容区分命名: I 是对积分函数求导, II 是由导数积分恢复绝对连续函数.
\end{remark}

\originalsection{积分形式的 Minkowski 不等式}

\sourcelecture{24}
下面回到 $L^p$ 范数. 有限和的 Minkowski 不等式可以推广为一族函数的积分, 但必须先证明这个积分对几乎处处的点存在.

\begin{theorem}[广义 Minkowski 不等式]\label{ra:c8:gmink}
设 $(X,\mathcal M,\mu)$, $(Y,\mathcal N,\nu)$ 都 $\sigma$-有限, $1\le p<\infty$, $u:X\times Y\to\C$ 乘积可测. 令 $a(y)=\norm{u_y}_{L^p(\mu)}$, 并设 $C=\int_Ya(y)\,\dd\nu(y)<\infty$. 则对几乎处处的 $x$, $u^x\in L^1(\nu)$, 函数 $U(x)=\int_Yu(x,y)\,\dd\nu(y)$ 可测, 且
\[
 \norm{U}_{L^p(\mu)}
 \le\int_Y\norm{u_y}_{L^p(\mu)}\,\dd\nu(y).
\]
\end{theorem}
\begin{proof}
由 Tonelli, $a(y)^p=\int_X|u(x,y)|^p\,\dd\mu(x)$ 可测. 若 $p=1$, 对 $|u|$ 用 Tonelli, 得 $\int_X\int_Y|u|=C<\infty$, 所以截面绝对可积且不等式成立.

设 $p>1$. 若 $C=0$, 则 $a=0$ 几乎处处, 从 Tonelli 可知 $u=0$ 乘积几乎处处, 结论直接成立. 以下设 $C>0$. 可在 $u$ 上作不影响结论的零测修改, 使 $a(y)=0$ 或 $a(y)=\infty$ 的截面为零. 对 $a=\infty$ 的参数集合, $\nu$ 测度为零, 它与 $X$ 的乘积零测, 这里使用 $X$ 的 $\sigma$-有限性. 对 $a=0$ 的参数, $u_y=0$ $\mu$-几乎处处; 对这些截面非零处的指示函数用 Tonelli, 可知删去的部分同样是乘积零测.

在 $0<a(y)<\infty$ 处定义 $v(x,y)=|u(x,y)|a(y)^{-1/p'}$, 其余置零. 则对几乎处处的 $y$,
$\int_Xv(x,y)^p\,\dd\mu(x)=a(y)$.
对几乎处处的 $x$, 关于 $y$ 的 Hölder 不等式给出
\[
 H(x):=\int_Y|u(x,y)|\,\dd\nu(y)
 \le\left(\int_Yv(x,y)^p\,\dd\nu(y)\right)^{1/p}C^{1/p'}.
\]
用 Tonelli 将右侧的 $p$ 次幂积分, 得
$\int_XH(x)^p\,\dd\mu(x)\le C^{p-1}\int_Y(\int_Xv^p\,\dd\mu)\dd\nu=C^p$.
因此 $H<\infty$ 几乎处处, 从而 $U$ 绝对存在. 对实部、虚部的正负部分作截面积分, 可知 $U$ 可测; 又 $|U|\le H$, 给出所需估计. 将零测修改还原, 对修改集使用 Tonelli 可知积分和结论仍几乎处处不变.
\end{proof}

\begin{example}
重新证明一个 Young 不等式. 若 $f\in L^1(\R^n)$, $g\in L^p(\R^n)$, $1\le p<\infty$, 令 $u(x,y)=f(y)g(x-y)$. 对固定 $y$, 平移不变性给出 $\norm{u_y}_p=|f(y)|\norm{g}_p$. 因此广义 Minkowski 不等式给出 $\norm{f*g}_p\le\norm{f}_1\norm{g}_p$, 并同时证明卷积几乎处处绝对存在. 当 $Y$ 是有限集合并取计数测度时, 同一定理退化为普通有限和的 Minkowski 不等式.
\end{example}

\originalsection{分布函数与最大算子}

\begin{definition}[分布函数]
对测度空间上的可测函数 $f$, 定义 $d_f(t)=\mu\{x:|f(x)|>t\}$, $t\ge0$. 它是非增函数, 记录函数超过每个高度的区域大小.
\end{definition}

\begin{theorem}[分层积分公式]\label{ra:c8:layercake}
若 $f:X\to[0,\infty]$ 可测, 则 $\int_Xf\,\dd\mu=\int_0^\infty\mu\{f>t\}\,\dd t$. 对 $p>0$,
\[
 \int_X|f|^p\,\dd\mu
 =p\int_0^\infty d_f(t)t^{p-1}\,\dd t.
\]
这些等式允许取无穷值.
\end{theorem}
\begin{proof}
若 $(X,\mu)$ $\sigma$-有限, 将 $f(x)$ 写成 $\int_0^\infty\ind_{\{0<t<f(x)\}}\,\dd t$, 直接用 Tonelli 即得第一式. 为说明此公式其实不需要 $\sigma$-有限性, 也可以从非负简单函数出发: 若 $s=\sum_ja_j\ind_{E_j}$, $E_j$ 两两不交, 则
$\int_0^\infty\mu\{s>t\}\,\dd t=\sum_j\mu(E_j)\int_0^{a_j}\dd t=\int s$.
选择简单函数 $s_k\uparrow f$. 对每个 $t$, $\{s_k>t\}\uparrow\{f>t\}$, 两边分别用单调收敛, 得第一式.

对 $|f|^p$ 使用第一式, 得 $\int|f|^p=\int_0^\infty\mu\{|f|>s^{1/p}\}\,\dd s$. 代换 $s=t^p$, 即得第二式.
\end{proof}

更一般地, 若 $\Phi:[0,\infty)\to[0,\infty)$ 局部绝对连续, $\Phi(0)=0$, 且 $\Phi'\ge0$ 几乎处处, 则
$\int_X\Phi(f)\,\dd\mu=\int_0^\infty\mu\{f>t\}\Phi'(t)\,\dd t$.
在 $\sigma$-有限空间上, 这是对 $\Phi(f(x))=\int_0^{f(x)}\Phi'(t)\,\dd t$ 用 Tonelli 的直接结果. 一般空间可先对简单函数验证, 再单调逼近. 在 $f=\infty$ 时, 令 $\Phi(\infty)$ 为递增极限. 原讲仅写“$\Phi$ 可微”并不充分: 还需要处理 $\Phi(0)$ 以及有符号积分的存在性. 此处给出的是所有积分均非负、足以支撑本章应用的版本.

\begin{theorem}[最大函数的 $L^p$ 有界性]\label{ra:c8:strong}
设 $1<p<\infty$, $f\in L^p(\R^n)$, 则 $Mf\in L^p(\R^n)$, 且
$\norm{Mf}_p\le C(n,p)\norm{f}_p$.
可取 $C(n,p)=(3^np)^{1/p}p/(p-1)$. 这个上界当 $p\to\infty$ 时趋于 $1$, 当 $p\downarrow1$ 时趋于无穷. 对 $p=\infty$, 有 $\norm{Mf}_\infty\le\norm{f}_\infty$.
\end{theorem}
\begin{proof}
有限测度集上的 Hölder 说明 $f$ 局部可积, 所以最大函数定义有效. 因为 $Mf=M|f|$, 可设 $f\ge0$. 对 $t>0$ 以及固定 $0<c<1$, 分解
$f=g_t+h_t$, 其中 $g_t=f\ind_{\{f>ct\}}$, $h_t=f\ind_{\{f\le ct\}}$.
这里 $0\le h_t\le ct$, 而 $g_t\in L^1$, 因为在 $\{f>ct\}$ 上 $f\le(ct)^{1-p}f^p$. 这是使用弱 $(1,1)$ 估计前必须核对的条件.

最大函数的次可加性及 $L^\infty$ 估计给出 $Mf\le Mg_t+ct$. 所以
\[
 \lambda\{Mf>t\}
 \le\lambda\{Mg_t>(1-c)t\}
 \le\frac{3^n}{(1-c)t}\int_{\{f>ct\}}f(x)\,\dd x.
\]
将此代入分层积分公式, 再对非负积分使用 Tonelli, 得
\begin{align*}
 \norm{Mf}_p^p
 &\le\frac{3^np}{1-c}\int_0^\infty t^{p-2}
       \left(\int_{\{f>ct\}}f(x)\,\dd x\right)\dd t\\
 &=\frac{3^np}{1-c}\int_{\R^n}f(x)
       \left(\int_0^{f(x)/c}t^{p-2}\,\dd t\right)\dd x\\
 &=\frac{3^npc^{1-p}}{(1-c)(p-1)}\norm{f}_p^p.
\end{align*}
内层积分在零端有限恰因 $p>1$. 右边有限, 因此此计算同时证明 $Mf\in L^p$, 未预先假定其可积性.

为选择较好的常数, 对 $c^{1-p}/(1-c)$ 的对数求导, 得 $(1-p)/c+1/(1-c)=0$, 即 $c=(p-1)/p$. 代入给出
$\norm{Mf}_p^p\le3^np[p/(p-1)]^p\norm{f}_p^p$,
取 $p$ 次根得到所述常数. 无穷端点则已由命题 \ref{ra:c8:max-meas} 证明.
\end{proof}

\begin{remark}
上面的分解证明是原讲 24 所用的 Marcinkiewicz 插值方法: 对大于阈值的部分用弱 $L^1$ 控制, 对其余部分用 $L^\infty$ 控制, 再通过分布函数还原 $L^p$ 范数. 此处证明的是最大算子的这个具体结论, 没有把它当成已证明的一般插值定理. 常数在 $p\downarrow1$ 时发散与最大函数通常不在 $L^1$ 中一致, 但并不声称该上界是最优算子范数.
\end{remark}

\originalsection{习题与解答}

\begin{exercise}
设 $f\in L^1([a,b])$, $F(x)=\int_a^xf$. 证明 $|F(x)-F(y)|\le\int_{\min(x,y)}^{\max(x,y)}|f|$, 并在 $f\in L^p$, $1<p\le\infty$, 时给出连续性模的估计.
\end{exercise}
\begin{solution}
积分可加性给出第一式. 若 $1<p<\infty$, 对该区间用 Hölder, 得
$|F(x)-F(y)|\le\norm{f}_p|x-y|^{1-1/p}$.
若 $p=\infty$, 得 $|F(x)-F(y)|\le\norm{f}_\infty|x-y|$, 即 Lipschitz 连续. 仅有 $L^1$ 条件时, 第一式与积分绝对连续性仍保证 $F$ 连续, 但未提供统一的幂次模.
\end{solution}

\begin{exercise}
设 $x$ 为局部可积函数的 Lebesgue 点. 对任意包含 $x$ 的轴平行立方体 $Q_k$, 若其边长趋于零, 证明 $\lambda(Q_k)^{-1}\int_{Q_k}f\to\widetilde f(x)$.
\end{exercise}
\begin{solution}
若边长为 $\ell_k$, 则 $Q_k\subset\overline{B(x,\sqrt n\,\ell_k)}$. 可用稍大的球, 如半径 $2\sqrt n\,\ell_k$, 保证包含关系对边界也成立. 球的体积是 $v_n(2\sqrt n)^n\ell_k^n$, 立方体体积为 $\ell_k^n$, 比值有与 $k$ 无关的上界. 因而这些立方体正则地趋近 $x$, 由定理 \ref{ra:c8:regular} 得结论.
\end{solution}

\begin{exercise}
设 $f=\ind_{B(0,1)}$. 证明 $Mf$ 不可积, 但对每个 $p>1$ 有 $Mf\in L^p$. 说明这与弱 $(1,1)$ 估计的关系.
\end{exercise}
\begin{solution}
$f$ 是非零 $L^1$ 函数, 前述命题说明 $Mf\notin L^1$. 另一方面 $f\in L^p$ 对所有有限的 $p$ 成立, 所以定理 \ref{ra:c8:strong} 给出 $Mf\in L^p$ 当 $p>1$; 对 $p=\infty$ 也有 $Mf\le1$. 弱 $(1,1)$ 只控制超水平集的测度为 $O(1/t)$, 将这个上界在所有高度上直接积分不能得到有限的 $L^1$ 范数. 它允许此例的无穷远尾部.
\end{solution}

\begin{exercise}
设 $X$ 的测度有限, $f\ge0$, $0<p<q<\infty$, 且 $\int f^q<\infty$. 用分布函数公式直接证明 $f\in L^p$.
\end{exercise}
\begin{solution}
对 $0<t\le1$, 用 $d_f(t)\le\mu(X)$; 对 $t>1$, 由 Chebyshev, $d_f(t)\le t^{-q}\int f^q$. 因此
$\int f^p=p\int_0^\infty d_f(t)t^{p-1}\,\dd t
\le\mu(X)+[p/(q-p)]\int f^q<\infty$.
这是有限测度空间中 $L^q\subset L^p$ 的另一种证明; 它不声称给出 Hölder 所得的最精确范数常数.
\end{solution}

\originalchapter{积分密度与测度分解}
\label{ra:c9:chapter}
\noindent 本章为编者补充. Radon--Nikodym 定理列在原课程教学大纲中, 但不在已公开的24份讲义正文中. 以下证明只使用前文已经建立的积分理论和 $L^2$ 完备性; 所需的 Hilbert 空间表示引理也在这里证明.

\originalsection{绝对连续与奇异}
如果 $f\ge0$ 可测, 则 $\nu(E)=\int_E f\dd\mu$ 定义一个测度. 单调收敛定理保证可列可加性. 这种测度不会给 $\mu$ 的零集增加质量, 因而自然提出: 是否每个具有这一性质的测度都由一个密度表示?

\begin{definition}
设 $\mu,\nu$ 为同一可测空间 $(X,\mathcal A)$ 上的正测度. 若 $\mu(E)=0$ 总蕴涵 $\nu(E)=0$, 称 $\nu$ 对 $\mu$ 绝对连续, 记为 $\nu\ll\mu$. 若存在 $S\in\mathcal A$ 使 $\mu(S)=0$ 且 $\nu(X\setminus S)=0$, 称 $\nu$ 与 $\mu$ 互相奇异, 记为 $\nu\perp\mu$.
\end{definition}
这里的绝对连续首先是零集之间的关系, 不要求存在一个固定常数使 $\nu(E)\le C\mu(E)$. 奇异也不要求两者的拓扑支集不交. 原子测度可以集中在稠密的可数集合上, 而该集合仍为 Lebesgue 零集.

\begin{proposition}\label{ra:c9:uniform-ac}
若 $\nu$ 为有限正测度且 $\nu\ll\mu$, 则对每个 $\eps>0$ 都存在 $\delta>0$, 使 $\mu(E)<\delta$ 蕴涵 $\nu(E)<\eps$.
\end{proposition}
\begin{proof}
反设结论不成立. 存在 $\eps_0>0$ 和 $E_n\in\mathcal A$, 满足 $\mu(E_n)<2^{-n}$, 但 $\nu(E_n)\ge\eps_0$. 令 $F_k=\bigcup_{n\ge k}E_n$, 则 $F_k$ 递减且 $\mu(F_k)\le\sum_{n\ge k}2^{-n}\to0$. 所以 $F=\bigcap_kF_k$ 满足 $\mu(F)=0$. 绝对连续性给出 $\nu(F)=0$. 另一方面, 因 $\nu$ 有限, 从上连续性给出 $\nu(F)=\lim_k\nu(F_k)\ge\eps_0$, 矛盾.
\end{proof}

\begin{example}
在 $(0,1)$ 上取 $\mu$ 为 Lebesgue 测度, $\nu(E)=\int_E x^{-1/2}\dd x$. 说明绝对连续不蕴涵 $\nu\le C\mu$.
\end{example}
\begin{solution}
零集上的非负积分为零, 所以 $\nu\ll\mu$. 对 $E=(0,t)$, 有 $\nu(E)=2\sqrt t$, 而 $\mu(E)=t$. 两者之比为 $2/\sqrt t$, 当 $t\downarrow0$ 时无界. 尽管没有线性的统一常数, 命题\ref{ra:c9:uniform-ac}仍适用, 因为 $\nu((0,1))=2$.
\end{solution}

\originalsection{一个表示引理}
积分关于函数是线性的. 为从测度恢复密度, 我们把它看成 $L^2$ 上的线性泛函. 下面的证明说明表示函数为何存在, 不把一般泛函分析的表示定理作为未证黑箱.

\begin{lemma}\label{ra:c9:hilbert}
设 $H$ 为实 Hilbert 空间, 即具有内积 $\ip{\cdot}{\cdot}$ 且在相应范数下完备的实向量空间. 若 $T:H\to\R$ 为有界线性泛函, 即存在线性映射 $T$ 与常数 $C$ 使 $|T(h)|\le C\norm h$ 对所有 $h$ 成立, 则存在唯一 $g\in H$, 使所有 $h\in H$ 都满足 $T(h)=\ip{h}{g}$.
\end{lemma}
\begin{proof}
记 $\norm T=\sup_{\norm h\le1}|T(h)|$. 有界性蕴涵连续性, 因为 $|T(h)-T(k)|\le\norm T\norm{h-k}$. 若 $T=0$, 取 $g=0$. 以下设 $T\ne0$, 并令 $A=\{u:T(u)=1\}$. 该集合非空且闭. 由 $1=|T(u)|\le\norm T\norm u$, 数 $a=\inf_{u\in A}\norm u$ 为正. 选 $u_n\in A$ 使 $\norm{u_n}\to a$. 中点仍属于 $A$, 因而由平行四边形恒等式,
\[
 \norm{u_n-u_m}^2
 =2\norm{u_n}^2+2\norm{u_m}^2
   -4\norm{(u_n+u_m)/2}^2
 \le2\norm{u_n}^2+2\norm{u_m}^2-4a^2.
\]
右端趋于零, 所以 $u_n$ 是 Cauchy 列. 完备性给出极限 $u\in A$, 且 $\norm u=a$. 对任意 $v\in\ker T$ 和实数 $t$, 有 $u+tv\in A$. 展开 $\norm{u+tv}^2\ge\norm u^2$, 得 $2t\ip u v+t^2\norm v^2\ge0$. 让 $t$ 分别从正负两侧趋于零, 得 $\ip u v=0$.

任意 $h$ 都有 $h-T(h)u\in\ker T$, 因而 $\ip h u=T(h)\norm u^2$. 取 $g=u/\norm u^2$ 即得表示. 若 $g_1,g_2$ 都表示 $T$, 代入 $h=g_1-g_2$ 得 $\norm{g_1-g_2}^2=0$, 故唯一.
\end{proof}
在下面的使用中, $H=L^2(X,\lambda;\R)$, 内积为 $\ip h k=\int hk\dd\lambda$. 乘积可积由 Hölder 不等式保证, 完备性已在第5章建立. 这里只需要实空间, 不涉及复内积的共轭约定.

\originalsection{Radon--Nikodym 定理}
先把两种测度加起来. 相对于 $\lambda=\mu+\nu$, 二者都受到 $\lambda$ 控制, 因而可以用上一引理表示. 这个中间测度避免预先猜测密度的大小.

\begin{theorem}[Radon--Nikodym]\label{ra:c9:rn}
设 $\mu,\nu$ 是同一可测空间上的 $\sigma$ 有限正测度, 且 $\nu\ll\mu$. 则存在非负可测函数 $f$, 使每个 $E\in\mathcal A$ 满足 $\nu(E)=\int_E f\dd\mu$. 函数 $f$ 在 $\mu$ 几乎处处意义下唯一, 且可以选成几乎处处有限. 记作 $f=\dd\nu/\dd\mu$.
\end{theorem}
\begin{proof}
先设 $\mu(X),\nu(X)<\infty$. 令 $\lambda=\mu+\nu$. 对 $h\in L^2(\lambda;\R)$ 定义 $T(h)=\int h\dd\nu$. 由于 $\nu\le\lambda$, $h$ 的 $\lambda$ 几乎处处代表也确定 $\nu$ 几乎处处代表, 所以定义与代表选择无关. Cauchy--Schwarz 给出
\[
 |T(h)|\le\int |h|\dd\nu
 \le\nu(X)^{1/2}\left(\int |h|^2\dd\nu\right)^{1/2}
 \le\nu(X)^{1/2}\norm h_{L^2(\lambda)}.
\]
由引理\ref{ra:c9:hilbert}, 存在实可测 $g\in L^2(\lambda)$ 使 $T(h)=\int hg\dd\lambda$. 对所有可测 $E$ 代入 $h=\ind_E$, 得 $\nu(E)=\int_E g\dd\lambda$, 以及 $\mu(E)=\int_E(1-g)\dd\lambda$.

正测度不能在某个集合上得到负积分. 对集合 $\{g\le-1/n\}$ 和 $\{g\ge1+1/n\}$ 分别使用上述两个等式, 得这些集合的 $\lambda$ 测度为零. 改变一个零集上的取值后, 可令 $0\le g\le1$ 处处成立. 令 $S=\{g=1\}$. 则 $\mu(S)=0$, 又因 $\nu\ll\mu$, 有 $\nu(S)=0$, 所以 $\lambda(S)=0$. 在 $X\setminus S$ 上令 $f=g/(1-g)$, 在 $S$ 上令 $f=0$.

对任意非负可测 $u$, 从指示函数的等式、简单函数线性和单调逼近依次推出
\[
 \int u\dd\mu=\int u(1-g)\dd\lambda,
 \qquad
 \int u\dd\nu=\int ug\dd\lambda.
\]
取 $u=f\ind_E$ 并注意 $S$ 为 $\lambda$ 零集, 就得到 $\int_E f\dd\mu=\int_E g\dd\lambda=\nu(E)$. 有限情形成立.

一般情形中, 分别用递增的有限测度集合覆盖 $X$, 再取两列集合的逐项交, 得到递增 $X_n\uparrow X$ 且 $\mu(X_n),\nu(X_n)<\infty$. 令 $D_1=X_1$, $D_n=X_n\setminus X_{n-1}$. 对每个 $D_n$ 上的限制测度应用有限情形, 得到有限的非负可测密度 $f_n$. 在两两不交的 $D_n$ 上拼接 $f=f_n$. 对任意 $E$, 可列可加性和非负积分的可列可加性给出 $\nu(E)=\sum_n\nu(E\cap D_n)=\int_Ef\dd\mu$.

最后证明唯一性. 若 $f,h$ 都是密度, 在任何同时使两测度有限的 $X_k$ 上, 集合 $A_{k,n}=X_k\cap\{f\ge h+1/n\}$ 满足
$\int_{A_{k,n}}f\dd\mu=\int_{A_{k,n}}h\dd\mu<\infty$. 由 $\int_{A_{k,n}}f\dd\mu\ge\int_{A_{k,n}}h\dd\mu+\mu(A_{k,n})/n$, 得 $\mu(A_{k,n})=0$. 对 $k,n$ 取可数并得 $\mu(\{f>h\})=0$; 交换 $f,h$ 得另一方向. 这里有限性的局部化防止从两个无穷积分相等错误地相减.
\end{proof}

\begin{proposition}\label{ra:c9:density-int}
若 $\nu=f\mu$ 表示上面的密度关系, 则对非负可测 $u$ 有 $\int u\dd\nu=\int uf\dd\mu$. 对实或复函数 $u$, 可积性等价于 $\int |u|f\dd\mu<\infty$, 此时同一等式成立.
\end{proposition}
\begin{proof}
指示函数情形就是密度定义. 简单函数情形由线性性成立. 对非负可测函数取递增简单函数逼近, 两边分别使用单调收敛定理. 对可积实函数分正负部, 对复函数分实虚部即可.
\end{proof}

\begin{example}
令 $\mu$ 为 $[0,1]$ 上的 Lebesgue 测度. 将 $\nu(E)=\int_E2x\dd x$ 与 $\rho=\nu+\delta_{1/2}$ 相比较.
\end{example}
\begin{solution}
$\nu\ll\mu$, 其密度为 $2x$. 测度 $\rho$ 在单点 $\{1/2\}$ 上的质量为1, 所以 $\rho$ 不对 $\mu$ 绝对连续, 不能写成一个相对于 $\mu$ 的函数密度. 相对于 $\lambda=\mu+\delta_{1/2}$, 它却有密度: 除 $1/2$ 外取 $2x$, 在 $1/2$ 取1. 这说明密度必须说明相对于哪一个参照测度.
\end{solution}

\originalsection{Lebesgue 分解}
上面的证明同时找出了障碍所在: $g=1$ 的集合对 $\mu$ 为零, 但没有绝对连续假设时, $\nu$ 可以在这里携带质量.

\begin{theorem}[Lebesgue 分解]\label{ra:c9:decomposition}
设 $\mu,\nu$ 为 $\sigma$ 有限正测度. 存在唯一的正测度分解 $\nu=\nu_a+\nu_s$, 其中 $\nu_a\ll\mu$, $\nu_s\perp\mu$. 绝对连续部分具有密度 $\nu_a=f\mu$.
\end{theorem}
\begin{proof}
先设两测度有限. 与上一定理相同, 以 $\lambda=\mu+\nu$ 得到 $0\le g\le1$ 和 $\nu=g\lambda$, $\mu=(1-g)\lambda$. 取 $S=\{g=1\}$, 并定义 $\nu_s(E)=\nu(E\cap S)$, $\nu_a(E)=\nu(E\setminus S)$. 因 $\mu(S)=0$, 有 $\nu_s\perp\mu$. 令 $f=g/(1-g)$ 于 $X\setminus S$, 于 $S$ 上取0. 上节的积分计算给出 $\nu_a(E)=\int_Ef\dd\mu$, 故 $\nu_a\ll\mu$.

对 $\sigma$ 有限情形使用上节的两两不交有限块 $D_n$. 在各块构造分解, 把密度拼接, 把奇异支承零集取可数并 $S=\bigcup_nS_n$. 可列可加性给出 $\nu_a,\nu_s$, 且 $\mu(S)=0$, $\nu_s(X\setminus S)=0$.

为证唯一性, 设另有 $\nu=\widetilde\nu_a+\widetilde\nu_s$, 且两奇异部分分别集中于 $\mu$ 零集 $S,\widetilde S$. 令 $Z=S\cup\widetilde S$. 两个绝对连续部分在 $Z$ 上均为零, 两个奇异部分在 $X\setminus Z$ 上均为零. 因而对任意 $E$, 两绝对连续部分都等于 $\nu(E\setminus Z)$, 两奇异部分都等于 $\nu(E\cap Z)$. 不需对无穷总质量相减.
\end{proof}

\begin{remark}
密度 $f$ 的存在并不表示奇异部分一定由原子组成. Cantor 集上存在没有原子的奇异概率测度. 本册已讨论 Cantor 零集与 Cantor 函数; 这里的分解解释了为什么某些单调函数能几乎处处导数为零而仍有总增长.
\end{remark}

\originalsection{绝对连续函数}
\noindent 本节补齐原第22讲只陈述的绝对连续逆向积分定理. 绝对连续和积分函数几乎处处求导已在第8章讨论; 此处用测度密度证明相反方向, 不把几乎处处存在导数误当作充分条件.

\begin{lemma}\label{ra:c9:variation}
设实函数 $F$ 在 $[a,b]$ 上绝对连续. 则它具有有限全变差, 函数 $V(x)=\operatorname{Var}(F;[a,x])$ 也绝对连续且递增. 因而 $F-F(a)$ 可以写成两个递增绝对连续函数之差.
\end{lemma}
\begin{proof}
全变差是所有有限分划 $a=t_0<\cdots<t_k=x$ 的和 $\sum_j|F(t_j)-F(t_{j-1})|$ 的上确界. 在绝对连续定义中取允许的总振幅为1, 并取得 $\delta>0$. 把 $[a,b]$ 分为有限个长度小于 $\delta$ 的固定区间. 任意分划与这些固定端点一起加细后, 每个固定区间内的振幅和小于1, 所以全区间变差受固定区间个数控制. 有限全变差因此成立.

在任意 $a\le s<t\le b$ 处插入分划端点, 并从两侧选择任意接近上确界的分划, 得到变差的可加性, 即 $V(t)-V(s)=\operatorname{Var}(F;[s,t])$. 给定 $\eps>0$, 在 $F$ 的绝对连续定义中选择振幅阈值 $\eps/2$ 对应的 $\delta$. 对有限个互不相交区间 $(s_i,t_i)$, 若总长度小于 $\delta$, 则各区间内任何有限分划的所有小区间仍互不相交且总长度不变. 因而其总振幅小于 $\eps/2$. 逐个取分划上确界得 $\sum_i(V(t_i)-V(s_i))\le\eps/2<\eps$. 这正是 $V$ 的绝对连续性.

令 $P=(V+F-F(a))/2$, $Q=(V-F+F(a))/2$. 因 $V(t)-V(s)\ge|F(t)-F(s)|$, $P,Q$ 递增. 它们是绝对连续函数的线性组合, 因而绝对连续, 且 $F-F(a)=P-Q$.
\end{proof}

\begin{lemma}\label{ra:c9:stieltjes}
若 $G$ 在 $[a,b]$ 上递增且绝对连续, 则存在有限 Borel 测度 $\nu_G$, 满足 $\nu_G((s,t])=G(t)-G(s)$, 且 $\nu_G\ll m$, 其中 $m$ 为 $[a,b]$ 上的 Lebesgue 测度.
\end{lemma}
\begin{proof}
下面直接把 Lebesgue 测度推到 $[a,b]$, 不调用未证的测度延拓定理. 在 $J=(G(a),G(b)]$ 上定义
$\phi(u)=\min\{x\in[a,b]:G(x)\ge u\}$. 连续性和紧性保证最小值存在. $\phi$ 递增, 所以是 Borel 可测的. 对 Borel 集 $E\subset[a,b]$, 令 $\nu_G(E)=m_1(\phi^{-1}(E))$, 其中 $m_1$ 是实轴 Lebesgue 测度. 逆像保互不相交并, 故这确为有限测度. 若 $G(a)=G(b)$, 取零测度即可.

对 $x\in[a,b]$, 连续性保证 $\phi(u)\le x$ 当且仅当 $u\le G(x)$, 因而 $\nu_G([a,x])=G(x)-G(a)$. 相减得到区间公式. 同一公式和 $G$ 连续性还说明所有单点测度为零.

设 Borel 集 $E$ 满足 $m(E)=0$. 对给定 $\eps>0$, 取 $G$ 的绝对连续阈值 $\delta$. 用总长度小于 $\delta$ 的开集覆盖 $E$, 并把开集与 $[a,b]$ 的交写成至多可数个不交区间. 端点不产生 $\nu_G$ 质量, 每个区间的质量为 $G$ 的端点增量. 绝对连续性使任意有限个这样的增量之和小于 $\eps$, 可列可加性给出 $\nu_G(E)\le\eps$. 令 $\eps\downarrow0$, 得 $\nu_G(E)=0$. 对 Lebesgue 可测集, 利用 Borel 测度的完备化扩展 $\nu_G$; 因它对 $m$ 的 Borel 零集绝对连续, 扩展与代表选择无关. 严格地说, 先在 $\nu_G$ 的完备化上定义扩展, 再限制到 Lebesgue $\sigma$-代数; 两种测度各自的完整 $\sigma$-代数不必相同.
\end{proof}

\begin{theorem}[Lebesgue 积分的微积分基本定理]\label{ra:c9:ac-ftc}
对 $[a,b]$ 上的实或复函数 $F$, 下列条件等价: $F$ 绝对连续; 存在 $f\in L^1([a,b])$ 使 $F(x)=F(a)+\int_a^x f(t)\dd t$ 对所有 $x$ 成立. 此时 $F'$ 几乎处处存在, 属于 $L^1$, 且 $F'=f$ 几乎处处. 特别地, $F(x)-F(a)=\int_a^xF'(t)\dd t$.
\end{theorem}
\begin{proof}
若存在积分表示, 对有限个互不相交区间有
$\sum_i|F(t_i)-F(s_i)|\le\int_{\bigcup_i(s_i,t_i)}|f|\dd m$. 可积函数积分对集合测度的绝对连续性给出 $F$ 绝对连续. 第8章的微分定理给出 $F'=f$ 几乎处处.

反之, 先设 $F$ 实值且绝对连续. 由引理\ref{ra:c9:variation}, 写 $F-F(a)=P-Q$, 其中 $P,Q$ 递增且绝对连续, 且 $P(a)=Q(a)=0$. 对其区间增量使用引理\ref{ra:c9:stieltjes}, 得有限测度 $\nu_P,\nu_Q\ll m$. Radon--Nikodym 定理给出非负 $p,q\in L^1(m)$, 使 $\nu_P=p\,m$, $\nu_Q=q\,m$. 于是 $P(x)=\int_a^xp\dd m$, $Q(x)=\int_a^xq\dd m$, 所以取 $f=p-q$ 就得到积分表示. 已证方向说明 $F'=f$ 几乎处处, 从而导数可积. 复值情形分别对实部、虚部应用上述证明, 再合并两个可积密度.
\end{proof}

\begin{theorem}\label{ra:c9:everywhere-ftc}
设 $F:[a,b]\to\R$ 连续, 在 $(a,b)$ 每一点都可导, 且 $F'\in L^1((a,b))$. 则 $F$ 绝对连续, 并且 $F(x)-F(a)=\int_a^x F'(t)\dd t$ 对所有 $x\in[a,b]$ 成立. 复值函数对实虚部分别应用此结论.
\end{theorem}
\begin{proof}
\textbf{编者补充原第22讲的另一充分条件.} 证明利用第4章的半连续积分逼近; 处处可导的条件将在每一点的单侧差商中使用, 不能改成仅几乎处处可导.

先说明一个实分析事实. 若连续函数 $P$ 在一个区间的每个内点满足
$D^+P(x):=\limsup_{h\downarrow0}(P(x+h)-P(x))/h\le0$, 则 $P$ 非增. 若存在内点 $s<t$ 使 $P(t)>P(s)$, 选 $0<c<(P(t)-P(s))/(t-s)$. 连续函数 $Q(x)=P(x)-cx$ 在 $[s,t]$ 上取得最小值, 且因 $Q(t)>Q(s)$, 可在某个 $x_0<t$ 取得最小值. 所有足够小的正 $h$ 满足 $Q(x_0+h)\ge Q(x_0)$, 故 $D^+Q(x_0)\ge0$. 但 $D^+Q=D^+P-c<0$, 矛盾. 向端点取连续极限后, 全区间上的非增性也成立.

记 $h=F'$, 在端点赋值0. 导数是连续函数差商的点态极限, 所以 $h$ 可测. 对任意 $\eps>0$, Vitali--Carath\'eodory 定理提供下半连续 $v\ge h$ 处处, 且 $\int_a^b(v-h)\dd t<\eps$. 由于 $h\in L^1$ 而 $v-h$ 非负可积, $v\in L^1$, 虽然个别点允许 $v=+\infty$. 令 $G(x)=\int_a^xv(t)\dd t$. 它连续.

对任意内点 $x$, 下半连续性给出
\[
 \liminf_{r\downarrow0}\frac{G(x+r)-G(x)}r\ge v(x).
\]
确实, 若 $v(x)$ 有限, 对任何 $\eta>0$, 邻域内 $v(t)>v(x)-\eta$, 积分后再令 $\eta\downarrow0$. 若 $v(x)=+\infty$, 对任意有限 $M$ 使用同一论证, 得左边为 $+\infty$. 因 $F$ 在该点有有限导数, $P=F-G$ 满足 $D^+P(x)\le h(x)-v(x)\le0$. 前面的事实说明 $P$ 非增. 因而对 $a\le s<t\le b$,
\[
 F(t)-F(s)\le\int_s^tv\dd x
 \le\int_s^th\dd x+\eps.
\]
对 $-F$ 和导数 $-h$ 应用相同构造, 得反向不等式, 误差仍可任意小. 令 $\eps\downarrow0$, 得 $F(t)-F(s)=\int_s^th\dd x$. 积分表示以及定理\ref{ra:c9:ac-ftc}的已证方向给出绝对连续性.
\end{proof}

仅连续且几乎处处可导不足以保证这个公式. Cantor 函数连续递增, 几乎处处导数为零, 但端点总增量为1; 它不绝对连续. 本节的测度分解说明, 正是不能由 Lebesgue 密度表示的奇异增长阻止了从导数恢复函数.

\originalsection{练习与衔接}
\begin{exercise}
设 $\nu\ll\mu$, $\rho\ll\nu$, 且三者为 $\sigma$ 有限正测度. 写出 $\rho$ 相对于 $\mu$ 的密度, 并核对可能出现零密度的集合.
\end{exercise}
\begin{solution}
取非负且几乎处处有限的 $f=\dd\nu/\dd\mu$ 和 $g=\dd\rho/\dd\nu$. 因 $g$ 只在 $\nu$ 几乎处处意义下确定, 先在 $\{f=0\}$ 上令 $g=0$, 再在其余 $\nu$ 零集上取有限代表; 这不会改变积分. 由命题\ref{ra:c9:density-int}, 对每个 $E$ 有 $\rho(E)=\int_Eg\dd\nu=\int_Egf\dd\mu$. 因而 $\dd\rho/\dd\mu=gf$ 几乎处处. 如不作代表的调整, 则约定乘积在 $f=0$ 处为零; 这正是非负积分中 $0\cdot\infty=0$ 的约定.
\end{solution}

\begin{exercise}
说明测度的 $\sigma$ 有限条件不能无条件删除: 在不可数集合 $X=[0,1]$ 的 Lebesgue 可测集上, 取 $\mu$ 为计数测度, $\nu$ 为 Lebesgue 测度.
\end{exercise}
\begin{solution}
计数测度的零集只有空集, 所以 $\nu\ll\mu$. 如果存在非负密度 $f$, 单点等式会给 $f(x)=\nu(\{x\})=0$ 对每个 $x$ 成立. 于是 $f$ 恒为零, 从而 $\int_X f\dd\mu=0$, 与 $\nu(X)=1$ 矛盾. 这里 $\mu$ 不是 $\sigma$ 有限: 每个有限计数测度集合有限, 可数个有限集不能覆盖不可数的 $X$.
\end{solution}

\begin{exercise}
设 $\nu=f\mu$ 且 $\nu(X)<\infty$. 直接用 $f$ 证明命题\ref{ra:c9:uniform-ac}, 并说明有限质量的用途.
\end{exercise}
\begin{solution}
因 $f\in L^1(\mu)$, 可选 $M>0$ 使 $\int_{\{f>M\}}f\dd\mu<\eps/2$. 若 $\mu(E)<\eps/(2M)$, 则
\[
 \nu(E)\le\int_{E\cap\{f\le M\}}f\dd\mu
       +\int_{\{f>M\}}f\dd\mu
 <M\mu(E)+\eps/2<\eps.
\]
有限质量保证截尾积分趋于零. 没有它时, 零集绝对连续并不自动给出对全部可测集统一的小测度控制.
\end{solution}

前文的 Hölder 不等式说明, 若 $1<p<\infty$, $q=p/(p-1)$, 每个 $g\in L^q(\mu)$ 都定义 $L^p$ 上的有界线性泛函 $T_g(f)=\int fg\dd\mu$. 本章说明了从测度恢复积分密度的机制. 从一般有界泛函构造测度、处理复数共轭约定并证明对偶空间的等距同构, 属于泛函册的衔接内容; 这里不把尚未建立的完整 $L^p$ 对偶定理当作已证结果.

\originalchapter{来源对照}
\noindent 下表按实际下载的 PDF 正文建立对应, 页数为源文件物理页数. 官网主题表有错位: Cantor 集与非 Borel 可测集为第3章编者补充; 原第17讲没有平滑稠密证明, 该内容实际在第21讲. 原第19讲已经陈述一般 Fubini, 第20讲已经给出 Young 不等式. 原第24讲只证明最大算子的专门插值估计. 中文补证、勘误与新增练习见各章正文及独立核查记录.

\begin{longtable}{@{}p{.07\textwidth}p{.07\textwidth}p{.66\textwidth}p{.11\textwidth}@{}}
\toprule 讲 & 源页 & 实际内容 & 本册章\\\midrule\endhead
\href{https://ocw.mit.edu/courses/18-125-measure-and-integration-fall-2003/resources/18125_lec1/}{1} & 3 & 可测空间、可测复合与Borel集 & 1 \\
\href{https://ocw.mit.edu/courses/18-125-measure-and-integration-fall-2003/resources/18125_lec2/}{2} & 4 & 实可测函数、极限、简单函数与非负积分 & 1,2 \\
\href{https://ocw.mit.edu/courses/18-125-measure-and-integration-fall-2003/resources/18125_lec3/}{3} & 3 & Riemann积分判据、测度连续性、积分性质 & 1,2 \\
\href{https://ocw.mit.edu/courses/18-125-measure-and-integration-fall-2003/resources/18125_lec4/}{4} & 4 & 单调收敛、积分可加性、级数与Fatou & 2 \\
\href{https://ocw.mit.edu/courses/18-125-measure-and-integration-fall-2003/resources/18125_lec5/}{5} & 5 & 复积分、主导收敛、零集和完备化 & 2 \\
\href{https://ocw.mit.edu/courses/18-125-measure-and-integration-fall-2003/resources/18125_lec6/}{6} & 4 & 矩形与多面体、开紧集、内外测度 & 3 \\
\href{https://ocw.mit.edu/courses/18-125-measure-and-integration-fall-2003/resources/18125_lec7/}{7} & 4 & 可测集定义和Lebesgue测度可列可加性 & 3 \\
\href{https://ocw.mit.edu/courses/18-125-measure-and-integration-fall-2003/resources/18125_lec8/}{8} & 3 & Carathéodory准则、开集紧集逼近、完备性 & 3 \\
\href{https://ocw.mit.edu/courses/18-125-measure-and-integration-fall-2003/resources/18125_lec9/}{9} & 4 & 不变性、Vitali不可测集、线性变换 & 3 \\
\href{https://ocw.mit.edu/courses/18-125-measure-and-integration-fall-2003/resources/18125_lec10/}{10} & 5 & Riesz表示、由Riemann泛函构造Lebesgue测度 & 4 \\
\href{https://ocw.mit.edu/courses/18-125-measure-and-integration-fall-2003/resources/18125_lec11/}{11} & 3 & Lusin、Vitali–Carathéodory & 4 \\
\href{https://ocw.mit.edu/courses/18-125-measure-and-integration-fall-2003/resources/18125_lec12/}{12} & 3 & 连续函数逼近、几乎处处收敛和集合可加积分 & 2,4 \\
\href{https://ocw.mit.edu/courses/18-125-measure-and-integration-fall-2003/resources/18125_lec13/}{13} & 4 & Egorov、依测度收敛、子列及主导收敛 & 4 \\
\href{https://ocw.mit.edu/courses/18-125-measure-and-integration-fall-2003/resources/18125_lec14/}{14} & 4 & 凸函数、Jensen、Hölder、Minkowski & 5 \\
\href{https://ocw.mit.edu/courses/18-125-measure-and-integration-fall-2003/resources/18125_lec15/}{15} & 5 & Lp商空间、范数与Riesz–Fischer完备性 & 5 \\
\href{https://ocw.mit.edu/courses/18-125-measure-and-integration-fall-2003/resources/18125_lec16/}{16} & 2 & Cc在Lp稠密、Cc在C0稠密 & 5 \\
\href{https://ocw.mit.edu/courses/18-125-measure-and-integration-fall-2003/resources/18125_lec17/}{17} & 4 & Lp与lp包含、局部Lp、凸性 & 5 \\
\href{https://ocw.mit.edu/courses/18-125-measure-and-integration-fall-2003/resources/18125_lec18/}{18} & 6 & Rn非负及可积Fubini、截面论证和双积分例 & 6 \\
\href{https://ocw.mit.edu/courses/18-125-measure-and-integration-fall-2003/resources/18125_lec19/}{19} & 3 & 可测直积、一般乘积测度与Fubini & 6 \\
\href{https://ocw.mit.edu/courses/18-125-measure-and-integration-fall-2003/resources/18125_lec20/}{20} & 3 & 卷积、Young、共轭卷积一致连续 & 7 \\
\href{https://ocw.mit.edu/courses/18-125-measure-and-integration-fall-2003/resources/18125_lec21/}{21} & 4 & 近似恒等、磨光、平滑稠密 & 7 \\
\href{https://ocw.mit.edu/courses/18-125-measure-and-integration-fall-2003/resources/18125_lec22/}{22} & 5 & 绝对连续、Vitali覆盖与弱L1最大估计 & 8,9 \\
\href{https://ocw.mit.edu/courses/18-125-measure-and-integration-fall-2003/resources/18125_lec23/}{23} & 5 & Lebesgue微分、Lebesgue点、积分函数求导 & 8 \\
\href{https://ocw.mit.edu/courses/18-125-measure-and-integration-fall-2003/resources/18125_lec24/}{24} & 5 & 广义Minkowski、分布函数、最大算子Lp估计 & 8 \\
\bottomrule
\end{longtable}
第9章的积分密度与 Lebesgue 分解是编者补充, 不对应上述任何源 PDF; 其中绝对连续函数一节补证第22讲只陈述的逆向积分定理. 课程大纲中的 Hausdorff 测度、面积与余面积公式没有被包含在实际讲义中, 因此尚未译编. 课程描述提及 Fourier transform, 这24份文件也没有对应正文. 独立作业集与外部教材原题不属于本册的来源范围.

\begin{thebibliography}{9}
\bibitem{ocw} Jeff Viaclovsky (讲授), Ethan Brown (排录). \textit{18.125 Measure and Integration, Fall 2003}. MIT OpenCourseWare. \href{https://ocw.mit.edu/courses/18-125-measure-and-integration-fall-2003/pages/lecture-notes/}{24份讲义目录}; 逐份原件页码见来源对照.
\bibitem{terms} MIT OpenCourseWare. \href{https://ocw.mit.edu/pages/privacy-and-terms-of-use/}{Privacy and Terms of Use}. 核验日期2026-10-08. \href{https://creativecommons.org/licenses/by-nc-sa/4.0/}{CC BY-NC-SA 4.0 International}.
\end{thebibliography}
\end{document}
